%%%%%%%%%%%%%%%%%%%%%%%%%%%%%%%%%%%%%%%%%%%%%%%%%%%%%%%%%%%%%%%%%%%%%%%%%%%
% filename    : sp.tex
% author      : 
%%%%%%%%%%%%%%%%%%%%%%%%%%%%%%%%%%%%%%%%%%%%%%%%%%%%%%%%%%%%%%%%%%%%%%%%%%%

% BUGFIX: Save LaTeX kernel version of \@xfloat
\makeatletter
\let\my@xfloat\@xfloat
\makeatother

% use ateneo de naga etd class (actually modified virginia tech's etd class)
\documentclass[oneside]{etd}

% BUGFIX: Create a modified copy of \@xfloat using the kernel definition
\makeatletter
\def\@xfloat#1[#2]{
	\my@xfloat#1[#2]%
	\def\baselinestretch{1}%
	\@normalsize \normalsize
}
\makeatother

% use these packages
\usepackage{tabularx}
\usepackage{booktabs}
\usepackage{graphicx}
\usepackage{amssymb}
\usepackage{textcomp}
\usepackage{gensymb}
\usepackage{url}
\usepackage{hyperref}
\usepackage{doi}
\usepackage{latexsym}
\usepackage{amsmath}
%\usepackage[lineno5,noindent]{lgrind}
%\usepackage{rotating} 
%\usepackage{makeidx}
%\usepackage{stmaryrd}
\usepackage{float}
\usepackage{subcaption}
%\usepackage{cite}
%\usepackage{moreverb}
%\usepackage{pictexwd,dcpic}
\usepackage{algorithm,algorithmic}
\renewcommand{\algorithmicrequire}{\textbf{Input:}}
\renewcommand{\algorithmicensure}{\textbf{Output:}}
\usepackage{tikz}
\usetikzlibrary{positioning,calc}

\usepackage{listings}
\definecolor{cppred}{rgb}{0.6,0,0} % for strings
\definecolor{cppgreen}{rgb}{0.25,0.5,0.35} % comments
\definecolor{cpppurple}{rgb}{0.5,0,0.35} % keywords
\definecolor{cppdocblue}{rgb}{0.25,0.35,0.75} % doc
 
\lstset{language=C++,
basicstyle=\linespread{0.8}\ttfamily\small,
keywordstyle=\color{cpppurple}\bfseries,
stringstyle=\color{cppred},
commentstyle=\color{cppgreen},
morecomment=[s][\color{cppdocblue}]{/**}{*/},
numbers=left,
numberstyle=\tiny\color{black},
% stepnumber=2,
numbersep=10pt,
tabsize=4,
showspaces=false,
showstringspaces=false}

\graphicspath{{figures/}{./}}

%\renewcommand{\floatpagefraction}{0.8}

\makeindex

%%%%%%%%%%%%%%%%%%%%%%%%%%%%%%%%%%%%%%%%%%%%%%%%%%%%%%%%%%%%%%%%%%%%%%%%%%%

\title{MARQ: A Multi-Agent Rainbow Deep Q-Networks Framework using AAREN for Mastering Imperfect Information Games}
\department{Computer Science}
\BSthesis % options: \BSthesis \BSproject \BShonors \Mdegree \MSdegree \PhDdegree
\degreemonth{September 5}  % date of final defense
\degreeyear{2025}

\twostudents % options: \onestudent \twostudents \threestudents \fourstudents
{James Gabriel P. Mabagos}
{Mark Lawrence M. Quibot}


\studentsheader
{Mabagos and Quibot}

\twodegrees % options: \onedegree \twodegrees \threesdegrees \fourdegrees
{Computer Science}
{Computer Science}

\advisor
{Marianne A. Tolentino}
\threemembers
{Estrella H. Montealegre}
{
Michelle Elija B. Santos}
{Ian Peter L. Lastimosa}
\deanandchair
{Joshua C. Martinez MIT }
{Lowie Vincent S. Bisana MIT}

\keywords{Imperfect Information, DQN, Multi-Agent}

\begin{document}

%edit etd.cls to accommodate lesser number of proponents
\maketitle
\makerecomm
\makeacceptance
\makedeclaration

%%%%%%%%%%%%%%%%%%%%%%%%%%%%%%%%%%%%%%%%%%%%%%%%%%%%%%%%%%%%%%%%%%%%%%%%%%%
%                    ABSTRACT/EXECUTIVE SUMMARY                           %
%%%%%%%%%%%%%%%%%%%%%%%%%%%%%%%%%%%%%%%%%%%%%%%%%%%%%%%%%%%%%%%%%%%%%%%%%%%

% can use execsummary or abstract
\begin{abstract}
This research presents MARQ, a Multi-Agent Rainbow Deep Q-Networks framework for imperfect information games, and evaluates it on the board game Stratego. MARQ combines a Rainbow Deep Q-Network with a curriculum designed to end in league-based self-play, to address the size of Stratego's game tree and its hidden information. The architecture integrates an Attention as Recurrent Neural Network (AAREN) module that produces learned embeddings from observed opponent behavior. AAREN is trained to predict the identity of revealed pieces, and the Q-network reads its embeddings as inputs for implicit piece identity inference. The system employs Distributional Reinforcement Learning with Noisy Networks in place of the scalar value estimate and the $\epsilon$-greedy exploration of a standard DQN baseline. The methodology involves developing a game environment and training the AI agent through multi-channel state representations, incorporating AAREN-generated embeddings for hidden information modeling and piece values computed analytically from the combat rules. The agents are evaluated through agent-versus-agent games, by win rate, draw rate and game length. In self-play, the draw rate fell from 50.04\% for the baseline to 6.69\% for MARQ. Against fixed heuristic opponents, however, MARQ won 11.66\% of 75,047 games and drew 77.02\%, and no ablation isolates the contribution of any component. Follow-up runs with a corrected learner and AAREN switched off found that the agent played at the level of the scripted player it imitated before training: cloning a scripted player that looks one move ahead gave 29 wins in 40 games against the heuristic, and 5,000 episodes of reinforcement learning did not improve on it. The study does not show that MARQ plays Stratego at a competitive level. It contributes a tested environment and training system, a set of checks that reversed results first read as validation, and a measured account of why a value-based learner stalled on this game at this scale of computation.
\end{abstract}

\include{misc/dedication}
\include{misc/acknowledgements}

\begingroup
\renewcommand*{\addvspace}[1]{}
\tableofcontents
\listoffigures
\listoftables
\endgroup

\beginbody
\pagenumbering{arabic}

\chapter{Introduction}
\section{Project Context}
Games have long served as benchmarks for measuring Artificial Intelligence (AI) progress \cite{StudeGame}. Many real-world problems, from military tactics to multiplayer games, require intelligent decision-making in environments where all information is not available. A common example is the ``fog of war'' mechanic found in many multiplayer games \cite{FogofWar}, where the game environment is not fully revealed to any single player, placing these games in the category of imperfect information games \cite{DRLImperfectInfo2023}. In such games, players must develop strategies without full knowledge of the game's environmental state \cite{FogofWar}. Determining the optimal strategy under these conditions remains a standing challenge for Artificial Intelligence \cite{li2024opponent}.

This study focuses on mastering imperfect information games using Multi-Agent Reinforcement Learning (MARL) with Deep Q-Networks (DQN). In MARL, multiple agents learn and adapt in the same environment, making each decision based on its own observations and rewards \cite{2509_03682v1}. With DQN, the agents use deep neural networks to approximate the Q-function that maps state-action pairs to expected reward. This study introduces the MARQ (Multi-Agent Rainbow Deep Q-Networks) framework, an extension of the DQN architecture with Rainbow enhancements: distributional value estimation (C51) and noisy exploration networks \cite{RainbowJMLR2021}. An AAREN (Attention as Recurrent Neural Network) module \cite{Aaren} processes extended observation sequences and produces learned embeddings meant to carry what the opponent's moves reveal about its hidden pieces. AAREN is trained on the pieces that battles reveal, and the Q-network reads its embeddings as inputs. League-based self-play training \cite{perolat2022stratego} introduces agents to different historical opponents to prevent similar strategies. Stratego is a good example of a test environment for our framework, agents compete against each other and must develop strategies through repetition of experience while reasoning under uncertainty about the enemy's setup and the environments hidden states.



\section{Purpose and Description}
This research evaluates whether Multi-Agent Reinforcement Learning (MARL) with Deep Q-Networks (DQN) can produce agents that can play imperfect information games at a competitive level. The board game Stratego serves as the test domain because it combines hidden piece identities, a game tree of about $10^{535}$ nodes \cite{perolat2022stratego}, and long game horizons. Existing analytical tools for Stratego have little to no support for the deeper tactical reasoning that we want to achieve. 

The MARQ framework was developed to test whether targeted architectural changes, specifically distributional value estimation, attention-based history encoding, and learned exploration, can overcome the performance ceiling observed in a standard DQN baseline. The experiments, detailed in later chapters, do not show that they did. In self-play the combined changes cut the draw rate from 50.04\% to 6.69\%, and the self-play win rate rose from 24.97\% to 49.88\%. Both sides of a self-play game are the same agent, however, so that rise follows from the fall in draws and does not measure strength (Section~\ref{finding:7}). Against fixed heuristic opponents the Rainbow agent with AAREN won 11.66\% of 75,047 games and drew 77.02\%. The runs made after those experiments, with AAREN switched off, found that the agent played at the level of the scripted player it had imitated before training, and that reinforcement learning did not raise that level (Appendix~\ref{app:lessons}). The thesis therefore reports what was built, what was measured, and a diagnosis of why the learner stalled.
\subsection{Research Questions}
This study is guided by three main research questions:

\begin{enumerate}
   
    \item How does the MARQ framework perform in terms of convergence speed and win-rate against heuristic baselines in a partially observable environment?
    \item Does the MARQ framework demonstrate statistically significant superiority over standard DQN and rule-based agents across varying opponent difficulty tiers?
    \item What is the impact of integrating the AAREN module on the agent's ability to infer hidden opponent piece configurations compared to approaches without explicit belief state modeling?
\end{enumerate}

Chapter~\ref{chap:recommendations} answers each question as far as the data allow.

\section{Objectives}
The main objective of this study is to develop and evaluate the MARQ framework for Stratego, where agents must operate under hidden state variables. The specific objectives are:
\begin{itemize}
    \item Identify and evaluate reinforcement learning algorithms for competitive play in Stratego.
    \item Design and implement a game environment for Stratego
    \item Build and implement a training model using Multi-Agent Reinforcement Learning with Deep Q-Networks and AAREN.
    \item Train and optimize the AI agent to respond to varying game states, hidden information, and opposing strategies.
\end{itemize}

\section{Scope and Limitation}

The study is limited to the development of an AI agent for imperfect information games, with Stratego as the test environment. The scope includes the design and implementation of Stratego as a game environment, the development of a training framework using Multi-Agent Reinforcement Learning (MARL) with Deep Q-Networks (DQN) and AAREN, and the training and optimization of the AI agent to adapt to varying game states. Agent performance is evaluated through agent-vs-agent testing within the Stratego environment. Each configuration was trained once, and no ablation that removes one component at a time was run, so the results describe the combined system and not the contribution of each part. Training is limited to the personal devices of the researchers and does not rely on cloud-based computing resources.

\section{Significance of the Study}
Current imperfect information game systems either require large computational resources or use model-free approaches with limited strategic depth. The MARQ framework contributes a specific combination of techniques: implicit opponent modeling through AAREN-generated embeddings, distributional value estimation via C51, and learned exploration via Noisy Networks. The experiments do not establish that this combination plays better than the standard DQN baseline: the two were never evaluated against each other or against the same fixed opponents, and against heuristic opponents the combined agent won 11.66\% of its games.

The study contributes three things that hold whatever the outcome. The first is a tested Stratego environment and training system on which a value-based learner can be run and examined on a single consumer GPU, with a network that can be played in a browser. The second is a set of checks that reversed results the study had first read as validation: a self-play win rate that measured the draw rate, a loss--win correlation that measured a shared trend, an average Q-value taken over illegal actions, an accuracy compared with the wrong baseline, and a gradient path that reached a single bias (Chapter~\ref{chap:results}, Appendix~\ref{app:lessons}). The third is a measured account of where the agent's strength came from: the demonstrator it imitated set its level, a behaviour cloning loss over the legal moves transferred the demonstrator's decisions where a large-margin loss had not, and the policy was lost when the cloning term was withdrawn (Appendix~\ref{app:lessons}, Section~\ref{sec:formal14}). These apply to other partially observable or adversarial domains in which a value-based learner is started from demonstrations.

\section{Definition of Terms}

\noindent\textbf{AAREN (Attention as Recurrent Neural Network)} \\
A neural architecture that combines the parallel training efficiency of Transformers with the sequential processing of RNNs. In the MARQ framework, AAREN processes observation sequences to produce learned embeddings that encode implicit piece identity inference. \\[6pt]

\noindent\textbf{Counterfactual Regret Minimization (CFR)} \\
An iterative algorithm that approximates Nash equilibrium by minimizing regret, widely applied in imperfect information games such as Poker and Stratego. \\[6pt]

\noindent\textbf{Deep Q-Networks (DQN)} \\
A reinforcement learning algorithm that combines Q-learning with deep neural networks, enabling agents to approximate optimal decision-making in complex environments. \\[6pt]

\noindent\textbf{Distributional Reinforcement Learning} \\
An approach that models the full distribution of returns rather than just the expected value. The MARQ framework uses C51 (Categorical DQN) with 51 atoms to capture uncertainty in value estimates. \\[6pt]

\noindent\textbf{Fog of War} \\
A game mechanic used in strategy games that conceals parts of the map or opponent actions, requiring players to utilize prediction and information-gathering strategies to reduce uncertainty. \\[6pt]

\noindent\textbf{Imperfect Information Games} \\
Games where players have incomplete knowledge of the current game state, requiring them to make strategic decisions under uncertainty (e.g., Stratego, Poker, Multiplayer Online Battle Arenas). \\[6pt]

\noindent\textbf{Information Set Monte Carlo Tree Search (IS-MCTS)} \\
A search algorithm for imperfect information games that samples consistent game states from belief distributions and performs tree search on information sets rather than individual states. \\[6pt]

\noindent\textbf{League Training} \\
A self-play training paradigm where agents compete against a diverse pool of historical opponents, preventing strategy cycling and ensuring robust generalization across different playing styles. \\[6pt]

\noindent\textbf{Multi-Agent Reinforcement Learning (MARL)} \\
An extension of reinforcement learning in which multiple agents interact within the same environment, learning coordinated or competitive strategies through simultaneous training. \\[6pt]

\noindent\textbf{Nash Equilibrium} \\
A game-theoretic concept where no player can gain an advantage by unilaterally changing their strategy, often used as a benchmark for decision-making in imperfect information games. \\[6pt]

\noindent\textbf{Perfect Information Games} \\
Games in which all players have complete knowledge of the current game state and available moves, such as Chess and Go. \\[6pt]

\noindent\textbf{AAREN Embeddings} \\
Embedding vectors produced by the AAREN module from observed opponent action sequences. \\[6pt]

\noindent\textbf{Rainbow DQN} \\
An enhanced DQN architecture that combines multiple improvements including distributional RL, noisy networks for exploration, and dueling network architecture for improved stability and sample efficiency. \\[6pt]

\noindent\textbf{Recurrent Neural Network (RNN)} \\
A type of neural network designed to process sequential data by maintaining memory of previous inputs, useful in imperfect information games for recalling past moves or revealed information. \\[6pt]

\noindent\textbf{Reinforcement Learning (RL)} \\
A machine learning paradigm where agents learn optimal strategies through trial-and-error interactions with the environment, guided by rewards and penalties. \\[6pt]

\noindent\textbf{Stratego} \\
A two-player strategy board game characterized by hidden information and asymmetric piece abilities, where the objective is to capture the opponent's flag or render them immobile. \\[6pt]
\include{chapters/review}
\chapter{Theoretical Frameworks}

\section{Stratego Game Domain}

Stratego is an imperfect information board game with European roots, popularized in the Netherlands in the 1940s and released in North America in 1961 \cite{perolat2022stratego}. The game is played on a 10x10 board that includes two 2x2 ``lake'' areas in the center, which are impassable. Each player commands an army of 40 pieces of varying ranks, and during the game, the identity and rank of the pieces are kept secret from the opposing player. The objective is to capture the opponent's flag or eliminate all of their movable pieces.

Each army consists of 12 different types of pieces, with ranks ranging from 1 (the Spy) to 10 (the Marshal), alongside Bombs and a Flag. Higher-ranked pieces defeat lower-ranked ones during combat, except in special situations: the Spy can defeat a Marshal if attacking first, the Miner is the only piece that can defuse Bombs, and Scouts can move multiple squares in a straight line. There are 3 ways to win the game capturing the flag, immobilizing all movable enemy pieces, or forcing the opponent to run out of time.


\begin{figure}[ht]
    \centering
    \includegraphics[width=0.5\textwidth]{figures/stratego_piece.png}
    \caption{Stratego pieces.}
\end{figure}

\begin{figure}[ht]
    \centering
    \includegraphics[width=0.7\textwidth]{figures/stratego.png}
    \caption{Stratego Game board.}
    \label{fig:stratego_tech}
\end{figure}

The number of ways to start a game of Stratego already shows how difficult the game is to solve. Each player arranges 40 pieces of 12 types (one Flag, Spy, Marshal and General, 2 Colonels, 3 Majors, 4 Captains, 4 Lieutenants, 4 Sergeants, 5 Miners, 8 Scouts and 6 Bombs) on 40 squares:

\begin{equation}
\frac{40!}{1!\,1!\,1!\,1!\,2!\,3!\,4!\,4!\,4!\,5!\,8!\,6!} \approx 1.4 \times 10^{33}
\end{equation}

The two players' setups together give about $10^{66}$ starting positions, and Perolat et al. \cite{perolat2022stratego} put the size of the game tree at about $10^{535}$ nodes. The agent must reason with limited information rather than try to look at every possible outcome, which is intractable.


\section{MARQ Framework Overview}

The MARQ framework is a Multi-Agent Rainbow Deep Q-Networks system that trains agents through self-play and deep reinforcement learning \cite{AlphaStratego}. Agents learn by competing against many different opponents, including their own previous versions and exploratory variants, within a league-based self-play system.

The central problem is reasoning about hidden information without enumerating every possible state. The game tree has about $10^{535}$ nodes, and the two players' initial setups give about $10^{66}$ starting positions \cite{perolat2022stratego}. MARQ handles this scale through implicit belief representations, probabilistic reasoning, and pattern recognition. Deep reinforcement learning estimates values while attention mechanisms process game history and mixed strategies prevent exploitation, all within real-time computational constraints.

The architecture connects these components in a pipeline. The AAREN module takes raw game positions and move sequences and produces learned embeddings that encode opponent piece identity inferences. These embeddings are the input for all subsequent reasoning. A league-based training method trains and evaluates MARQ agents simultaneously.



\begin{table}[htbp]
    \centering
    \caption{Key symbols and descriptions.}
    \label{tab:MARQ-symbols}
    \begin{tabular}{|l|p{4cm}|p{6cm}|}
        \hline
        \textbf{Symbol} & \textbf{Description} & \textbf{Example (MARQ Framework)} \\
        \hline
        $\pi_i(a|I_i)$ & Policy for agent i choosing action a given information set $I_i$ & Probability of moving a piece forward given current visible board state and belief about opponent pieces \\
        \hline
        $\mathcal{B}_t$ & Implicit belief representation at time t & Probability distribution over possible identities and positions of hidden opponent pieces inferred by AAREN \\
        \hline
        $v^{\pi}_i(h, \mathcal{B}_t)$ & Value function for agent i following policy $\pi$ given history h and belief $\mathcal{B}_t$ & Expected reward for current board position considering uncertainty about opponent setup \\
        \hline
        $Q_{\text{network}}(s', a)$ & Deep Q-Network approximation of action-value function & Neural network estimate of value for attacking with a suspected Marshal vs opponent piece \\
        \hline
        $R^{\text{capture}}{i,t+1}$ & Immediate reward from piece captures weighted by piece values & +9 points for capturing opponent Marshal, -1 for losing a Scout \\
        \hline
        $O = \{o_1, \dots, o_t\}$ & Observation sequence processed by AAREN attention mechanism & Sequence of board states, moves, and attack outcomes throughout the game \\
        \hline
        $\alpha_i$ & Attention weights for historical observations in AAREN & Higher weight on turn when opponent's Marshal was revealed, lower on routine Scout moves \\
        \hline
        $\pi^*$ & Nash equilibrium policy profile & Optimal mixed strategy that cannot be exploited by opponent strategy changes \\
        \hline
        $Q_{eval}(\mathcal{B})$ & Embedding quality evaluator output & Confidence score for learned representation accuracy \\
        \hline
        $H(\pi)$ & Policy entropy measure & Ensures minimum unpredictability to prevent exploitation \\
        \hline
        $\epsilon$ & Exploration rate for agent behavior & Controls exploration vs exploitation balance in multi-agent learning \\
        \hline
        $\gamma$ & Discount factor for future rewards & Importance of long-term vs immediate rewards in value estimation \\
        \hline
    \end{tabular}
\end{table}


\section{MARQ Policy Definition}

Policy optimization in reinforcement learning improves an agent's strategy to maximize total reward \cite{DRLImperfectInfo2023}. Games with hidden information and multiple players require policies that perform well against fixed opponents and also handle uncertainty and shifting strategies from other agents. MARQ combines belief state reasoning with value-based decisions.

Policy learning in MARQ serves two goals. The first is winning games through tactical choices. The second is improving the quality of the embeddings that the AAREN module produces. The two objectives reinforce each other: better embeddings sharpen decision-making, and good decisions reveal opponent information that in turn refines the embeddings.

\subsection{Core Policy Framework}

A policy in the MARQ framework, denoted $\pi_i(a|I_i)$, defines a probability distribution over actions for agent $i$ conditioned on its current information set $I_i$. In perfect information games, policies operate on complete game states. MARQ policies operate under uncertainty about the true state of the environment. The information set $I_i$ contains all observations available to the agent: visible board positions, revealed piece identities, and the history of observed movements and combat outcomes.

The core difficulty is that agents must act without knowing the opponent's private information. In Stratego, this means the rank and position of every unrevealed opponent piece remain unknown. The policy needs to take account for variations of the possible hidden pieces alongside the opponents strategy.

MARQ uses a implicit belief representations $\mathcal{B}_t$ rather than calculating every possible game state. The AAREN module learns and updates these representations as new information arrives. The policy takes the following form:

\begin{equation}
\pi_i(a|I_i, \mathcal{B}_t) = P(\text{action} = a \mid \text{information set } I_i, \text{implicit belief } \mathcal{B}_t)
\end{equation}

The policy learns to weight actions based on their expected value across the distribution of opponent configurations rather than overfitting to a single enemy configuration.

\subsection{Implicit Representation Optimization}

MARQ treats representation accuracy as a training objective and uses revealed piece identities as supervised labels for embedding refinement.

When combat reveals the piece identities, the system calculates the prediction error using Kullback-Leibler (KL) divergence.

\begin{equation}
\mathcal{E}_{\text{pred}}(t) = \sum_{p \in \text{Pieces}_{\text{revealed}}(t)} D_{\text{KL}}\left(\text{Representation}_t(p) \| \delta_{s^*_p}\right)
\end{equation}

Here, $\delta_{s^*_p}$ is the Dirac delta distribution concentrated on the true piece type. The error quantifies the distance of the predicted distribution to the actual identity.

This prediction error becomes is the reward signal that encourages the policy to better predict actual piece configurations:

\begin{equation}
R_{\text{accuracy}}(t) = -\mathcal{E}_{\text{pred}}(t) = -\sum_{p \in \text{Pieces}_{\text{revealed}}(t)} \left[\text{Representation}_t(p=s^*_p) > 0\right] \cdot \log \text{Representation}_t(p=s^*_p)
\end{equation}

The optimal policy maximizes:

\begin{equation}
\pi^*_i = \arg\max_{\pi_i} \mathbb{E}_{\tau \sim \pi_i}\left[\sum_{t=0}^T \gamma^t \left(R_{\text{game}}(t) + \eta \cdot R_{\text{accuracy}}(t)\right)\right]
\end{equation}

The hyperparameter $\eta \in \mathbb{R}^+$ controls the accuracy versus immediate game rewards. The value function in MARQ extends the standard Bellman equation to account for game state belief, expressing the expected cumulative reward as a function of the current known information set and implicit belief.

\begin{equation}
v^{\pi}_i(I, \mathcal{B}_t) = \mathbb{E}_{\pi}\left[R^{\text{capture}}_{i,t+1} + \lambda_1 R^{\text{info}}_{i,t+1} + \lambda_2 R^{\text{position}}_{i,t+1} + \gamma v^{\pi}_i(I_{t+1}, \mathcal{B}_{t+1}) \mid I_t = I, \mathcal{B}_t\right]
\end{equation}

Three components make up the reward function. The capture reward $R^{\text{capture}}$ for capturing a piece. The information reward $R^{\text{info}}$ for information revelation. The position reward $R^{\text{position}}$ for territorial control.

MARQ training aims to reach a Nash equilibrium. At equilibrium, no player can improve their expected payoff.

\begin{equation}
v_i(\pi^*) = \max_{\pi_i} v_i(\pi_i, \pi^*_{-i})
\end{equation}

Reaching Nash equilibria in a large imperfect information game is computationally intractable. MARQ seeks approximate equilibria through self-play training with diverse opponent populations.


\section{Deep Q-Network Architecture}

The Deep Q-Networks in MARQ operates on implicit belief representations instead of a complete game states \cite{DRLImperfectInfo2023}. The network approximates Q-values over probability distributions of opponent configurations, $\mathcal{B}_t$:
$$Q(s, a) = \mathbb{E}_{s' \sim \mathcal{S}_t \subseteq \mathcal{B}_t}[Q_{network}(s', a)]$$

A standard replay buffer stores experience. Transitions $(\mathcal{B}_t, a_t, r_t, \mathcal{B}_{t+1})$ are sampled uniformly to train the value function.

\subsection{Rainbow DQN Components}
The implementation uses the standards of the Rainbow DQN architecture \cite{RainbowJMLR2021}.

\paragraph{Distributional RL (Categorical DQN).}
The Categorical DQN (C51) \cite{DistributionalRL2023} models the value distribution directly. It captures the uncertainty of the game by approximating the distribution with 51 atoms spanning the range from $V_{min}$ to $V_{max}$.

\begin{equation}
Z_\theta(s, a) = \sum_{i=1}^{N} p_i(s, a) \delta_{z_i}
\end{equation}
Training minimizes the Kullback-Leibler divergence between the predicted distribution and the projected target distribution:
\begin{equation}
D_{KL}(\Phi \hat{\mathcal{T}} Z_{\theta'} || Z_\theta)
\end{equation}
where $\Phi$ is the projection operator that maps the Bellman update onto the fixed support atoms.

\paragraph{Noisy Networks for Exploration.}
Noisy Linear layers \cite{Yang2023Exploration} adds a learned noise into the end of the network layers. This creates an unpredictable strategy without a fixed exploration variable, and the noise parameters are tuned during training for each state.

\begin{equation}
y = (\mu_w + \sigma_w \odot \epsilon_w)x + (\mu_b + \sigma_b \odot \epsilon_b)
\end{equation}

\paragraph{Dueling Network Architecture.}
Compared to a single network in deep Q-networks, the dueling architecture \cite{Dueling2022} branches to 2 streams for state value estimation $V(s)$ and advantage estimation $A(s, a)$:
\begin{equation}
Q(s, a; \theta, \alpha, \beta) = V(s; \theta, \beta) + \left( A(s, a; \theta, \alpha) - \frac{1}{|\mathcal{A}|} \sum_{a'} A(s, a'; \theta, \alpha) \right)
\end{equation}
This separation improves policy evaluation when many actions carry similar value.

An exploration bonus is added to the Q-values. The bonus is based on the information gathered by garnering territory, and identifying opponent piece identity.

\begin{equation}
Q_{augmented}(s, a) = Q_{base}(s, a) + \lambda \cdot U(a)
\end{equation}
Here $U(a)$ represents the entropy of the learned distribution at the target location. This term encourages the agent to investigate unknown opponent pieces.

\subsection{Training Stability via Regularization}

The recent study of the imperfect information game, Ataraxos reached superhuman performance \cite{Ataraxos2025}. Its designers found that training stability matters as much as the algorithm itself. Deep reinforcement learning agents face two recurring problems which is wasting updates on low-information experiences, a waste of sampling and training, and cycling between strategies without settling on a robust one, non-stationary.

\paragraph{Advantage Filtering.}
Advantage filtering hand picks the most informative transitions to learn from \cite{RegularizationStrategyExploration}. In reinforcement learning, the temporal-difference (TD) error measures the gap between predicted values and outcomes. A standard replay buffer samples past transitions at random and stores it in the buffer, treating every move equally. Advantage filtering instead samples the buffer and retains only the transitions with the largest TD errors:
\begin{equation}
\mathcal{B}_{filtered} = \text{Top}_{n/k}\left(\mathcal{B}_{sampled}, |\delta_i|\right)
\end{equation}
where $\delta_i = r_i + \gamma^n V(s'_i) - V(s_i)$ is the $n$-step TD error for transition $i$.

\section{AAREN and Learned Embedding System}

The AAREN module \cite{Aaren} will serve as the memory and history system within MARQ. It produces learned embeddings from observations that encode the inferences from those observations. AAREN handles long sequences, retaining details over hundreds of moves while avoiding the gradient problems that degrade older models such as LSTMs. An attention mechanism lets AAREN look back across the full sequence of moves at once \cite{Aaren}.

\subsection{Recurrent Attention Computation}
AAREN implements an attention-based recurrent computation that maintains a state triplet $(a_t, c_t, m_t)$ across the observation sequence. Traditional recurrent neural networks suffer from vanishing gradients over long sequences. AAREN avoids this entirely by providing a direct access to the full history through attention weights while maintaining $O(1)$ inference complexity per timestep.

Given a sequence of observations $O = \{o_1, \ldots, o_t\}$ where each $o_i$ encapsulates the visible board state and public action outcomes, the AAREN cell first projects each observation into key and value representations:

\begin{equation}
k_t = W_k(o_t), \quad v_t = W_v(o_t)
\end{equation}

The projection matrices $W_k$ and $W_v$ are learned during training. The key $k_t$ determines how relevant the current observation is for belief updates. The value $v_t$ contains the content to be integrated into the belief representation. The attention score $s_t$ is a numerical representation of each step in the sequence relative other steps and to the learned query vector $q$:

\begin{equation}
s_t = q^T k_t
\end{equation}

The accumulator $a_t$ aggregates value representations across the 1 episode sequence:

\begin{equation}
a_t = a_{t-1} \exp(m_{t-1} - m_t) + v_t \exp(s_t - m_t)
\end{equation}

The term $\exp(m_{t-1} - m_t)$ adjusts the accumulator when a new maximum appears. The term $\exp(s_t - m_t)$ computes the relative attention weight for the current observation. 


\subsection{Embedding Updates}

When a battle reveals a piece's true type, the AAREN embeddings is updated with ground truth. The inferred distribution at the revealed position collapses to certainty:
\begin{equation}
\text{Representation}_{pos}[t] = \begin{cases}
1.0 & \text{if } t = g \\
0.0 & \text{otherwise}
\end{cases}
\end{equation}

The system updates the revealed piece count and adjusts remaining beliefs every time a revelation happens. For unrevealed positions $p \neq pos$:
\begin{equation}
\mathcal{B}_p[t] \leftarrow \mathcal{B}_p[t] \cdot \frac{\text{max\_count}[t] - \text{revealed\_count}[t]}{\text{remaining\_count}[t]}
\end{equation}


\section{Multi-Agent Learning and Strategy Mixing}

\subsection{League-Based Training Architecture}
The league has a changing pool of agents so that the model encounters many different strategies \cite{perolat2022stratego}. The main agent is the current policy under training. The historical pool stores earlier versions of the main agent, each representing a different stage of development. Opponents are sampled during training according to a fixed probability distribution, to have less bias on the opponents fought.
\begin{equation}
P(\text{opponent}) = 
\begin{cases} 
0.5 & \text{Main Agent (Self-Play)} \\
0.5 & \text{Uniform}(\text{Historical Pool})
\end{cases}
\end{equation}
Preventing the agent from cycling through the same strategies repeatedly. Reaching mastery of the past versions of itself.

\subsection{Strategy Mixing for Exploitation Resistance}

Deterministic strategies are suboptimal in imperfect information games where opponents have varying strategies \cite{perolat2022stratego, MARLGametheoretic}. An adversary who knows the deterministic policy of the opponent reduces the game to a perfect information setting enabling the perfect strategy for an opponent, eliminating the strategic uncertainty that defines games like Stratego. MARQ counters this at two levels.

At the action selection level, Noisy Linear layers inject learned noise into the policy. The network learns the appropriate degree of randomization for each game state rather than relying on fixed exploration rate such as $\epsilon$-greedy.

At the strategy level, diverse pool of opponent exposure during training achieves a more generalized learning. The league system maintains a population of historical agent checkpoints, each representing a different strategic approach discovered during training. 

\subsection{Reward Structure}

The reward function $R(s, a, s')$ guides learning and shaping agent goals in Stratego's long horizon environment \cite{PBRSModern2022}:

\begin{equation}
R_{total} = R_{move} + R_{capture} + R_{tactical} + R_{info}
\end{equation}

The move reward $R_{move}$ gives small bonuses for advancing pieces (approximately $+0.05$) to encourage the agent to prioritize the territory gain and control. The capture reward $R_{capture}$ scales by the rank of the captured piece according to $0.15 \times \frac{\text{Rank}}{11.0}$. The tactical reward $R_{tactical}$ assigns specific bonuses for key strategic captures: $+0.5$ for Miner vs Bomb defusals and $+1.0$ for Spy vs Marshal assassinations. The information reward $R_{info}$ provides $+0.3$ for revealing high-value opponent pieces.


\section{Game Environment and State Management}

The Stratego environment handles board state management, action validation, and observation filtering for agent training and evaluation.

\subsection{Information Set Management}
The environment maintains distinct representations for each information context:

\begin{equation}
\mathcal{I}_{player} = \{s : s \text{ is consistent with player observations}\}
\end{equation}

Each player receives only the subset of game state information available through its observation function.

\subsection{Temporal State Evolution}
Game states evolve according to an update function.

\begin{equation}
s_{t+1} = \text{EnvStep}(s_t, a_t^1, a_t^2, \theta_t)
\end{equation}

Here $\theta_t$ represents stochastic elements such as hidden information revelation and time limits.

\subsection{Valid Action Filtering}
The environment enforces game rules through valid action constraints:

\begin{equation}
\mathcal{A}_{valid}(s, player) = \{a : \text{IsLegalMove}(a, s, player) \land \text{RespectsGameRules}(a, s)\}
\end{equation}

Only legal moves consistent with Stratego's rules and piece capabilities are available to the agent.


\section{MARQ Move Selection and Decision-Making}

Algorithm~\ref{alg:marq_decision} outlines the MARQ move selection process and loop. The first four lines establish the information context. AAREN runs over the public observation sequence to produce an updated history embedding $\mathcal{H}_t$, which is then projected through a softmax layer to yield a probability distribution $\mathcal{B}_t$ over the possible identities of each unrevealed opponent piece. The visible board state $S$ and the history embedding $\mathcal{H}_t$ are concatenated into a single 79-channel tensor, and the Rainbow DQN processes that tensor to produce a Q-value distribution over all actions.

The algorithm then diverges based on move type. Non-attack moves take $Q(s, a_i)$ directly as their final score, since those moves carry no contact uncertainty. Attack moves receive an additional correction. The agent queries $\mathcal{B}_t$ at the target square, retrieves the rank-probability vector $P$, and computes an expected utility $V_{exp}$ by summing over all possible opponent ranks weighted by the outcome of each encounter. The final score blends the network's learned estimate with this hand-computed expectation through the mixing coefficient $\lambda$. A high $\lambda$ trusts the DQN; a low $\lambda$ defers to the probabilistic expectation. This separation is necessary because the DQN's value estimate for an attack encodes long-run strategic considerations, while $V_{exp}$ captures the immediate piece-trade arithmetic more explicitly.

The final selection step depends on whether the agent is in training mode and whether the temperature $T$ is above zero. During self-play, a Boltzmann distribution over $V_{final}$ governs sampling, which injects controlled stochasticity and prevents the agent from cycling into a fixed attack pattern that opponents can predict and counter. At evaluation, the temperature is set to zero and the agent commits to the greedy $\arg\max$ over legal actions.

\begin{algorithm}
\caption{MARQ: Decision-making and Move Selection}
\begin{algorithmic}[1]
\REQUIRE Public Observations $O$, Board State $S$, DQN weights $\theta$, AAREN weights $\phi$, Temperature $T$
\ENSURE Optimal Action $a^*$
\STATE $\mathcal{H}_t \gets \text{UpdateAggregation}(O, \phi)$ 
\STATE $\mathcal{B}_t \gets \text{GetSoftmaxBeliefs}(\mathcal{H}_t)$ 
\STATE $X_t \gets \text{Concatenate}(S, \mathcal{H}_t)$ 
\STATE $Q(s, \cdot) \gets \text{RainbowInference}(X_t, \theta)$ 
\STATE $\mathcal{A} \gets \text{LegalMoves}(S)$
\FORALL{$a_i \in \mathcal{A}$}
    \IF{$\text{IsAttack}(a_i, S)$}
        \STATE $target\_pos \gets \text{TargetSquare}(a_i)$
        \STATE $P \gets \mathcal{B}_t(target\_pos)$ 
        \STATE $V_{exp} \gets \sum_{r \in \mathcal{R}} P(r) \cdot \mathbb{U}(\text{Attacker}, r)$ 
        \STATE $V_{final}(a_i) \gets \lambda Q(s, a_i) + (1-\lambda) V_{exp}$
    \ELSE
        \STATE $V_{final}(a_i) \gets Q(s, a_i)$
    \ENDIF
\ENDFOR
\IF{Training/Self-Play and $T > 0$}
    \STATE $\pi(a_i) \gets \frac{\exp(V_{final}(a_i) / T)}{\sum_{a \in \mathcal{A}} \exp(V_{final}(a) / T)}$
    \STATE \textbf{return} $a^* \sim \pi(\cdot)$
\ELSE
    \STATE \textbf{return} $a^* \gets \arg\max_{a_i \in \mathcal{A}} V_{final}(a_i)$
\ENDIF
\end{algorithmic}
\label{alg:marq_decision}
\end{algorithm}



\section{Generalizability}

\subsection{Mathematical Convergence Properties}
\subsubsection{Incremental Learning Bounds}
The algorithm bounds value function estimation error as follows:
\begin{equation}
\|V_t - V^*\|_\infty \leq \frac{C}{\sqrt{t}} + \epsilon_{approx}
\end{equation}
where $C$ depends on reward bounds and the discount factor, and $\epsilon_{approx}$ represents function approximation error.

\subsubsection{Implicit Representation Accuracy Bounds}
Representation accuracy improves with additional observations:
\begin{equation}
\mathbb{E}[\text{Accuracy}(\mathcal{B}_t)] \geq 1 - \exp(-\alpha \cdot N_{obs}(t))
\end{equation}
where $N_{obs}(t)$ is the cumulative number of informative observations and $\alpha$ is a positive constant.

\subsubsection{Strategy Convergence in Self-Play}
The league training methodology converges to $\epsilon$-Nash equilibria under standard monotonic improvement assumptions:
\begin{equation}
\lim_{T \to \infty} \| \pi_T - \pi^* \| \leq \epsilon
\end{equation}
for some $\epsilon > 0$ depending on the exploration parameters and opponent diversity. The convergence depends on the diversity of the opponent pool and the promotion thresholds that control which agents enter the main pool.

The experimental results in Chapter 5 are consistent with these bounds. The Rainbow agent's loss continued to decrease over 75,000 episodes (Finding 13), and the win rate in self-play reached 49.88\% (Finding 7), though the system was not run long enough to test convergence to approximate Nash equilibrium.

%\section{AI Transparency, Inclusion, and Originality of Intent}

%The development of the MARQ framework for Stratego was governed by intentionality and human-led architectural design. While large language models and generative AI have become ubiquitous in modern software development, their role in this research was strictly limited to tasks such as grammatical refinement, high-level brainstorming, and code debugging for modules and networking. Every core component of the MARQ system, including the AAREN history aggregator and the ResNet-based policy heads, was researched, evaluated, and implemented with the intent of mastering imperfect information games in the Stratego domain.

\section{Algorithmic Evolution and Architecture Ablations}

The final MARQ architecture emerged from systematic experimentation and the refinement of several reinforcement learning components. Multiple established algorithms were evaluated and replaced because of performance bottlenecks or insufficient strategic depth in the Stratego domain.

\subsection{Value Optimization and Exploration.} 
Early iterations used a normal DQN with $\epsilon$-greedy exploration. This approach proved to be a challenge early on for Stratego's sparse rewards and led to a desync of loss to seing model improvement though winrate and reward accumlated. The implementation was upgraded to a Rainbow DQN architecture incorporating Noisy Networks for a more dynamic approach to exploration. The standard DQN configuration was then replaced with a Dueling Head architecture, which improved value estimation accuracy in states where multiple actions carry similar utility. To mitigate overfitting to deterministic opponent policies during self-play, the system also employs \textit{Boltzmann Exploration} (Temperature Scaling). The opponent's action selection uses a temperature-scaled softmax over Q-values, injecting controlled stochasticity to produce a more robust training curriculum.

\subsection{History and Memory Representation.} 
Initial attempts to manage hidden information in Stratego relied on explicit Probabilistic Belief States, a direct representation using another dqn to estimate the piece value. These states introduced prohibitive computational complexity and memory overhead. They required $O(N^2)$ updates that throttled training throughput. The framework moved to an implicit belief system powered by the AAREN architecture. Long Short-Term Memory cells for trajectory embeddings were also replaced by AAREN's attention-based recurrence, which provided more stable gradient flow over long sequences and better handled the non-Markovian nature of piece identity inference.

\subsection{Architectural Refinements.} 
The view of the agent and how it sees the enviroment had several iterations, moving from standard deep convolutional stacks to a dedicated ResNet architecture. Testing identified 6 Residual Blocks as the optimal depth. This configuration provides full receptive field coverage of the 10x10 board while maintaining low inference latency. The optimizer was transitioned from standard ADAM to ADAMW to improve weight decay and policy loss stability.

\subsection{Strategic Initialization and Efficiency.} 
The piece placement phase was initially modeled using an agentic SetupAgent trained through deep reinforcement learning. This approach introduced exploitability risks and computational overhead. The project transitioned to a non-agentic, fixed Heuristic Setup Agent with randomized flag positioning. The change resolved the exploitability concern and provided substantial computational savings. The agentic setup required approximately 2 seconds per episode; the heuristic method completes in under 100 milliseconds. This optimization effectively halved the total training compute requirement. The architectural iterations and the theoretical justifications are systematically documented, linking each component's theoretical necessity directly to empirical validation.

\chapter{Methodology}\label{chap:methodology}
This chapter presents the development and evaluation of the MARQ framework, covering the system architecture, training procedure, and the metrics used to compare agent performance against heuristic and self-play baselines.

\section{Project Planning}
This research followed a rapid prototyping method for the design and construction of the MARQ framework. Development ran from September 2025 to February 2026 and covered building the game environment and training the agent. Algorithmic evaluation and ablation testing ran from March to April 2026.


\section{System Development}
System development consists of two components: the game environment and the MARQ framework. The game environment provides the simulation in which agents operate and train. The MARQ framework produces the trained agent whose performance is evaluated against heuristic baselines and through self-play.

\begin{figure}[htbp]
    \centering
    \includegraphics[width=1.1\textwidth]{figures/MARQArchi.png}
    \caption{MARQ framework architecture.}
    \label{fig:marq_architecture}
\end{figure}

\subsection{Piece Value Computation}
The reward function requires a quantitative assessment of piece value to guide learning. This section presents an analytical method for computing piece values based on probability theory, as a deterministic alternative to simulation-based estimation.

\subsubsection{Combat Dominance Score}
The Combat Dominance Score (CDS) represents the probability that a given piece type defeats a uniformly random opponent piece. Let $\mathcal{P}$ denote the set of all piece types and $n_i$ the count of piece type $i$ in standard Stratego. The CDS for piece $p$ is defined below.

\begin{equation}
\text{CDS}(p) = \frac{\sum_{i \in \mathcal{P}} \mathbb{1}_{p \succ i} \cdot n_i}{N}
\end{equation}

\noindent where $N = \sum_{i \in \mathcal{P}} n_i = 40$ is the total piece count, $\mathbb{1}_{p \succ i}$ is an indicator function equal to 1 if piece $p$ defeats piece $i$ according to Stratego combat rules, and 0 otherwise. The combat rules specify that higher-ranked pieces defeat lower-ranked pieces, with specific exceptions. The Spy defeats the Marshal when attacking, and the Miner defeats Bombs.

\subsubsection{Survival Rate}
The Survival Rate (SR) measures the expected probability of a piece surviving an attack from a uniformly random attacker. Let $\mathcal{M} \subset \mathcal{P}$ denote the set of movable piece types (excluding Flag and Bomb). The SR for piece $p$ is calculated as follows.

\begin{equation}
\text{SR}(p) = \frac{\sum_{a \in \mathcal{M}} \mathbb{1}_{p \succeq a} \cdot n_a}{\sum_{a \in \mathcal{M}} n_a}
\end{equation}

The indicator function equals 1 if piece $p$ survives an attack from piece $a$.

\begin{table}[h]
\centering
\caption{Computed piece values with combat and survival metrics.}
\label{tab:piece_values}
\begin{tabular}{lccc}
\hline
\textbf{Piece} & \textbf{Value} & \textbf{CDS} & \textbf{SR} \\
\hline
Marshal & 8.4 & 0.825 & 0.939 \\
General & 8.0 & 0.800 & 0.939 \\
Colonel & 7.5 & 0.750 & 0.879 \\
Major & 6.7 & 0.675 & 0.788 \\
Captain & 5.7 & 0.575 & 0.667 \\
Lieutenant & 4.8 & 0.475 & 0.545 \\
Miner & 4.3 & 0.400 & 0.273 \\
Sergeant & 3.8 & 0.375 & 0.424 \\
Bomb$^{\dagger}$ & 2.0 & --- & --- \\
Spy & 1.1 & 0.050 & 0.000 \\
Scout & 1.0 & 0.050 & 0.030 \\
Flag$^{\ddagger}$ & 0.1 & --- & --- \\
\hline
\multicolumn{4}{l}{\footnotesize $^{\dagger}$Immovable defensive piece; eliminates any non-Miner attacker.} \\
\multicolumn{4}{l}{\footnotesize $^{\ddagger}$Immovable objective piece; capture results in game loss.} \\
\end{tabular}
\end{table}

\subsubsection{Auxiliary Value Components}


Certain pieces have capabilities that other than their combat rank. Table \ref{tab:strategic_bonus} enumerates these bonuses.

\begin{table}[h]
\centering
\caption{Strategic ability bonuses for special-capability pieces.}
\label{tab:strategic_bonus}
\begin{tabular}{lcp{6cm}}
\hline
\textbf{Piece} & \textbf{Bonus} & \textbf{Description} \\
\hline
Miner & 2.0 & Capable of removing Bombs \\
Spy & 1.5 & Defeats Marshal when initiating attack \\
Scout & 1.0 & Multi-square movement capability \\
Marshal & 0.5 & Highest combat rank \\
\hline
\end{tabular}
\end{table}


\subsection{AAREN and Game History}
The challenge in Stratego is the partial observability. The framework uses the same architectural approach as DeepNash. In deepnash the technique used was a U-Net Convolutional Torso with Residual Blocks and Skip-Connections \cite{perolat2022stratego}. The Attention-as-a-Recurrent-Neural-Network (AAREN) module aggregates history and generates embeddings which then is passed to the CNN and Resnet.

\paragraph{Action Feature Extraction.}
For each step observed opponent action, AAREN extracts a 24-dimensional feature vector encoding movement features, position features, behavioral indicators, and temporal features.


\section{Model Training}

\begin{figure}[htbp]
    \centering
    \includegraphics[width=1.1\textwidth]{figures/Training Architecture.drawio.png}
    \caption{MARQ training architecture.}
    \label{fig:training_architecture}
\end{figure}

Training the MARQ agent is computationally intensive due to the large state-action space and the need for accurate belief state representations over extended game horizons.

\paragraph{Rainbow DQN Architecture.}
The MARQ system uses the six components of Rainbow DQN \cite{RainbowJMLR2021}. Categorical DQN (C51) \cite{DistributionalRL2023} handles value distribution estimation. Noisy Networks \cite{Yang2023Exploration} drive exploration. A Dueling Network architecture \cite{Dueling2022} decouples state values from action advantages. Double DQN chooses the bootstrap action with the online network and takes its value from the target network. Prioritized experience replay samples transitions in proportion to their loss, and $n$-step returns ($n = 5$) carry rewards several moves back in one update. The noisy layers introduce state-dependent stochasticity, removing the need for manual epsilon-greedy scheduling.

\begin{figure}[htbp]
    \centering
    \includegraphics[width=0.8\textwidth]{figures/Dueling Heads.png}
    \caption{Dueling network heads.}
    \label{fig:dueling_heads}
\end{figure}

\paragraph{Network Structure.}
The Q-network features a ResNet backbone for spatial feature extraction. An initial convolutional layer feeds into six residual blocks. This configuration provides a sufficient receptive field to evaluate global board state interactions.

\begin{figure}[htbp]
    \centering
    \includegraphics[width=0.8\textwidth]{figures/Resnet.png}
    \caption{ResNet backbone.}
    \label{fig:resnet_backbone}
\end{figure}

A spatial self-attention layer follows the ResNet backbone. It captures relationships between pieces across the full board, even at long range. This layer feeds into the dueling heads, allowing the network to weight distant piece relationships in value estimation.

\begin{figure}[htbp]
    \centering
    \includegraphics[width=0.8\textwidth]{figures/Input layer.png}
    \caption{79-channel input representation.}
    \label{fig:input_layer}
\end{figure}

The 79-channel state representation is partitioned into two functional sets, as summarized in Table \ref{tab:state_channels}.

\begin{table}[h]
\centering
\caption{State representation input channels.}
\label{tab:state_channels}
\begin{tabular}{clp{8cm}}
\hline
\textbf{Channel} & \textbf{Category} & \textbf{Description} \\
\hline
\multicolumn{3}{l}{\textit{Board Features (Channels 0--14)}} \\
0--11 & Own Pieces & Binary spatial maps indicating positions of each piece type (Flag, Spy, Scout, Miner, Sergeant, Lieutenant, Captain, Major, Colonel, General, Marshal, Bomb) \\
12 & Enemy Mask & Binary map indicating positions of all visible enemy pieces \\
13 & Obstacles & Binary map of lake squares (impassable terrain) \\
14 & Empty Squares & Binary map of unoccupied traversable positions \\
\hline
\multicolumn{3}{l}{\textit{AAREN Embedding Features (Channels 15--78)}} \\
15--78 & Learned Embeddings & 64-dimensional learned embedding vectors per board position, produced by the AAREN history aggregator from observed opponent action sequences. AAREN learns them from a supervised loss on revealed pieces; the Q-network reads them as inputs. \\
\hline
\end{tabular}
\end{table}

The board feature channels provide deterministic knowledge of the agent's own piece positions and observed enemy locations. The AAREN embedding channels encode learned representations of hidden enemy piece identities based on observed behavior. These embeddings are dense learned vectors, not explicit probability distributions. AAREN learns them by predicting the type of an enemy piece when a battle reveals it. The Q-network learns how to use the embeddings, but its loss does not train AAREN (see the Gradient Bridge below). The AAREN embedding channels are populated in all training phases; the agent always receives embeddings from the history aggregator rather than ground truth one-hot encodings. The board feature channels carry different levels of enemy information depending on the observability setting. During full observability (Phase 1), the board tensor reveals all enemy piece identities, providing a strong signal in the piece-specific board channels. During partial observability phases, unrevealed enemy pieces appear only as a ``hidden piece'' presence marker (affecting Channel 12, the enemy mask), and the agent must rely on AAREN embeddings to infer their true identities. An earlier version suffered from a bug where the hidden piece marker was filtered out of Channel 12, blinding the network to the presence of enemy pieces. 
\paragraph{Training Hyperparameters.}
The training configuration specifies a learning rate $\alpha = 3 \times 10^{-5}$ with AdamW optimizer and weight decay $\lambda = 0.01$, and a discount factor $\gamma = 0.995$ to capture long-term strategic consequences. A batch size of 512 transitions and a replay buffer capacity of 150,000 transitions with prioritized sampling are used. Soft target network updates ($\tau = 0.001$) occur every 5,000 steps. Model updates run every 2 steps across 8 parallel lanes with 5-step returns for deeper credit assignment in long-horizon games. The AAREN history aggregator retains up to 50 actions per position in its training buffer to preserve early-game movement patterns across 200 to 500 move games. Turn normalization is aligned with the maximum game length of 1,000 turns. Horizontal flip augmentation exploits the game's symmetry to improve sample efficiency. Training proceeds for 35,000+ episodes with checkpoints saved every 1,000 episodes.

\paragraph{Configuration after the thesis experiments.} The settings in this chapter are those of the experiments reported in Chapter~\ref{chap:results}. The training system was revised after them (Appendix~\ref{app:lessons}). The later formal runs use a 30-plane observation (32 from formal\_6 on, with two planes marking the hidden enemy pieces that have moved; 41 from formal\_9 on, with nine battle-outlook planes that give each attack's chances of winning and losing) with a lossless 3,600-action encoding and fog of war in every stage. They use $\gamma = 0.999$, a batch of 128, a replay capacity of 75,000 and a hard copy into the target network every 1,000 updates, with a learning rate of $10^{-4}$ from formal\_10 on. From formal\_11 on the shaping potential counts the enemy's material exactly from public information, and from formal\_12 on it adds a term for the number of movable pieces. They run 16 parallel lanes and start from heuristic demonstrations. In their first phase the AAREN planes are switched off, so those runs test the Rainbow learner without AAREN. Their results are reported in Appendix~\ref{app:lessons}, not in Chapter~\ref{chap:results}.

\paragraph{Computational Optimizations.}
Mixed-precision inference accelerates training throughput. During action selection, network forward passes operate in FP16 precision via \texttt{torch.amp.autocast}, reducing memory bandwidth requirements and leveraging Tensor Core acceleration on NVIDIA GPUs.

Integrating the recurrent AAREN module into the replay pipeline presents a significant bottleneck. Full sequential reconstruction of action histories during replay increases iteration time from 15s to 60s. To maintain high throughput, MARQ stores pre-computed embedding snapshots in the replay buffer, detached from AAREN's computation graph. This gave a fourfold speedup. The thesis code also added a ``Gradient Bridge'' to the stored snapshots during replay: an epsilon-scaled projection of a zero vector, $\epsilon \cdot W_{proj}(\mathbf{0})$, meant to give the DQN loss a gradient path back to AAREN. The projection is affine, so $W_{proj}(\mathbf{0}) = \mathbf{b}_{proj}$ and the term is the layer's bias alone. The DQN loss therefore reached only one of AAREN's 21 parameter tensors, and none that depends on the move history. In the thesis experiments AAREN was trained, in effect, only by its supervised signal from piece reveals. The current code drops the bridge and trains AAREN by that supervised loss alone (Appendix~\ref{app:lessons}).

\paragraph{Multi-Phase Curriculum.}
Training follows a multi-phase curriculum. Phase 1 establishes foundational rules under full observability. As win rates stabilize, the curriculum introduces mixed observability, requiring the agent to use AAREN-based inference before full information occlusion is enforced.

Phase 1 runs for 5,000 to 10,000 episodes against a progressively challenging opponent mix: initially random opponents, then basic heuristic agents, and finally a multi-factor heuristic opponent. Phase advancement requires $\geq$95\% win rate against random opponents, $\geq$75\% against the basic heuristic, and $\geq$60\% against the strategic heuristic. All win-rate metrics are computed relative to decisive games (wins and losses) rather than total episodes to account for the high frequency of early-stage draws.

Phase 2 (Partial Observability) disables full observability and requires reliance on AAREN-generated embeddings for 5,000 to 10,000 episodes. The opponent distribution allocates 40\% to frozen heuristic agents and 60\% to the strategic heuristic. Phase advancement requires embedding-based inference accuracy $\geq$75\% and an overall win rate $\geq$60\% relative to decisive games.

Phase 3 (Self-Play) trains against a mixture of delayed agent copies (60\%) and strategic heuristic opponents (40\%) for 5,000 to 15,000 episodes to discover novel strategies while preventing policy collapse. Phase 4 (League Training) activates the full opponent distribution with 40\% historical league agents, 30\% strategic heuristic, 20\% greedy agents, and 10\% self-play. This phase also activates independent parallel-learning (MARL) for Agent 2, transitioning the opponent from a fixed policy pool to a continuously updating student agent. Phase 5 (Scenario Drills) introduces specialized board configurations every 1,000 episodes for targeted skill training.


\subsection{Test-Time Search}

The trained Q-network provides a strong action-selection policy through single-pass inference. Tactical performance in partially observable environments improves with a formal \textit{Expectamax} search, originally proposed by Michie \cite{novikov2024analysis} and recently analyzed in complex game environments by Novikov and Yanovskyi \cite{novikov2024analysis}. This algorithm extends standard minimax by incorporating stochasticity through ``Chance Nodes'' that represent the hidden identity of opponent pieces.

The search refinement uses the piece-identity probabilities generated by the AAREN history aggregator to calculate expected utility of potential moves. For a given state $s$ and a set of legal moves $\mathcal{A}(s)$, the refined value $V^*(a_i)$ is computed as follows:

\begin{enumerate}
    \item \textbf{Max Node evaluation.} The agent evaluates each candidate move $a_i$ using its positional prior $Q_{dqn}(s, a_i)$.
    \item \textbf{Chance Node summation.} If $a_i$ results in an attack on position $j$, the identity of the target piece is treated as a chance node. The expected combat utility $\mathbb{E}[U_{combat}]$ sums the utility of all possible outcomes, weighted by their probabilities:
    \begin{equation}
        \mathbb{E}[U_{combat}] = \sum_{r \in \mathcal{R}} P(\text{rank}_j = r | \mathcal{H}) \cdot \text{Utility}(rank_{attacker}, r)
    \end{equation}
    where $P(\text{rank}_j = r | \mathcal{H})$ is the AAREN-derived probability that the target piece at position $j$ has rank $r$, and $\mathcal{H}$ denotes the observed interaction history.
    \item \textbf{Move Selection.} The final decision combines the positional prior and the expected tactical outcome:
    \begin{equation}
        a^* = \arg\max_{a_i \in \mathcal{A}(s)} \left( \lambda \cdot Q_{dqn}(s, a_i) + (1-\lambda) \cdot \mathbb{E}[U_{combat}] \right)
    \end{equation}
\end{enumerate}

The Expectamax logic prevents the agent from committing high-value pieces to squares with high probabilities of containing Bombs or superior ranks, even when the raw Q-network output is optimistic about the board position. The exact summation over rank probabilities grounds the risk assessment in observed opponent behavior.

\subsection{Deployment}
The best-performing MARQ model weights are saved after training and validation. During evaluation, exploratory noise is deactivated, and the model selects the action with the highest predicted Q-value. Each evaluation run uses the same checkpoint to maintain consistency across test conditions.

\begin{figure}[H]
    \centering
    \includegraphics[width=.7\textwidth]{figures/Game diagram.png}
    \caption{Game sequence diagram.}
    \label{fig:sequence_diagram}
\end{figure}

The game loop runs until one agent captures the opposing flag. On each turn, the active agent receives the current board state and updates the AAREN module, which produces learned embeddings from observed action history. The DQN evaluates Q-values for all legal actions given this embedding-enhanced state representation and selects the action with the highest expected value.


\section{Game Environment}
The game environment is a digital implementation of Stratego developed in Python and Pygame. It manages the board state, enforces game rules, handles combat resolution, and provides the observation interface for agent perception. All agent-vs-agent experiments in this study use this environment.

\subsection{Evaluation}
The evaluation tracks learning progress and agent performance through agent-vs-agent testing. Win rate measures the proportion of decisive games won against each opponent tier. Average game length, computed as the mean number of moves per game, indicates decision efficiency. Loss curves show the progression of distributional and value losses over episodes. They show whether the critic fits its own targets, not whether the policy improves, so they are read next to the win rates rather than in their place. The Rainbow DQN agent is compared with a standard DQN baseline. The two configurations differ in more than one component at a time, so the comparison shows the effect of the combination. Isolating each component needs an ablation that removes one at a time, which has not yet been run (Chapter~\ref{chap:recommendations}). Empirical findings are organized thematically into designated subsections to clearly connect each architectural modification to its resulting behavioral outcome.

Developing the MARQ framework is iterative. Each version represents an improvement over the last. Appendix B shows the initial results from 1,000 episodes. During this early stage, the agent is still exploring strategies. The early win rates show many draws as pieces move without decisive engagement. With Noisy Networks, the agent learns how much to explore on its own, adapting better as training continues.

%%%%%%%%%%%%%%%%%%%%%%%%%%%%%%%%%%%%%%%%%%%%%%%%%%%%%%%%%%%%%%%%%%%%%%%%%%%
\chapter{Results and Discussion}\label{chap:results}
%%%%%%%%%%%%%%%%%%%%%%%%%%%%%%%%%%%%%%%%%%%%%%%%%%%%%%%%%%%%%%%%%%%%%%%%%%%

\section{Normal DQN Performance Ceiling}

The multi-phase curriculum concluded with 37,213 episodes of self-play training. The normal DQN agent used an LSTM-based history encoder and trained against random, greedy heuristic, league, and self-play opponents. Table~\ref{tab:vanilla_performance} reports the results.

\begin{table}[htbp]
    \centering
    \caption{Normal DQN performance after 37,213 episodes of self-play training.}
    \label{tab:vanilla_performance}
    \begin{tabular}{|l|c|c|c|c|}
        \hline
        \textbf{Opponent Type} & \textbf{Episodes} & \textbf{Wins} & \textbf{Win Rate} & \textbf{Draw Rate} \\
        \hline
        Greedy Heuristic & 9,319 & 1,418 & 15.2\% & 41.3\% \\
        Self-Play        & 9,565 & 2,370 & 24.8\% & 36.8\% \\
        Random           & 9,294 & 2,260 & 24.3\% & 37.7\% \\
        League           & 9,035 & 2,231 & 24.7\% & 38.2\% \\
        \hline
        \textbf{Aggregate} & \textbf{37,213} & \textbf{8,279} & \textbf{22.2\%} & \textbf{38.5\%} \\
        \hline
    \end{tabular}
\end{table}

The training data contains six findings, each linked to an architectural limitation.

\subsection{Win Rate Plateau}\label{finding:1}
Agent1's cumulative win rate stayed near 22\% after 2,000 episodes, it stayed more or less the same over the next 35,000. The 100-episode rolling average swung between 15\% and 30\% with no deviation.

\begin{figure}[htbp]
    \centering
    \includegraphics[width=0.85\textwidth]{figures/Results/cropped/cumulative_wins_winrate.png}
    \caption{Cumulative win count and win rate over 37,213 episodes. The win rate plateaus early and remains flat, indicating policy stagnation.}
    \label{fig:cumulative_winrate}
\end{figure}

\subsection{Loss--Performance Decoupling}\label{finding:2}
The TD error fell to roughly 0.001 however, the win rate did not follow. The network learned to predict the value of passive, draw-seeking play accurately. It did not discover better strategies.

\begin{figure}[htbp]
    \centering
    \includegraphics[width=0.85\textwidth]{figures/Results/cropped/policy_loss.png}
    \caption{Policy loss (TD error).}
    \label{fig:policy_loss}
\end{figure}

\subsection{Q-Value Overestimation}\label{finding:3}
Average Q-values climbed from 0.2 to 1.2. Win rates did not move with the rising value estimates.

\begin{figure}[htbp]
    \centering
    \includegraphics[width=0.85\textwidth]{figures/Results/cropped/q_value.png}
    \caption{Average Q-values climb.}
    \label{fig:q_value}
\end{figure}

\subsection{LSTM Prediction Accuracy Ceiling}\label{finding:4}
The LSTM piece-identity module accuracy stayed around 18--22\%, a 2.2--2.6$\times$ gain over the 8.3\% random baseline but is still far short of reliable tactical inference. The improvement stopped after episode 5,000 and stayed at that level.

\begin{figure}[htbp]
    \centering
    \includegraphics[width=0.85\textwidth]{figures/Results/cropped/prediction_accuracy.png}
    \caption{Piece-identity prediction accuracy.}
    \label{fig:prediction_accuracy}
\end{figure}

\subsection{Opponent-Specific Performance Asymmetry}\label{finding:5}
Performance varies by opponent given randomly in the training process. Win rate against the greedy heuristic was 15.2\%; in self-play it reached 24.8\%. The policy fit its own behavioral distribution and did not adapt to the others opponent types.

\subsection{Draw Dominance}\label{finding:6}
The game results was mostly draws. Draws accounted for 38.5\% of outcomes, far above the 22.2\% win rate. The agent learned to shuffle pieces passively and avoid uncertain engagements. It minimized short-term losses at the expense of long-term winning probability, lacking the long horizon planning.




\section{Behavioral Patterns from Training Logs}

\subsection{Passive Shuffling Convergence}

The agent grew conservative. Average episode length rose throughout training as it avoided combat and cycled pieces between adjacent squares. Scalar value collapse (Finding~2) explains why the Q-network converges to the lowest-variance action class, playing safe moves instead of winning moves.

\subsection{Opponent-Specific Adaptation Failure}

Table~\ref{tab:vanilla_performance} shows the agent overfitting to its own play style, not generalizing to other agents, since the drop-off of performance against greedy opponents is significant (15.2\%). The greedy heuristic has the lowest win rate (15.2\%) due to its rank-agnostic attack pattern directly punishing passive shuffling. The heuristic set of instructions attacks the nearest piece, forcing repeated engagements that the agent's risk-averse policy is designed to avoid. When forced into combat, the agent unreliable piece-identity inference needed to choose favorable matchups fails, so each engagement is approximately a coin flip weighted by the opponent's rank distribution. The result is a consistent loss rate that the agent win with further passivity.

Self-play produced the highest win rate (24.8\%), but this figure can be misleading. Both agents share the same weights and therefore the same behavioral tendencies. When two equally passive and cautious policies face each other, games resolve through symmetric attrition rather than strategic superiority. The apparent 24.8\% win rate reflects the probability that stochastic tiebreakers favor one copy over another, not genuine tactical competence. The high draw rate in self-play (36.8\%) confirms that most games end in mutual passivity.

Random opponents produced a 24.3\% win rate, nearly identical to self-play despite being the weakest opponent class. A random policy does not exploit the agent's weaknesses, but it also provides no consistent pattern for the agent to learn from. Transitions generated against random play carry high variance and low signal, contributing gradient noise without driving policy improvement. The similar win rates against random and self-play opponents indicate that the agent has not learned to exploit weak play, further evidence of policy stagnation.

League opponents (24.7\%) are a historical snapshots of the agent past versions. Because the agent's policy plateaued early (Finding~1), these snapshots are nearly indistinguishable from the current policy. A league of stagnant policies cannot force adaptation and the fixed $\epsilon$-greedy decay adds to this problem by providing no mechanism for adapting exploration to different opponent since the agent has stopped trying to find different strategy.

\subsection{Loss-Reward Divergence}

The TD loss dropped from 0.05 to 0.001 by episode 37,000. Average episode reward stayed flat near 0.08, a possible confusion in the reward signal. The network learned to predict outcomes of its existing passive play with high precision, knowing that a safe move will net a safe outcome. With the bottoming $\epsilon$-greedy, it did not search for strategies with higher expected returns.

\subsection{LSTM Accuracy}

The piece prediction accuracy hit 20\% after episode 5,000 and stagnated for the rest of the training. One correct inference in five is not enough to evaluate combat outcomes reliably, so the agent tend to avoid combat altogether


\section{Draw Dominance}

Draws reached 38.5\% of all games, and the win rate was only 22.2\%. This rate matters not just as a performance number but as a drag on learning itself, a symptom of a bigger problem.

The scalar value collapse (Finding~2) was first offered as the reason the model played safe rather than well. A scalar Q-value cannot represent a multi-modal outcome distribution, but that does not by itself favour the draw. A mean-based estimator ranks moves by their expected return. A move with a 40\% chance of winning and a 60\% chance of losing has an expected value of about $-0.2$, so it ranks below a certain draw only when the draw is worth more than $-0.2$; against a draw worth $-0.3$ it ranks above it. An estimator that had learned those means correctly would therefore choose the attack. The C51 agent also chooses by the mean of its distribution, so the distributional head would not change this ranking. The draws need another explanation. The later runs point to what a drawn game costs, and to how weakly the outcome of a game reaches the moves that led to it (Appendix~\ref{app:lessons}).

\section{Architectural Comparison: Unified versus Dual-Model Pipeline}

A different approach was use from the initial architecture, which was using two separate networks, a RLNN SetupAgent for piece placement and a DQN for movement. That dual-model design was abandoned, due to the computation overhead and signal interference between the two networks. The current unified pipeline uses a single DQN for all movement and a randomized heuristic for placement. Table~\ref{tab:architecture_comparison} compares both.

\begin{table}[htbp]
    \centering
    \caption{Dual-model versus unified architectural pipelines.}
    \label{tab:architecture_comparison}
    \begin{tabular}{|p{3cm}|p{5.5cm}|p{5.5cm}|}
        \hline
        \textbf{Aspect} & \textbf{Dual-Model (Setup + Attack)} & \textbf{Unified DQN} \\
        \hline
        Architecture & Separate C51-based SetupAgent for placement, separate DQN for attack & Single DQN for movement; heuristic setup with randomized flag \\
        \hline
        Setup Latency & Approximately 2 seconds per episode & Under 100 milliseconds per episode \\
        \hline
        Training Overhead & Doubled compute requirement (setup and attack training phases) & Halved total training compute \\
        \hline
        Exploitability & Deterministic setup patterns learned by opponents & Randomized placement resists opponent exploitation \\
        \hline
        Draw Pattern & Draws due to gradient interference between competing objectives & Draws due to risk-averse policy under uncertainty \\
        \hline
    \end{tabular}
\end{table}

Two lessons emerged from this transition. The SetupAgent did not converged on deterministic placements, randomizing setup while lowering the TD-error of the model a mismatch in result. Placing flag in front of the lineup and trapping big pieces behind the bomb. A randomized heuristic eliminated this vulnerability and cut setup latency by 20$\times$.

The dual-model design also suffered from gradient interference. Setup gradients rewarded defensive formations while attack gradients rewarded aggressive piece use. Optimal defense often places pieces in positions bad for attack. Removing the setup phase from the learning loop resolved this conflict.

Both architectures converged on draw-heavy outcomes. In the dual case, gradient interference between objectives prevented coherent learning. In the unified case, scalar value collapse and poor exploration drove the agent into passive shuffling. The draw problem is architectural and rooted in the long horizon of the game, not an unforeseen accident.


\section{Architectural Bottleneck Analysis}

Each of the findings above can be explained by a known structural property of the DQN architecture. The explanations are hypotheses that the later sections test; they do not rule out other causes, such as the reward design (Appendix~\ref{app:lessons}).

\subsection{Q-Value Overestimation}

Finding~3 is a known property of single-estimator Q-learning. The network selects and evaluates actions with the same weights:

\begin{equation}
y_t = r_t + \gamma \max_{a'} Q_\theta(s_{t+1}, a')
\end{equation}

Noisy estimates plus a max operator produce a positive bias, reducing the ability of the agent to learn from its mistake \cite{RainbowJMLR2021}:

\begin{equation}
\mathbb{E}\left[\max_{a} (Q^*(s, a) + \epsilon_a)\right] \geq \max_{a} Q^*(s, a)
\end{equation}

The bias grows with noise variance and action-space size $|\mathcal{A}|$. Stratego positions can have over in theory a 220 to 240 legal moves, amplifying the effect. Double DQN splits selection from evaluation across two estimators, provably reducing this bias \cite{RainbowJMLR2021}.

\subsection{Scalar Value Collapse}

Finding~2 (loss convergence without performance gain) follows from distributional RL theory. The standard DQN learns a scalar $\mathbb{E}[Z(s,a)]$ and minimizes:

\begin{equation}
\mathcal{L} = \mathbb{E}\left[\left(r + \gamma \max_{a'} Q(s', a') - Q(s, a)\right)^2\right]
\end{equation}

In Stratego the return distribution $Z(s,a)$ is multi-modal \cite{RainbowJMLR2021}: wins near $+1$, losses near $-1$, exchanges in $[-0.5, +0.5]$, draws near $0$. A scalar mean collapses all modes into one number that may not match any real outcome. The loss converges when this mean is estimated correctly, even if the policy cannot tell aggressive high-variance moves from passive low-variance ones.


\subsection{LSTM Information Bottleneck}

Finding~4 follows from sequential compression. The LSTM encodes observations $o_1, \ldots, o_T$ through a 64-dimensional hidden state:

\begin{equation}
h_t = \text{LSTM}(o_t, h_{t-1})
\end{equation}

A Stratego game runs 200 to 500 moves. The hidden state must track belief estimates for 40 opponent pieces across 12 types. A 64-dimensional vector can hold enough information in theory, but every successive nonlinear update risks losing early observations. Information from move $k$ must survive $T - k$ transformations to reach the final representation. Gating helps, but retention still degrades with sequence length \cite{Aaren}.

AAREN sidesteps this bottleneck. It computes weighted aggregations directly over the full history with $O(1)$ per-timestep complexity:

\begin{equation}
h_t = \frac{a_t}{c_t} = \frac{\sum_{k=1}^{t} v_k \exp(s_k - m_t)}{\sum_{k=1}^{t} \exp(s_k - m_t)}
\end{equation}

No sequential compression, no information decay. The network can attend to any prior timestep directly.

\subsection{Exploration Collapse}

The normal DQN relies on a hard-coded $\epsilon$-greedy exploration schedule that decays from 1.0 down to 0.01. Once it hits this floor, the schedule stops adapting. From that point on, the agent explores uniformly at a flat 1\% rate regardless of the game state or its confidence in the current Q-values.

Stratego's game tree has about $10^{535}$ nodes \cite{perolat2022stratego}. A uniform random exploration rate of 1\% is a blunt mechanism that will almost never stumble on the complex sequences of moves required to discover optimal strategies. The agent becomes trapped in local optima because its exploration policy cannot recognize which parts of the state space require further investigation. This limitation necessitates a mechanism like Noisy Networks \cite{RainbowJMLR2021}, which replaces hard-coded randomness with learnable parameters that concentrate exploration where the value function is most uncertain.

\subsection{Sparse Reward Propagation}

The normal DQN cannot develop long-horizon strategic behavior (Findings~1, 5, 6), a limitation compounded by single-step TD learning. Terminal reward information at step $T$ propagates backward through the Bellman backup at a rate set by the discount factor $\gamma$:

\begin{equation}
Q(s_k, a_k) \leftarrow Q(s_k, a_k) + \alpha \left( r_k + \gamma Q(s_{k+1}, a_{k+1}) - Q(s_k, a_k) \right)
\end{equation}

For a decisive game outcome occurring at step $T$ in a game of length 300, a mid-game decision at step 100 receives an effective signal of $\gamma^{200}$. With $\gamma = 0.99$, this yields $0.99^{200} \approx 0.134$, meaning that only 13.4\% of the terminal reward signal reaches mid-game decisions through single-step bootstrapping. Multi-step returns ($n$-step TD) \cite{RainbowJMLR2021} propagate terminal signals $n$ times faster by directly incorporating rewards over $n$ consecutive steps:

\begin{equation}
y_t^{(n)} = \sum_{k=0}^{n-1} \gamma^k r_{t+k} + \gamma^n Q(s_{t+n}, a_{t+n})
\end{equation}

This reduces the effective bootstrap depth and accelerates the propagation of sparse terminal rewards to earlier decision points.


\section{Empirical Validation: Rainbow DQN with AAREN}

The bottleneck analysis above predicted that each failure mode has a specific architectural fix. The following data mining analysis tests those predictions empirically. Three training datasets were analyzed: the same normal DQN at 15,302 episodes, the Rainbow DQN with AAREN at 9,778 episodes of self-play, and a separate extended Rainbow run of 75,047 episodes against heuristic opponents. Table~\ref{tab:arch_comparison} presents the head-to-head comparison. The complete training progress charts appear in Appendix B (Figures~\ref{fig:lstm15k_progress}--\ref{fig:rdqn75k_aaren}).

Three limits apply to every comparison in this section. First, the runs differ in more than the architecture: the opponents, curriculum, reward and run length differ too. Each configuration was also trained once, so a difference cannot be attributed to a single component. Second, the self-play figures come from the training games, in which the agent plays a copy of itself, not from evaluations against fixed opponents. Third, the average Q-value logged by the Rainbow trainer was the maximum over all 3,600 actions. Most of those actions are illegal and never trained, and Appendix~\ref{app:lessons} shows this statistic read about $+2$ while the returns were about $-1$. The Q-values below therefore do not measure the values of the moves that were played.

\begin{table}[htbp]
\centering
\caption{Self-play performance comparison between Vanilla DQN with LSTM and Rainbow DQN with AAREN.}
\label{tab:arch_comparison}
\begin{tabular}{l c c}
\hline
\textbf{Metric} & \textbf{DQN+LSTM (15k)} & \textbf{Rainbow+AAREN (9.5k)} \\ \hline
Total Episodes         & 15,302   & 9,778    \\
Total Steps            & 10,306,433 & 5,169,892 \\
Win Rate (\%)          & 24.97    & 49.88    \\
Draw Rate (\%)         & 50.04    & 6.69     \\
Loss Rate (\%)         & 24.99    & 43.43    \\
Steps per Win          & 2,697    & 1,060    \\
Avg Episode Length     & 1,345    & 1,055    \\
Flag Capture Ratio (\%)& 18.79   & 8.88     \\
Loss--Win Rate Corr.   & 0.051   & $-$0.787 \\
Mean Q-Value           & 0.33    & 0.70     \\
Model Size (per ckpt)  & $\sim$39\,MB & $\sim$433\,MB \\
\hline
\end{tabular}
\end{table}

\subsection{Win Rate }\label{finding:7}
The Rainbow agent won 49.88\% of its self-play games at 9,778 episodes, against 24.97\% for the DQN at 15,302 episodes. In self-play this number is set mostly by the draw rate. When both sides play the same policy, the decisive games split about evenly between them, as they did in both runs: the DQN won 24.97\% and lost 24.99\%, and the Rainbow agent won 49.88\% and lost 43.43\%. The win rate doubled because the share of decisive games nearly doubled (from 50\% to 93\%), not because the agent became stronger than its opponent. Showing that the Rainbow agent plays better than the DQN would need games between the two, or against the same fixed opponents.

\subsection{Draw Elimination}\label{finding:8}
The normal DQN drew 50.04\% of its games, finishing the 1500 moves per agent in an episode. The Rainbow agent drew 6.69\% over 1500 moves. It lost more games outright (43.43\% versus 24.99\%), but it produced definitive outcomes. An agent that draws half of its games has learned to avoid playing passively and conservative in calculating outcomes in the short term. Section~\ref{finding:6} (Draw Dominance) attributed the draws to scalar value collapse. These runs cannot tell which change removed them, because the configurations differ in more than the value head. The effect also did not carry over to the 75k run against heuristic opponents, which drew 77.02\% of its games (Table~\ref{tab:rdqn75k}).

A 43.43\% loss rate alongside a 6.69\% draw rate could suggest that the Rainbow agent takes excessive risks. In self-play, however, every decisive game has one winner and one loser, and both are copies of the same agent. The 93.31\% decisive outcomes mean that the policy forces resolutions rather than stalling. They say nothing about whether its play is reckless or sound against an opponent that is not itself. The self-play games were also shorter (1,055 moves on average, against 1,345 for the normal DQN).

\subsection{Loss--Performance Correlation}\label{finding:9}
The correlation between windowed policy loss and win rate was 0.051 for the normal DQN and $-$0.787 for the Rainbow agent. Two curves that both drift with training time correlate largely through that common drift, whatever their relation. The later formal\_4 run (Rainbow DQN, AAREN off) shows how far this goes. Over blocks of 100 episodes its loss and training win rate correlated at $-0.49$, but the correlation fell to $-0.14$ once a linear trend was removed from both. Its loss fell from 3.63 to 3.05 while its greedy evaluations against fixed opponents did not improve (Appendix~\ref{app:lessons}). The $-0.787$ was not computed on detrended curves, and in self-play the win rate mostly reflects the draw rate (Section~\ref{finding:7}). The coefficient therefore does not show that the Rainbow loss tracks the quality of play. The decoupling of Section~\ref{finding:2} appears in the Rainbow learner as well.

\subsection{Sample Efficiency}\label{finding:10}
The Rainbow agent required 1,060 environment steps per win compared to 2,697 for the normal DQN. At the episode level, the Rainbow agent won once every 2.0 episodes while the normal DQN won once every 4.0, meaning the DQN architecture stumbled its way to a win due to the maximum step is 3000 steps per episode. The Rainbow model consumed 5.17 million total steps; the normal DQN consumed 10.31 million. Steps per win is the total number of steps divided by the number of wins, and in self-play the win rate rose with the share of decisive games (Section~\ref{finding:7}). The metric therefore reflects fewer draws and shorter games more than it reflects learning more from each sample.

\subsection{AAREN and LSTM Accuracy}\label{finding:11}
The AAREN module achieved 24.23\% piece identity inference accuracy, compared to 17.74\% for the LSTM-based PBS module. Three things limit what this gap shows.

First, the 8.3\% baseline of Section~\ref{finding:4} is a uniform guess over 12 types, but the types are not equally common. Scouts are 8 of the 40 pieces (20\%) and 8 of the 33 movable ones (24\%). A predictor that always names the most common type would therefore score about 20 to 24\%, depending on which pieces battles reveal. Until both accuracies are compared with that majority baseline on the same reveals, 24.23\% is not yet evidence of inference from behaviour.

Second, the two accuracies come from different runs, with different opponents, lengths and networks. Third, in the thesis code AAREN read the padding step of every history shorter than the longest in its batch instead of its last real step, and the DQN loss did not train it (Chapter~\ref{chap:methodology}, Gradient Bridge). Section~\ref{finding:4} (LSTM Accuracy Ceiling) predicted that attention would help, but this result neither confirms nor refutes that prediction.

The Q-values cannot be compared either. The LSTM Q-values settled near 0.18 and the Rainbow ones reached 0.84, but the Rainbow trainer's average Q was a maximum over mostly untrained actions.

\subsection{Learned Exploration Active}\label{finding:12}
The NoisyNet entropy rose steadily from 0.034 to 0.049 throughout Rainbow training, adapting to the need for exploration during the training as the agent have not converge to a solution. The normal DQN has a fixed entropy of zero (Section~\ref{finding:5}, Exploration Collapse), due to the constant decay of $\epsilon$-greedy. Noisy Networks replaced frozen $\epsilon$-greedy decay with adaptive, state-dependent exploration.

\subsection{Extended Training at 75,000 Episodes}

A separate Rainbow DQN with AAREN training run reached 75,047 episodes against heuristic opponents. This run used a different curriculum configuration and did not advance to self-play, so it cannot be directly compared to the self-play results above. It does, however, validate two architectural properties over a longer horizon. Table~\ref{tab:rdqn75k} summarizes its metrics.

\begin{table}[htbp]
\centering
\caption{Rainbow DQN with AAREN extended training (75k episodes) against heuristic opponents.}
\label{tab:rdqn75k}
\begin{tabular}{l c}
\hline
\textbf{Metric} & \textbf{Value} \\ \hline
Total Episodes          & 75,047 \\
Win Rate                & 11.66\% \\
Draw Rate               & 77.02\% \\
vs Smart Heuristic WR   & 11.74\% (38,462 games) \\
vs Frozen Heuristic WR  & 11.59\% (43,331 games) \\
Policy Loss (late)      & 2.84 \\
Q-Value (late)          & 1.84 \\
AAREN Accuracy          & 20.78\% \\
Flag Capture Ratio      & 99.96\% \\
Avg Episode Length      & 378 \\
\hline
\end{tabular}
\end{table}

\subsubsection{Loss Convergence Over Long Horizons}\label{finding:13}
The distributional loss fell from 3.89 to 2.84 across 75,000 episodes. The normal DQN showed no comparable long-term loss reduction; its loss either plateaued or inflated in late training. A C51 loss of $\ln 51 \approx 3.93$ is what a prediction spread evenly over the 51 atoms scores, so the run started from an almost uniform prediction. The falling loss shows the critic fitting its own targets more closely. A gradual increase in win rate was seen across the second half of the run, but the run's win rate against heuristic opponents was 11.66\% overall. As in formal\_4 (Section~\ref{finding:9}), the loss can fall while the agent keeps losing, so the falling loss alone is not evidence that extended training improves play.

\subsubsection{Flag Pursuit Behavior}\label{finding:14}
The average episode length in the 75k run settled near 378 steps, compared to 1,345 for the normal DQN and 1,055 for the Rainbow self-play run. The shorter episodes indicate that the agent ended games faster. Of its wins, 99.96\% came by capturing the opponent flag rather than by attrition. The normal DQN won by flag capture only 18.79\% of the time, relying instead on piece depletion. The logged average Q-value grew from near zero to 1.84. With a win worth $+1$, a value of 1.84 at an 11.66\% win rate cannot be the value of the moves played; it is the maximum over mostly untrained actions described above. The attention layer may help the agent locate the flag, but no run has isolated its contribution.

The 75k run is also the more informative measure of strength. In self-play a win rate near 50\% follows from symmetry (Section~\ref{finding:7}). Against fixed heuristic opponents, the agent won about one game in nine (11.66\%) and drew 77.02\% across 75,047 episodes. It did as well against the smart heuristic (11.74\%) as against the frozen one (11.59\%). The later formal runs (Appendix~\ref{app:lessons}) also struggled against these opponents. formal\_4 still lost as many games as it won against a random opponent at episode 3,000, mostly because it attacked revealed pieces it could not beat.

\section{Evidence for Architectural Necessity}

Each component in the thesis title maps to a specific failure mode observed in the normal DQN baseline. Table~\ref{tab:title_evidence} summarizes this mapping and the evidence so far from the Rainbow DQN runs.

\begin{table}[htbp]
    \centering
    \caption{MARQ title components mapped to normal DQN failure modes and the Rainbow DQN evidence so far.}
    \label{tab:title_evidence}
    \begin{tabular}{|p{2.8cm}|p{4.5cm}|p{6.5cm}|}
        \hline
        \textbf{Title Component} & \textbf{Observed Failure Mode} & \textbf{Empirical Evidence} \\
        \hline
        Multi-Agent & Symmetric MARL Nash lock-in & P1 and P2 win rates were 25.1\% and 24.9\% in the normal DQN. In Rainbow self-play they were 49.88\% and 43.43\%, still close, as self-play implies; what changed is that draws fell from 50\% to 6.69\% (Section~\ref{finding:7}). \\
        \hline
        Rainbow DQN & Q-value overestimation, scalar value collapse, exploration death, sparse reward propagation & Self-play draw rate dropped from 50\% to 6.69\% (Section~\ref{finding:8}), but the 75k run drew 77\%. NoisyNet entropy rose throughout training (Section~\ref{finding:12}). The loss--win correlation ($r = -0.787$) was not detrended and does not show coupling (Section~\ref{finding:9}). \\
        \hline
        AAREN & Sequential information bottleneck & AAREN accuracy 24.23\% versus LSTM's 17.74\% (Section~\ref{finding:11}), not yet compared with a majority-class baseline of about 20--24\%. \\
        \hline
        Imperfect Information & Risk-averse draw-seeking policy & The self-play draw rate fell from 38.5\% (Section~\ref{finding:6}) to 6.69\% (Section~\ref{finding:8}). Against heuristic opponents the 75k run drew 77.02\% and won 11.66\%, 99.96\% of its wins by flag capture (Section~\ref{finding:14}). \\
        \hline
    \end{tabular}
\end{table}

The combined findings are consistent with the MARQ architecture, but they do not yet establish what each component contributes. Each component targets a distinct bottleneck. Draw dominance (Section~\ref{finding:6}) is plausibly a compound effect of scalar value collapse, exploration death and sparse reward propagation. Under the combined changes the self-play draw rate fell from 50.04\% to 6.69\% (Section~\ref{finding:8}). No ablation separates the effects of those changes, however, and the run against heuristic opponents still drew 77\% of its games.


\section{Supporting Evidence from Related Work}

These limitations are consistent with findings elsewhere in deep RL for imperfect information games.

P\'{e}rolat et al. (2022) solved full-size Stratego with DeepNash \cite{perolat2022stratego} using game-theoretic regularization (R-NaD), not Q-value maximization. DeepNash shows Stratego is solvable by deep RL, but only when the architecture explicitly handles non-stationarity and strategic diversity. The Ataraxos system \cite{Ataraxos2025} extended this to no-press Diplomacy with KL-divergence stabilization, a mechanism similar to MARQ's dynamic damping but absent from the normal DQN baseline.

Hessel et al. (2021) showed in the Rainbow DQN study that Double DQN, PER, Dueling Networks, Noisy Nets, C51, and multi-step returns each contribute additive gains \cite{RainbowJMLR2021}. Removing any single component degraded performance. No one fix is enough.


\section{Implications for MARQ}

The analysis attributes the normal DQN's performance ceiling to identifiable architectural properties. Table~\ref{tab:limitation_mapping} maps each finding to its theoretical cause, the MARQ component meant to address it, and how far the data so far supports that.

\begin{table}[htbp]
    \centering
    \caption{Observed limitations mapped to theoretical causes, MARQ components, and validation status.}
    \label{tab:limitation_mapping}
    \begin{tabular}{|p{3cm}|p{3.5cm}|p{3.5cm}|p{2.5cm}|}
        \hline
        \textbf{Finding} & \textbf{Theoretical Cause} & \textbf{MARQ Component} & \textbf{Status} \\
        \hline
        Q-Value Overestimation & Maximization Bias & Double DQN / Dueling & Not tested \\
        \hline
        Loss--Perf. Decoupling & Scalar Value Collapse & C51 Distributional RL & Not supported (Sec.~\ref{finding:9}) \\
        \hline
        LSTM Accuracy Ceiling & Sequential Bottleneck & AAREN Attention & Inconclusive (Sec.~\ref{finding:11}) \\
        \hline
        Exploration Collapse & Fixed $\epsilon$-greedy & Noisy Networks & Supported (Sec.~\ref{finding:12}) \\
        \hline
        Reward Propagation & 1-step Bootstrapping & $n$-step Returns & Not tested \\
        \hline
        Draw Dominance & Compound Effect & Full MARQ Pipeline & Self-play only (Sec.~\ref{finding:8}) \\
        \hline
    \end{tabular}
\end{table}

``Not tested'' means that no run isolated the component. Overestimation would need the values of the moves played compared with the returns that followed. The trainer's logged Q was a maximum over mostly untrained actions, and the win rate of Section~\ref{finding:7} does not measure bias. The steps-per-win figure of Section~\ref{finding:10} reflects the draw rate, not how far rewards propagate.

The normal DQN baseline in this chapter served as a reference point. In self-play the Rainbow agent with AAREN produced far fewer draws (6.69\% against 50.04\%) and shorter games, and its exploration noise stayed active. Against heuristic opponents it won 11.66\% of 75,047 games. Two questions remain open. What each component contributes needs ablations evaluated against fixed opponents. Whether AAREN supplies information the observation does not needs runs with and without it, compared on the same evaluations. The formal runs of Appendix~\ref{app:lessons} begin that work, so far with AAREN switched off.

%%%%%%%%%%%%%%%%%%%%%%%%%%%%%%%%%%%%%%%%%%%%%%%%%%%%%%%%%%%%%%%%%%%%%%%%%%%
\chapter{Conclusion and Recommendations}\label{chap:recommendations}
%%%%%%%%%%%%%%%%%%%%%%%%%%%%%%%%%%%%%%%%%%%%%%%%%%%%%%%%%%%%%%%%%%%%%%%%%%%

\section{Conclusion}

The experiments presented in this thesis set out to determine whether the performance ceiling of a normal DQN agent in Stratego arises from insufficient training or from structural properties of the architecture itself. The data do not settle that question in favour of the architecture. Across 37,213 episodes of curriculum training, the normal DQN held a stable 22\% win rate, its TD loss fell with no corresponding gain in play, its logged Q-values rose, its piece prediction accuracy stayed near 20\%, and 38.5\% of its games were drawn through passive piece shuffling. Chapter~\ref{chap:results} explained these observations by scalar value estimation, a decayed $\epsilon$-greedy schedule and the LSTM history encoder. Those explanations remain hypotheses. The Rainbow agent replaced each of the three, and it showed the same decoupling of loss from play. Its win rate against heuristic opponents (11.66\%) was no higher than the baseline's against the greedy heuristic (15.2\%), in runs that are not directly comparable.

The Rainbow DQN with AAREN changed four things at once. C51 distributional value estimation replaced scalar Q-learning. Noisy Networks replaced $\epsilon$-greedy exploration, and their noise stayed active throughout training. AAREN replaced the LSTM history encoder; its piece-identity accuracy was 24.23\%, against 17.74\% for the LSTM. Multi-step returns replaced single-step targets. In self-play the combined changes cut the draw rate from 50.04\% to 6.69\% and shortened the games. The self-play win rate rose from 24.97\% to 49.88\%, but in self-play both sides are the same agent, so that rise reflects the decisive share of games and not greater strength (Section~\ref{finding:7}).

The shift from drawn to decisive games in self-play is the clearest behavioural change in the data. It did not carry over to other opponents. Against fixed heuristic opponents the extended 75,047-episode run won 11.66\% of its games and drew 77.02\%, while its loss fell from 3.89 to 2.84. Of its wins, 99.96\% came by flag capture, and its average episode lasted 378 steps. Whether the agent learned to identify and pursue the flag, or only ended the games it won by that route, these figures cannot tell.

The runs made after those experiments (Appendix~\ref{app:lessons}), with the learner corrected and AAREN switched off, locate the agent's strength in its starting point and not in its reinforcement learning. A network given only the cold start, 100 games of the heuristic and 5,000 pre-training updates, scored 27 wins of 40 against the random opponent and 7 against the heuristic. The run that went on to train for 2,000 episodes from such a start scored 17 to 30 and 8 to 12 (Section~\ref{sec:cold-start-level}). A stronger scripted demonstrator, with a behaviour cloning loss over the legal moves, raised the starting point to 29 wins of 40 against the heuristic before a single training game (Section~\ref{sec:behaviour_cloning}). Five thousand episodes of reinforcement learning from that start, 393,845 updates, did not add to it: the wins against the heuristic stayed between 18 and 26 of 40. Against the demonstrator itself the agent went from 3 or 4 losses in 40 games to between 12 and 27 while the cloning term was still in the loss. When the demonstrations left the batch, the gap between the best and the second-best legal move fell from 0.126 to 0.006 and the agent's agreement with the demonstrator from 45.8\% to 14.6\% (Section~\ref{sec:formal14}).

The thesis therefore does not show that MARQ plays Stratego at a competitive level, nor that it plays better than the normal DQN. It connects each architectural component to a failure mode of the baseline without measuring any component's contribution, for three reasons. The configurations differ in more than one component, each was trained once, and no ablation was run. Several of the measures first used (the loss--win correlation, the logged Q-values, the self-play win rate) turned out not to measure what they were taken to measure (Chapter~\ref{chap:results}). What the work does establish is where the learner's signal ran out, which the remaining sections set out.

\section{Answers to the Research Questions}

\paragraph{Question 1: convergence speed and win rate against heuristic baselines.} The loss converged; the play did not. Against heuristic opponents the Rainbow agent with AAREN won 11.66\% of 75,047 games and drew 77.02\%, with no difference between the smart heuristic (11.74\%) and the frozen one (11.59\%). In the later runs the best level against the heuristic, 29 wins, 9 losses and 2 draws in 40 games, was reached by imitation before any reinforcement learning, and training did not raise it. A convergence speed for the quality of play cannot be reported, because no run improved its play measurably through reinforcement learning.

\paragraph{Question 2: statistically significant superiority over standard DQN and rule-based agents.} This was not demonstrated. The Rainbow agent and the normal DQN never played each other or the same fixed opponents, and each configuration was trained once, so there is no sample on which to test a difference. Against rule-based agents the evidence points the other way. The 75k run lost or drew 88\% of its games against the heuristics. The outlook expert, a scripted player that looks one move ahead, beat the heuristic 79 to 11 in 100 games (Section~\ref{sec:expert}). The best network produced in this work was a clone of that player: it drew 39 of 40 games with it straight after cloning, and at episode 5,000 of the run that followed it won none of 100 and lost 52 (Section~\ref{sec:formal14}).

\paragraph{Question 3: the impact of AAREN on inferring hidden pieces.} This remains open. AAREN's accuracy of 24.23\% exceeded the LSTM's 17.74\%, but always naming the most common piece type would score about 20 to 24\%, the two figures come from runs that differ in more than the encoder, and in the thesis code AAREN read a padding step and was in effect not trained by the value loss (Section~\ref{finding:11}). The later runs switched AAREN off. In them the information about hidden pieces came from planes computed by rule from public information, the counts of the ranks not yet revealed, and was not learned (Appendix~\ref{app:lessons}, Section~\ref{sec:outlook}). No pair of runs with and without AAREN, evaluated on the same games, exists yet.

\section{What the Data Show About Each Component}

\subsection{Value Estimation}

The rising Q-values and the falling loss of the baseline motivated two Rainbow components. The Dueling Network decomposes action values into state-value and advantage streams, which suits a game in which many quiet moves share nearly the same value. The C51 head learns the distribution of the return instead of its mean \cite{C51}.

The agent still chooses the move with the highest mean of that distribution. The distributional head therefore changes what is learned, not the rule by which a move is chosen: a risky move and a safe move with the same expected return remain tied, and C51 cannot by itself make the agent prefer decisive play. The later runs point to the opposite difficulty. The cross-entropy fits the whole return distribution, and differences of a few hundredths between the means of two moves, which is what separates a won exchange from a lost one, were learned slowly enough that a squared error on the means was added to the loss (Appendix~\ref{app:lessons}, Section~\ref{sec:formal7}). In formal\_14 the differences between moves were held up by the cloning term and fell twentyfold when it was removed (Section~\ref{sec:formal14}).

The other evidence is consistent with the design without isolating it. The correlation between loss and win rate shifted from 0.051 (normal DQN) to $-$0.787 (Rainbow DQN with AAREN). Both curves drift with training time, though, and in the later formal\_4 run the correlation shrank from $-0.49$ to $-0.14$ once that drift was removed, while the loss fell without any gain in evaluation (Section~\ref{finding:9}). The self-play draw rate dropped from 50.04\% to 6.69\%, but the 75k run against heuristics drew 77.02\%. Whether distributional value estimation changes the draw-seeking behaviour needs an ablation that changes only the value head.

\subsection{LSTM to AAREN Transition}

The LSTM prediction accuracy ceiling of approximately 18 to 22\% motivated the AAREN module. A separate history encoder has concrete advantages over a full DRQN architecture in replay compatibility, per-position granularity, and gradient isolation. A recurrent encoder must carry early observations through every later update, and AAREN replaces that recurrent computation with attention-based aggregation over the history, without modifying the replay buffer or the DQN backbone.

AAREN achieved 24.23\% piece identity inference accuracy compared to 17.74\% for the LSTM, a 6.5 percentage point difference. This does not validate the transition. Always naming the most common piece type would score about 20 to 24\%, the two runs differed in more than the encoder, and the thesis code took AAREN's embedding at a padding step (Section~\ref{finding:11}). The separated encoder design itself is a useful structure. Independent hidden states for each of the 100 board positions produce a spatially structured $(64, 10, 10)$ embedding tensor that preserves piece-level temporal information, where a monolithic DRQN hidden state would mix the action histories of all 40 opponent pieces into a single vector. The per-position design also remains compatible with standard experience replay by storing compact embedding snapshots alongside the board tensors, avoiding the sequence burn-in that full DRQN architectures require.

\subsection{Exploration}

The normal DQN had no noisy layers. Its exploration came from an $\epsilon$-greedy schedule that decayed to 0.01 and stayed there, and the zero noise entropy logged for it is a property of that design, not a measurement. Noisy Networks \cite{NoisyNetworks} parameterize exploration as learnable per-weight perturbations. In the Rainbow run their entropy rose from 0.034 to 0.049 over training, so the learned noise did not collapse. How much this exploration contributed to the change in outcomes is not known without an ablation.

The later runs add a caution. Exploration noise competes with the differences between the values of the moves. In formal\_14 the noise changed the greedy move in 14\% of positions at episode 4,500 and in 31\% at episode 5,000, although its scale had fallen, because the differences it competes with had shrunk (Section~\ref{sec:formal14}).

\subsection{Long-Range Credit Assignment}

Single-step TD learning attenuates terminal rewards by $\gamma^{200} \approx 0.134$ for mid-game decisions at $\gamma = 0.99$. Five-step returns ($n = 5$), as specified in MARQ, shorten the bootstrap chain. The steps-per-win figure (1,060 versus 2,697 for the normal DQN) does not test this: in self-play it follows the draw rate (Section~\ref{finding:10}).

The later runs measured how short the chain of updates is at this scale. formal\_12 made 33,585 updates in 1,500 episodes, about 22 per episode, and copied the online network into the target every 1,000 updates, so the outcome of a game had been bootstrapped through about 34 target networks (Section~\ref{sec:cold-start-level}). Whether multi-step returns accelerate the propagation of the outcome remains to be measured; at this number of updates the outcome reaches few of the decisions that led to it, whatever the value of $n$.

\section{Computer Science Contributions}

The following techniques and lessons are not specific to Stratego. Each is stated with what the experiments did and did not show about it.

\subsection{Five-Phase Curriculum Design}

The curriculum transitions agents from full to partial observability in stages. Observability drops progressively while opponent difficulty rises and reward shaping decays. The intent is to mitigate the cold-start problem, so that agents do not need to learn basic mechanics and hidden-information reasoning at the same time. The 75k run never passed the gate of its second phase, and the later runs found a cost to the structure itself (Appendix~\ref{app:lessons}). Changing the observation and the shaping weight between phases made the network relearn the game at each switch. The current opponent schedule therefore keeps fog of war, the turn limit and the shaping fixed, and varies only the opponents.

\subsection{Embedding Snapshots for Replay Compatibility}

Hybrid architectures face a tension between replay accuracy and speed. Reconstructing full history sequences at replay time gives correct gradients but costs 4$\times$ throughput. Storing detached snapshot embeddings keeps throughput, but the DQN loss can then no longer train the encoder. The thesis code tried to restore that path with a ``gradient bridge'', adding $\epsilon \cdot W_{proj}(\mathbf{0})$ to the stored embeddings. That term equals the projection's bias, so it reached only that one parameter tensor and nothing that depends on the history (Chapter~\ref{chap:methodology}). The design that works is a separately supervised encoder: the snapshots are plain inputs to the Q-network, and the encoder is trained by its own loss on revealed pieces. Any encoder that is meant to learn from the value loss should be checked by listing which of its parameters actually receive a gradient.

\subsection{Expectamax Search Integration}

At test time, the system can blend the DQN's Q-values with combat utilities derived from rank probabilities through Expectamax search. The intent is to stop attacks on hidden pieces that the value estimate favours and the rank probabilities do not. Its effect on playing strength was not measured in this thesis. The later runs suggest why it could matter: a scripted player that scores each attack by its expected value over the target's rank distribution, with no learned component, beat the heuristic 79 to 11 in 100 games (Section~\ref{sec:expert}).

\subsection{Potential-Based Reward Shaping with Phase Annealing}

Potential-based shaping leaves the optimal policy unchanged \cite{NgPBRS} while providing dense signals during early training. A multiplicative schedule tapers the shaping coefficient to zero across curriculum phases, so the final training phase runs on terminal rewards alone.

The invariance holds at convergence, not along the way. In the later formal runs a potential that valued a hidden enemy piece at the mean rank of an army raised the potential whenever a lost attack revealed a Bomb, so throwing cheap pieces at unmoved ones was rewarded, and the agent lost by running out of movable pieces (Appendix~\ref{app:lessons}, Section~\ref{sec:formal11}). A potential should not reward revealing what it estimates; here the enemy's material could be computed exactly from public information.

\subsection{Demonstrations Set the Level}

A value-based learner started from demonstrations should have its starting point measured alone. In this work the cold start accounted for the whole level of the trained agent (Section~\ref{sec:cold-start-level}). The loss used for the demonstrations decided what was transferred: with the large-margin loss of DQfD \cite{DQfD} the network took up the demonstrator's appetite for battle and lost about half of its attacks on revealed pieces, where the demonstrator lost none; with a cross-entropy over the legal moves it lost 17\% and drew with its teacher (Sections~\ref{sec:expert} and~\ref{sec:behaviour_cloning}). A demonstration schedule that decays to zero removed the policy with it (Section~\ref{sec:formal14}).

\subsection{Checks That Reversed a Claimed Result}

Five checks, each cheap, reversed a result this study had first read as validation (Appendix~\ref{app:lessons}). A win rate from symmetric self-play measured the draw rate, not strength. A correlation between two curves that both drift with time measured the drift. An accuracy needed the majority-class baseline, not the uniform one. An average Q-value taken as a maximum over all actions measured the untrained illegal ones. A claimed gradient path reached one bias. A sixth belongs with them: a falling loss was repeatedly read as learning, and in formal\_14 the loss kept falling while the network stopped distinguishing moves.

\section{Limitations and Future Work}

The validation presented in this thesis was bounded by the training infrastructure, opponent diversity, and compute budget available during the research period. The following limitations qualify the conclusions drawn above and identify directions where additional work is needed.

The Rainbow DQN with AAREN produced far fewer draws than the normal DQN baseline in self-play. The two were never evaluated against each other or against the same fixed opponents, and against heuristic opponents the Rainbow agent won 11.66\% of its games. Each configuration was trained once, and the periodic evaluations used 40 games per opponent, for which the standard deviation of a win rate near one half is about 8 percentage points. Most differences between consecutive runs lie within that noise.

Combining C51, Noisy Nets, Dueling Networks, AAREN, and $n$-step returns in a multi-agent self-play loop introduces interaction effects that single-component theory does not predict. The extended 75k training run remained in the partial observability phase against heuristic opponents for its entire duration, so the later phases of the curriculum, including league training, were not exercised by it.

The full MARQ pipeline is heavier than the normal DQN. Each checkpoint occupies approximately 433\,MB compared to 39\,MB, an 11-fold increase. The Rainbow run used fewer environment steps (5.17 million against 10.31 million), but that reflects shorter games and fewer draws, not more learning from each sample (Section~\ref{finding:10}). The number of updates is the binding constraint: the later run with the most updates made 393,845 in 5,000 episodes on one consumer GPU.

The later runs show that, in this system, reinforcement learning has not yet improved an agent beyond its starting point. Three directions follow from where it failed. First, the policy should not depend on small differences between action values. In formal\_14 those differences were sustained by the cloning term and not by the value targets, so a learner that keeps a policy of its own, either a separate policy output trained beside the value or a cloning term that is kept instead of decayed, is the first thing to try. Second, the number of updates per hour should be raised, with a faster environment and a lighter value head, before a plateau is read as a limit of the method. Third, every run should report the strength of its starting point beside the trained agent's, from evaluations of at least 100 games against the same fixed opponents.

A first test of these directions was made after the thesis experiments and is reported in Appendix~\ref{app:lessons} (Sections~\ref{sec:vec-engine} to~\ref{sec:ppo}). A network with a policy output of its own, cloned from the scripted player and then trained by a policy gradient on self-play in a batched game engine, improved on its starting point within minutes. After 45 minutes it won 193 of 200 games against the heuristic and 177 of 200 against the scripted player it had been cloned from, in the thesis's own evaluation. That run changed the learner, the amount of experience and the reward of a draw together, so it shows that the plateau above was not a limit of reinforcement learning in this game, without yet showing which change removed it.

Ablation studies are needed. The data does not yet show that the combined architecture plays better than the normal DQN, nor what each component contributes. Future work should test each component in isolation, evaluated in greedy games against the same fixed opponents, with enough games and enough repeated runs for a difference to exceed its noise. AAREN in particular needs paired runs with and without it, and a majority-class baseline for its accuracy. Beyond fixed opponents, the agent's exploitability should be measured by training a fresh exploiter against the frozen policy. A deterministic greedy policy can be exploited in a game of hidden information, and methods built for such games (NFSP, PSRO, R-NaD) add an averaged or regularized policy for that reason.

\include{chapters/AppendixA}
\include{chapters/AppendixB}
\chapter{Implementation Log and Lessons Learned}
\renewcommand{\thefigure}{C.\arabic{figure}}
\renewcommand{\thetable}{C.\arabic{table}}
\label{app:lessons}

This appendix records the changes made to the training system after the thesis experiments,
each with the evidence that motivated it and what it taught. Entries are kept in the order the
work was done, and every entry follows the same pattern: the change, why it was made, and the
lesson. New entries are added at the end.

% ------------------------------------------------------------------------------------------
\section{Narrowing the C51 support to $\pm 2$}

\paragraph{Change.} The value support of the distributional head \cite{C51} was narrowed from
$[-10, 10]$ to $[-2, 2]$ (51 atoms, so 0.08 per atom instead of 0.4).

\paragraph{Why.} A return is the outcome ($+1$, $-1$ or the draw reward $-0.5$) plus the
potential-based shaping $\gamma\Phi(s') - \Phi(s)$ \cite{NgPBRS}, which telescopes over a game to
$-\Phi(s_0)$. With $\Phi \in [-0.4, 1]$ every return therefore lies in $[-2, 1.4]$; the returns
observed in the formal\_2 run spanned $-1.54$ to $0.71$. The $\pm 10$ support was inherited from
the thesis configuration, whose shaping was not potential-based. On it, the shaping of a capture
(about 0.02 to 0.05) was a tenth of an atom: the head could hardly represent the difference
between a good and a bad exchange.

\paragraph{Lesson.} The support of a distributional critic is part of the reward design, not a
default. It should be derived from the bounds of the return (here the telescoping of the
shaping makes that exact) and checked against the returns a run actually produces.

% ------------------------------------------------------------------------------------------
\section{Measuring the average Q-value on the actions taken}

\paragraph{Change.} The logged ``average Q'' is now the expected value of the action that was
taken in each sampled replay transition.

\paragraph{Why.} The thesis metric took the maximum over all 3600 actions. Nearly all of them are
illegal in any given state and are never trained, so the maximum was an artefact: it read about
$+2$ while the returns were about $-1$.

\paragraph{Lesson.} With a large, mostly illegal action space, any statistic taken over all
actions measures the untrained part of the network. Diagnostics must be restricted to actions
that can occur.

% ------------------------------------------------------------------------------------------
\section{Diagnosing a flat value function}

\paragraph{Observation.} After thousands of episodes of the formal\_2 run the agent still drew
most games, and its Q-values barely distinguished good moves from bad ones.

\paragraph{Diagnosis.} Two probes separated the possible causes. A \emph{clean-signal probe}
trained the same network on synthetic transitions in which only the combat outcome determined the
reward; it learned the difference quickly, so the learner itself (loss, target, optimiser) was
sound. A \emph{same-state probe} then compared, in real positions, the values of a winning and a
losing attack available in the same state. In the formal\_3 run at episode 1250 the winning attack
was preferred in 68\% of such positions, but by a margin of 0.003, against a true difference of
about 0.05. In real training the combat signal is small next to the noise: a winning attack is
worth $+0.019$ of shaping and a losing one $-0.033$, while an ordinary move's shaping varies with
a standard deviation of 0.02, and the outcome arrives hundreds of moves later.

\paragraph{Lesson.} When a reinforcement learner does not learn, first test the learner on a
problem whose answer is known; if it passes, the problem is in the signal (reward scale, support,
credit assignment), not in the optimisation. Probes that compare actions \emph{within the same
state} are far more informative than averages over states.

\paragraph{Operational note.} A run that ends with exit code 1 and no Python traceback was
killed from outside the process (here by memory pressure), not by a bug in the training code.

% ------------------------------------------------------------------------------------------
\section{Making the training loop faster}

\paragraph{Changes.}
\begin{itemize}
  \item \textbf{Batched evaluation.} The evaluation games now run side by side, with one network
        call per side per round instead of one call per position. An evaluation of
        $40 + 40$ games went from 363\,s to 143\,s.
  \item \textbf{Cheaper replay.} The C51 loss needs the distribution of the action taken only,
        and the target needs that of the Double-DQN action only. The network now returns raw
        logits; the two needed actions are gathered first and normalised afterwards, instead of
        normalising all $3600 \times 51$ outputs. One update went from 128.7\,ms to 97.5\,ms.
  \item \textbf{More lanes at a constant replay ratio.} The loop used to perform at most one
        gradient update per iteration, so adding parallel games silently lowered the number of
        updates per move. It now performs one update per \texttt{REPLAY\_UPDATE\_INTERVAL}
        learner moves however many lanes moved, which allowed 16 lanes instead of 8 (acting cost
        per move fell from 2.0\,ms to 1.33\,ms).
\end{itemize}

\paragraph{Lesson.} Profile before optimising: replay was half of the step time and evaluation
was serial, while the environment was below 10\%. And a ``speed'' setting such as the number of
lanes can change the learning algorithm (here the replay ratio) unless the loop is written so
that it cannot.

% ------------------------------------------------------------------------------------------
\section{Removing the curriculum's environment changes}

\paragraph{Change.} The five-phase curriculum no longer changes the game. Every stage is played
under fog of war, with the same turn limit (1500), the same no-battle draw rule (200 plies,
also used in evaluation and the GUI) and the same, constant shaping weight. The stages remain
only as an opponent schedule: random and scripted opponents first, then the stronger heuristic,
then self-play and the league, advancing on evaluation gates as before.

\paragraph{Why.} The thesis curriculum \cite{CurriculumLearning} changed the environment under
the network at every transition. Phase 1 had full observability (the agent saw every enemy
rank), and the shaping weight fell from 1.0 to 0.5, 0.2 and finally 0. Each of these changes the
optimal values of the states the network has already learned: a value learned with the enemy's
ranks visible is wrong under fog of war, and cutting the shaping weight changes every return.
The network must then relearn what it knew, with the replay memory still full of transitions
from the old rules, which is the setting in which catastrophic forgetting
\cite{CatastrophicForgetting} occurs. Annealing the shaping also buys nothing: potential-based
shaping leaves the optimal policy unchanged \cite{NgPBRS}, so there is no bias to remove. A
smaller inconsistency was the draw rule: training drew after 100 plies without a battle and
evaluation after 200, so the agent was graded under a rule it had not trained under.

\paragraph{Lesson.} A curriculum should make the \emph{task} easier (weaker opponents), not
change the \emph{rules} the values refer to. Anything the network's targets depend on
(observability, reward, termination) must stay fixed across stages, and evaluation must use the
same rules as training.

% ------------------------------------------------------------------------------------------
\section{Cold start from demonstrations}

\paragraph{Change.} Before its first game the agent now learns from games in which the scripted
heuristic plays its side, following Deep Q-learning from Demonstrations (DQfD) \cite{DQfD}; the
idea of a short supervised ``cold start'' before reinforcement learning is also central to recent
reasoning-model training \cite{DeepSeekR1}.
\begin{enumerate}
  \item The heuristic plays 100 games against the random and heuristic opponents under the
        training rules. Its transitions are recorded exactly as the agent's own are (move plus
        reply, PBRS shaping plus outcome, 5-step returns, mirror augmentation) into a separate
        demonstration buffer that the agent's own experience never overwrites. Each one also
        keeps the legal moves of its state.
  \item The network is pre-trained on the demonstrations alone (5000 updates).
  \item During training 10\% of each batch comes from the demonstrations, falling linearly to
        0\% by episode 5000.
  \item Demonstrated transitions also get the large-margin loss of DQfD,
        \[
          J_E(Q) = \max_{a \in \mathcal{A}(s)} \bigl[ Q(s,a) + \ell(a_E, a) \bigr] - Q(s, a_E),
          \qquad \ell(a_E, a) = \begin{cases} 0 & a = a_E \\ 0.1 & a \neq a_E \end{cases}
        \]
        with the maximum over the legal moves $\mathcal{A}(s)$ only, added to the C51 loss with
        weight 1.
\end{enumerate}
The earlier ``imitation'' option, which let the heuristic choose 10\% of the agent's moves in
its first 2000 games, is switched off.

\paragraph{Why.} A network trained from zero starts from random play, which against these
opponents mostly ends in draws: at episode 1000 the formal\_3 run drew 32 of 40 evaluation games
against the random opponent. Its first thousands of games therefore carry almost no outcome
signal. The imitation option did not fix this: the heuristic's moves entered the replay memory
like the agent's own, so the C51 loss learned what they were worth but nothing told the network
to prefer them over the moves that were not played. The margin loss does exactly that. In a unit
test with equal returns for every move, the C51 loss alone left the greedy move at chance, while
with the margin loss the demonstrated move became the greedy one in at least 90\% of states.
The margin is 0.1 (1.25 atoms of the $\pm 2$ support), comparable to the value of a capture; DQfD
used 0.8 with rewards clipped to $\pm 1$. The demonstration share decays because the heuristic is
a weak demonstrator, not an expert, and should not cap the agent.

\paragraph{Result.} The 100 demonstration games took about three minutes; the heuristic won
about half of them (52 wins, 26 losses, 22 draws in one collection) and they produced 38\,742
transitions, mirror copies included. Pre-training costs about 107\,ms per update, so the 5000
updates take about nine minutes. Whether the cold start speeds up learning has not yet been
measured in a full run.

% ------------------------------------------------------------------------------------------
\section{The formal\_3 run and the start of formal\_4}

\paragraph{formal\_3.} The formal\_3 run trained from zero with the $\pm 2$ support but still
with the curriculum's environment changes and without the cold start. It was stopped at episode
3500, never having left Phase~1. Its evaluations (40 greedy games per opponent every 500
episodes) did not improve; against the random opponent they got worse:
\begin{center}
\begin{tabular}{rcc}
\toprule
Episode & vs random (W/L/D) & vs heuristic (W/L/D) \\
\midrule
500  & 8/14/18 & 1/37/2  \\
1000 & 7/1/32  & 1/28/11 \\
1500 & 6/2/32  & 1/32/7  \\
2000 & 0/4/36  & 0/29/11 \\
2500 & 5/8/27  & 2/30/8  \\
3000 & 3/5/32  & 0/35/5  \\
3500 & 1/3/36  & 0/35/5  \\
\bottomrule
\end{tabular}
\end{center}
The agent learned to avoid losing to the random opponent mainly by drawing, and it never won
more than 2 of 40 games against the heuristic. Its games carried almost no outcome signal, which
is the problem the cold start addresses.

\paragraph{formal\_4.} formal\_4 starts from zero with every change above: the opponent schedule
without environment changes, the cold start, the asynchronous learner and 16 lanes. Its cold
start played the 100 demonstration games in 84\,s (the heuristic won 56, lost 24 and drew 20)
and stored 35\,544 transitions. The 5000 pre-training updates took 271\,s (54\,ms per update with
the GPU not shared with another run), and the large-margin loss fell to 0.068, below the margin
of 0.1.

\paragraph{formal\_4, first 1000 episodes.} The evaluations (40 greedy games per opponent) were,
next to formal\_3 at the same episodes:
\begin{center}
\begin{tabular}{rcccc}
\toprule
 & \multicolumn{2}{c}{formal\_4} & \multicolumn{2}{c}{formal\_3} \\
Episode & vs random & vs heuristic & vs random & vs heuristic \\
\midrule
500  & 16/16/8 & 10/17/13 & 8/14/18 & 1/37/2  \\
1000 & 14/22/4 & 8/21/11  & 7/1/32  & 1/28/11 \\
\bottomrule
\end{tabular}
\end{center}
(wins/losses/draws). Against the heuristic formal\_4 won 8 to 10 games where formal\_3 won one,
and it no longer drew most games against the random opponent. But it lost more than half of them
instead, almost always by losing all its movable pieces, and it did not improve from episode 500
to 1000, nor by episode 2500 (vs random 12/23/5, 12/25/3 and 15/22/3 at episodes 1500, 2000 and
2500; vs heuristic 4/14/22, 5/19/16 and 6/16/18). formal\_4 also ran about 1.4 times as fast (1000 episodes in 45 minutes, against about
64 for formal\_3).

\paragraph{Diagnosis: attacks on revealed pieces.} A probe played 30 greedy games against the
random opponent under the training rules and classified every attack by whether the attacked
piece's rank was already known to the attacker:
\begin{center}
\begin{tabular}{lccc}
\toprule
Player 1 & attacks per game & hidden target: won / lost & revealed target: won / lost \\
\midrule
heuristic (the demonstrator) & 39.1 & 39\% / 47\% & 77\% / 0\%  \\
formal\_4, episode 250       & 41.7 & 45\% / 45\% & 26\% / 62\% \\
formal\_4, episode 500       & 41.0 & 42\% / 43\% & 27\% / 63\% \\
formal\_4, episode 1000      & 39.8 & 43\% / 42\% & 28\% / 63\% \\
formal\_3, episode 1000      & 20.2 & 36\% / 48\% & 22\% / 72\% \\
\bottomrule
\end{tabular}
\end{center}
Against hidden pieces the agent does as well as the heuristic. Against pieces whose rank it can
see (the ranks of revealed enemy pieces are in the observation, and a test checks that they reach
the network) it loses most of its attacks, while the heuristic never attacks a known piece it
cannot beat. The cold start transferred how \emph{often} the heuristic attacks (about 40 attacks
per game, twice formal\_3's), but not \emph{which} known pieces it attacks, and this was so from
episode 250 on, so it was never learned rather than learned and forgotten.

\paragraph{Test: pre-training on the demonstrations alone.} To see whether more imitation would
teach the rule, a separate copy of the network was pre-trained on a fresh set of 100
demonstration games (35\,198 transitions) with the cold start's losses, for 20\,000 updates
instead of 5000, and probed in the same way. The agreement is the share of demonstration states
in which the network's greedy legal move is the heuristic's:
\begin{center}
\begin{tabular}{rcccc}
\toprule
Updates & agreement & margin loss & revealed target: won / lost & hidden target: won / lost \\
\midrule
0      & 8.4\%  & --    & --          & --          \\
5000   & 63.1\% & 0.072 & 28\% / 60\% & 41\% / 45\% \\
10\,000 & 75.7\% & 0.051 & 28\% / 60\% & 44\% / 42\% \\
20\,000 & 88.2\% & 0.027 & 31\% / 62\% & 42\% / 44\% \\
\bottomrule
\end{tabular}
\end{center}
The network came to reproduce the heuristic's move in 88\% of the demonstration states, yet in its
own games it lost as many attacks on revealed pieces as before. It also did worse against pieces
it could see than against hidden ones, as a player would that picks its targets without regard to
their visible rank: the enemy pieces that are revealed are mostly ones that survived an earlier
battle, so they are stronger than average. Fitting the demonstrations, even closely, did not
produce the rule that makes the heuristic's attacks good. (The agreement is measured on the
demonstrations the network was trained on, so part of it may be memorisation.)

\paragraph{Test: does the value of an attack respond to the target's rank?} In 300 real
positions (random play, fog of war) where a piece of rank Miner to General could attack an
adjacent enemy piece, the observation was edited at the target square only, showing the target
as a revealed Scout (the attack wins), a revealed Marshal (the attack loses) or as hidden, and the
network's value of that attack was compared across the three. The same edit was also applied
while reading the value of a quiet move far from the target, as a control.
\begin{center}
\begin{tabular}{lcccc}
\toprule
Network & $Q$ vs Scout & $Q$ vs Marshal & $Q$ vs hidden & Scout $-$ Marshal (positive in) \\
\midrule
untrained               & $-0.023$ & $-0.025$ & $-0.022$ & $+0.002$ (51\%) \\
formal\_3, episode 3500  & $-0.642$ & $-0.644$ & $-0.628$ & $+0.001$ (46\%) \\
formal\_4, episode 2500  & $+0.055$ & $+0.048$ & $+0.187$ & $+0.007$ (50\%) \\
pre-trained, 20\,000 updates & $-0.140$ & $-0.172$ & $+0.074$ & $+0.032$ (59\%) \\
\bottomrule
\end{tabular}
\end{center}
The control moved by less than 0.01 in every case. The networks trained by reinforcement value
an attack on a revealed Scout and on a revealed Marshal the same, no more apart than the untrained
network does, although one wins a piece and the other loses one. They do respond to \emph{whether}
the target is revealed (formal\_4 values attacking a hidden piece at $+0.19$ and a revealed one at
about $+0.05$), but not to \emph{which} rank is revealed. The pre-trained network shows a small
preference ($+0.03$), with a spread across positions (standard deviation 0.12) four times as
large. The environment keeps a revealed piece's rank visible when it moves, and the ranks are in
the observation, so the information is there: it is the training signal that has not made the
networks use it. This is the combat signal of the earlier diagnosis: the shaping difference
between winning and losing an exchange (a few hundredths) is no larger than the variation of
an ordinary move's shaping, and the outcome comes hundreds of moves later.

\paragraph{Measurement: the combat signal of the potential.} The shaping of Player-1
transitions (its move plus the reply, as the trainer builds them; terminal transitions left out)
was measured in 30 games against the random opponent, with the heuristic and with a random player
as Player 1, for the current potential and for its material term alone,
$\Phi_{\mathrm{mat}} = (\text{own} - \text{enemy material}) / 148$:
\begin{center}
\begin{tabular}{lcc}
\toprule
 & current $\Phi$ & material only \\
\midrule
quiet move, mean / sd (random Player 1)        & $+0.0008$ / 0.015 & $0.0000$ / 0.0001 \\
quiet move, mean (heuristic Player 1)          & $+0.0065$         & $0.0000$ \\
attack on a revealed piece, won $-$ lost (random) & 0.043          & 0.061 \\
\quad in standard deviations of a quiet move   & 2.9               & about 430 \\
\quad in atoms of the $\pm 2$ support          & 0.53              & 0.76 \\
attack on a hidden piece, won $-$ lost (heuristic) & 0.093         & 0.058 \\
\bottomrule
\end{tabular}
\end{center}
Under the current potential an ordinary move is not neutral: the proximity and penetration terms
reward advancing, so for the heuristic a quiet move is worth $+0.0065$ on average, and beating a
revealed piece only $+0.011$. Every 5-step return also adds four more transitions of this
variation (standard deviation about 0.02 each) to the difference between a won and a lost
exchange of about 0.04. With the material term alone, a move without a battle is worth exactly
nothing (up to the discount), so the only shaping left in a return is that of the battles in it,
and a won exchange differs from a lost one by 0.06, about three-quarters of an atom. The
information and proximity terms also explain why the networks prefer attacking hidden pieces:
revealing a piece raises the information term whatever the result.

\paragraph{Lesson.} A shaping potential with several heuristic terms makes every move carry
reward, and that variation is noise to the decision the agent most needs to learn. Before tuning
a learner, measure the reward difference between the decisions that matter, relative to the
variation of the rest of the reward stream, and in the units the critic resolves.

\paragraph{Lesson.} Imitation transfers the frequent, easy-to-see features of a demonstrator's
play (here, attacking a lot) before the conditional rule behind them (attack a known piece only
if it is weaker), and a high agreement with the demonstrator does not show that the rule was
learned. A demonstration is worth checking for the specific decisions it was meant to teach, with
a probe that classifies those decisions in the agent's own games, not only by win rate or by
agreement on the demonstrations.

\paragraph{Lesson.} An evaluation that only counts wins hides how a policy fails: formal\_3's
win rate against the random opponent stayed low throughout, but its draws rising from 18 to 36
of 40 showed that it was converging to passive play rather than failing to learn at all.

\paragraph{Lesson.} Demonstrations stored as ordinary experience teach values, not preferences;
to make a value learner imitate, the loss must rank the demonstrated action above the
alternatives. And with a large, mostly illegal action space, that ranking has to be restricted to
the legal moves, or it spends its effort pushing down actions that can never be played.

% ------------------------------------------------------------------------------------------
\section{A material-only potential (formal\_5)}

\paragraph{Change.} The weights of the shaping potential's four terms (material, penetration,
information, proximity) became a setting, \texttt{POTENTIAL\_WEIGHTS}, and were changed from
$(0.4, 0.15, 0.1, 0.35)$ to $(1, 0, 0, 0)$:
\[
\Phi(s) = \frac{\text{own material} - \text{enemy material}}{148},
\]
with a hidden enemy piece still counted at the mean army rank of 3.7. Its range is $[-1, 1]$, and
the support test now reads the range from the weights. $\Phi$ is 0 at the start of a game (both
armies have 148 material, the hidden one counted at the mean), so the shaping of a whole game
telescopes to zero, and a return still lies within the outcome reward plus or minus 1, inside the
$\pm 2$ support. Everything else is formal\_4's configuration, including the cold start and the
asynchronous learner. formal\_4 was stopped at episode 3000; its last evaluation had 17 wins,
17 losses and 6 draws against the random opponent and 3 wins, 24 losses and 13 draws against the
heuristic, with no gain since episode 1500.

\paragraph{Why.} The measurement above: under the old potential an ordinary move carried a
shaping reward that varied as much as the difference between winning and losing an exchange on a
revealed piece, and the network did not learn to tell those exchanges apart. With material only,
a move without a battle carries no shaping, and a won exchange is worth about three-quarters of
an atom more than a lost one. The cost is that nothing rewards advancing or scouting any more
until a battle happens; whether the agent still finds the enemy, or stalls into no-battle draws,
is what formal\_5 has to show.

\paragraph{Result, first 1000 episodes.} The cold start went as in formal\_4 (the heuristic won 58,
lost 20 and drew 22 of the 100 demonstration games; the margin loss fell to 0.071). The
evaluations were 20/16/4 and 9/30/1 against the random opponent and 3/18/19 and 9/19/12 against
the heuristic at episodes 500 and 1000 (wins/losses/draws), against 16/16/8 and 14/22/4, 10/17/13
and 8/21/11 for formal\_4. The agent did not stall: no-battle draws were 74 and 94 of the
training games in the two blocks of 500, and games lasted 340 to 346 moves. It still lost
mostly by running out of movable pieces (197 and 198 of 500 games), as formal\_4 did. So far the
material-only potential has not changed how the agent loses. Two evaluations 500 episodes apart
also differed by 11 wins out of 40 against the same opponent, which shows how wide the noise of a
40-game evaluation is.

\paragraph{Result, to episode 2750.} formal\_5 was stopped at episode 2750 to start formal\_6
(Section~\ref{sec:formal6}) with the fixes found in the meantime. Its later evaluations were
22/15/3, 17/22/1 and 16/20/4 against the random opponent and 9/24/7, 8/20/12 and 7/14/19 against
the heuristic at episodes 1500, 2000 and 2500. formal\_4 had 12, 12 and 15 wins against the random
opponent and 4, 5 and 6 against the heuristic at the same episodes. The material-only potential
won a few more games against the heuristic, but by less than the noise of 40-game evaluations
(about 3 wins), and neither run passed a gate. The potential was not the bottleneck it was
suspected to be; the faults of the next sections were found while formal\_5 ran.

\paragraph{Result, the revealed-piece probes.} The two probes of formal\_4 were repeated on the
networks of formal\_5 (episode 2750) and formal\_4 (episode 3000). Both have 30 input planes, so
they were run under the formal\_6 code with the agent built for 30 planes and given the first 30
planes of the new encoding, whose layout is unchanged. Each played 30 greedy games as Player~1
against the random opponent:
\begin{center}
\begin{tabular}{lccc}
\toprule
 & formal\_5, ep.\ 2750 & formal\_4, ep.\ 3000 & heuristic \\
\midrule
games won / lost / drawn               & 17 / 13 / 0 & 14 / 16 / 0 & 23 / 4 / 3 \\
attacks on revealed pieces             & 378         & 339         & 270 \\
\quad won / tied / lost                & 28 / 9 / 63\% & 30 / 11 / 59\% & 77 / 23 / 0\% \\
attacks on hidden pieces, won / lost   & 43 / 43\%   & 44 / 42\%   & 39 / 47\% \\
\bottomrule
\end{tabular}
\end{center}
A network whose checkpoint failed to load (a fresh one, through the silent fallback of
\texttt{load\_model}) lost 66 to 67\% of its attacks on revealed pieces in the same games, so on
this decision both trained networks are close to an untrained one. In the rank-swap test (300
positions, the target shown as a revealed Scout or a revealed Marshal, everything else unchanged)
formal\_5's value of the attack rose by $+0.0016$ for the Scout, positive in 49\% of positions,
and formal\_4's by $+0.0011$, positive in 51\%, against a spread of the attack's value across
positions of 0.06 to 0.09. With a won exchange worth about three-quarters of an atom more than a
lost one in every return that contains it, 2750 episodes of material-only shaping did not make
the network read the revealed rank at all. A reward signal that clear and that frequent should
have been learnt, so the cause has to be sought in the learner or in what reaches it rather than
in the reward. None of the faults of the next sections has yet been shown to explain it.

\paragraph{Lesson.} When a change aimed at one failure leaves that failure exactly where it was,
measure the failure directly before reading the win rates: the probe settled in minutes what
three noisy evaluations could not. And a diagnostic that loads a checkpoint must fail loudly when
the load fails; here a silent fallback to a fresh network gave plausible-looking numbers.

% ------------------------------------------------------------------------------------------
\section{Asynchronous learner}

\paragraph{Change.} The gradient updates now run on a background thread with its own CUDA
stream while the main loop chooses moves and steps the parallel games. This is the actor-learner
split of Ape-X \cite{ApeX} and IMPALA \cite{IMPALA}, also used for large-scale asynchronous
reinforcement learning of language models \cite{AReaL}, here within a single process:
\begin{itemize}
  \item The replay ratio is unchanged: the loop still owes one update per
        \texttt{REPLAY\_UPDATE\_INTERVAL} learner moves and hands them to the learner. Once more
        than \texttt{ASYNC\_MAX\_LAG} $= 2$ updates are owed it waits (back-pressure), so the
        policy that acts is at most three updates behind the one being trained.
  \item Acting reads a copy of the network that is refreshed after every update. The learner
        must not share its network with the actor: the optimiser changes the weights in place,
        and NoisyNet acting switches the network's noise and train/eval mode.
  \item The replay memory, which both threads use, is behind a lock. Because the learner runs on
        its own CUDA stream, it waits for the actor's stream before sampling and finishes its
        reads before releasing the lock; otherwise a transition evicted by the actor could be
        freed by the allocator while the learner's kernels were still reading it.
  \item Anything that reads or replaces the learner's weights or optimiser state (checkpoints,
        copying the network into the second agent, the learning-rate schedule) pauses the
        learner first, and an error in the learner thread is re-raised in the training loop.
\end{itemize}

\paragraph{Result.} In two otherwise identical short runs (8 lanes, the cold start on, measured
from the first gradient update over about 10\,000 learner moves, on a GPU shared with another
training run) the synchronous trainer made 43.4 learner moves per second and the asynchronous one
59.3, i.e.\ 1.37 times as fast. Both performed exactly 1313 gradient updates, so the replay ratio
was the same.

\paragraph{Lesson.} Asynchrony is only safe if every piece of shared state is accounted for: the
weights (a separate acting copy), the replay memory (a lock), GPU memory on two streams (stream
ordering and synchronisation), and the replay ratio (back-pressure). The last one matters most
for reproducibility: without back-pressure an asynchronous learner changes the algorithm, not
just its speed.

% ------------------------------------------------------------------------------------------
\section{What the falling loss of formal\_4 did and did not show}

\paragraph{Observation.} Over the 3000 episodes of formal\_4 the C51 loss of Player~1 fell
steadily while its evaluations did not improve. Per block of 500 training episodes (about 11\,500
gradient updates each):
\begin{center}
\begin{tabular}{rcccc}
\toprule
Episodes & mean loss & mean $Q$ (action taken) & mean return & training games won (random / heuristic) \\
\midrule
0--500     & 3.63 & $-0.01$ & $-0.66$ & 34\% / 16\% \\
500--1000  & 3.42 & $-0.02$ & $-0.59$ & 39\% / 18\% \\
1000--1500 & 3.32 & $-0.06$ & $-0.63$ & 34\% / 18\% \\
1500--2000 & 3.21 & $-0.14$ & $-0.60$ & 36\% / 20\% \\
2000--2500 & 3.11 & $-0.22$ & $-0.57$ & 39\% / 21\% \\
2500--3000 & 3.05 & $-0.28$ & $-0.51$ & 43\% / 22\% \\
\bottomrule
\end{tabular}
\end{center}
The greedy evaluations over the same span won 40, 35, 30, 30, 38 and 43\% of the 40 games
against the random opponent and 25, 20, 10, 13, 15 and 8\% against the heuristic (episodes 500
to 3000). The run had the AAREN history planes switched off (\texttt{AAREN\_ENABLED} false: 30
input planes), so it says nothing about AAREN.

\paragraph{Reading.} The loss is the cross-entropy between the predicted return distribution and
a target built from the network's own target copy; a distribution spread evenly over the 51 atoms
scores $\ln 51 = 3.93$, about where the cold start left it (3.93). Most of the fall is the
network making its distributions sharper and moving its values towards what it actually gets:
the mean $Q$ of the actions taken went from $-0.01$ to $-0.28$, towards the mean return of about
$-0.55$, i.e.\ it learned more accurately that it was losing. Across blocks of 100 episodes the
loss and the training win rate correlated at $-0.49$, but the loss correlated with time at
$-0.97$; once a linear trend was removed from both, the correlation was $-0.14$. The training
games (with NoisyNet exploration on) did improve somewhat against the random opponent, from 34\%
to 43\% of about 215 games per block, roughly two standard errors; the greedy evaluations of 40
games did not show it.

\paragraph{Lesson.} A falling TD or distributional loss shows that the critic fits its own
moving targets better, not that the policy plays better; with a pessimistic, well-calibrated
critic it falls while the agent keeps losing. Two quantities that both drift with training time
correlate whatever their relation, so a loss--performance correlation should be computed on
detrended (or differenced) curves before it is read as coupling. Progress has to be judged on
evaluations against fixed opponents, with enough games for the difference to exceed its noise.

% ------------------------------------------------------------------------------------------
\section{Correcting the claims of the thesis chapters}

\paragraph{Change.} The Methodology and Results chapters were revised where the code or the data
did not support what they said. No reported number was changed; the readings of them were.
\begin{itemize}
  \item AAREN was described as trained end to end with the Q-network through a ``gradient
        bridge'', $\epsilon \cdot W_{proj}(\mathbf{0})$ added to the stored embeddings. An affine
        layer maps the zero vector to its bias, so the DQN loss reached only one of AAREN's 21
        parameter tensors, and none that depends on the move history. AAREN was trained, in
        effect, only by its supervised loss on revealed pieces, and the text now says so. In the
        thesis code its embedding was also taken at the padding step of every shorter history
        rather than at the history's last real step.
  \item The self-play win rate of 49.88\% (against 24.97\% for the DQN) was read as the agent
        doubling its strength. In self-play both sides are the same policy, and in both runs the
        decisive games split almost evenly (24.97/24.99\% and 49.88/43.43\%). The win rate doubled
        because the decisive share did (50\% to 93\%), and steps per win follows from that too.
  \item The loss--win correlation of $-0.787$ was read as the loss becoming a reliable proxy for
        play. It was computed on raw windowed curves. formal\_4 shows the raw correlation
        ($-0.49$) shrinking to $-0.14$ once a linear trend is removed from both curves, while its
        evaluations did not improve.
  \item AAREN's 24.23\% accuracy was compared only with a uniform 8.3\% baseline over 12 types.
        Scouts are 8 of the 40 pieces and 8 of the 33 movable ones, so always naming the most
        common type would score about 20 to 24\%.
  \item The logged average Q-values (0.70, 0.84, 1.84) were read as value resolution. They were
        the maximum over all 3600 actions, the artefact of the second entry of this appendix.
\end{itemize}
The validation table now reads ``Not tested'', ``Not supported'', ``Inconclusive'', ``Supported''
or ``Self-play only'' instead of ``Validated'' for every row. The Methodology chapter now lists all six
Rainbow components and separates the thesis configuration from that of the formal runs.

The Introduction and the Conclusion were revised later, after formal\_14, because they still
stated what the Results chapter had withdrawn. The Introduction said the changes ``doubled the win
rate'' and ``resolved the scalar value collapse and exploration death'', and the Conclusion that
the evidence supported an architectural cause of the baseline's ceiling. Both now report the
self-play figures as a change in the draw rate, give the 11.66\% against heuristic opponents beside
them, and say that the later runs found the agent's level set by its demonstrator. The Conclusion
gained a section that answers the three research questions as far as the data allow (the second
was not demonstrated; the third remains open, with AAREN off in every formal run). Four further
statements were corrected in the same pass:
\begin{itemize}
  \item C51 was said to let the agent tell risky decisive moves from safe passive ones. The agent
        acts on the mean of the distribution, so two moves with the same expected return remain
        tied whatever the head learns.
  \item The zero noise entropy of the normal DQN was read as confirming that its exploration had
        been exhausted. That network had no noisy layers, so the zero is a property of the design.
  \item The lower step count of the Rainbow run (5.17 against 10.31 million) was called sample
        efficiency in the Limitations, against Section~\ref{finding:10}.
  \item $10^{535}$ was given in the Introduction, the Results and the technical background as the
        size of Stratego's state space; the cited source gives it as the size of the game tree
        \cite{perolat2022stratego}. The background chapter derived it from a formula that omitted
        the Bombs and evaluates to about $10^{58}$. It now gives the number of setups, $40!$
        divided by the factorials of the twelve piece counts, about $1.4 \times 10^{33}$ for one
        player and $10^{66}$ for both, and cites the game tree separately.
\end{itemize}
The Expectamax search is now described as not evaluated: no run in this thesis measured its
effect on playing strength.

The Results chapter's example for scalar value collapse was also wrong as arithmetic. It said that
a move with a 40\% chance of winning and a 60\% chance of losing, worth $-0.2$ in expectation,
looks worse to a mean-based estimator than a certain draw worth $-0.3$. It looks better, since
$-0.2 > -0.3$, so the example does not explain the draws, and the section now says so. The abstract
described a study with human participants and usability assessments that no chapter reports, and
aimed at play comparable to human opponents; it now states the evaluation that was done
(agent-versus-agent games) and its results.

\paragraph{Why.} Re-reading formal\_4 against the chapters (previous entry). That run had AAREN
switched off, yet the chapters presented the falling loss and the self-play results as validation
of both the Rainbow learner and AAREN.

\paragraph{Lesson.} A metric from symmetric self-play measures the draw rate, not strength; a
correlation between two curves that both drift with time measures the drift; an accuracy needs
the majority-class baseline, not the uniform one; and a claimed gradient path should be checked
by listing which parameters actually receive a gradient. Each of these was cheap to check and
reversed a claimed validation.

% ------------------------------------------------------------------------------------------
\section{Replay priorities from the cross-entropy}

\paragraph{Observation.} The prioritized replay \cite{RainbowJMLR2021} sets a transition's priority to
$(\ell + 10^{-6})^{0.6}$, with $\ell$ the C51 loss of that transition, which is the cross-entropy
$-\sum_i m_i \log p_i$ between the projected target $m$ and the prediction $p$. The
cross-entropy is the KL divergence plus the entropy of the target,
$\mathrm{CE} = \mathrm{KL}(m \,\|\, p) + H(m)$. Rainbow prioritizes by the KL divergence, the part
that measures how wrong the prediction is. $H(m)$ is large for every non-terminal transition,
because its target is the network's own diffuse prediction one step on. It is at most $\ln 2$ for a
terminal one, whose target is the outcome itself, split over the two nearest atoms.

\paragraph{Measurement.} formal\_4's network of episode 3000 played 32 greedy games (against the
random opponent and the heuristic, 8 won, 16 lost and 8 drawn). They were recorded through the
trainer's 5-step and mirroring pipeline, which gave 13\,328 transitions; the rewards came from the
current material-only potential. For each transition:
\begin{center}
\begin{tabular}{lrccc}
\toprule
Transitions & count & CE & $H(m)$ & KL \\
\midrule
terminal                              & 320   & 4.15 & 0.51 & 3.64 \\
non-terminal, return with a battle    & 6790  & 3.76 & 3.73 & 0.02 \\
non-terminal, no battle               & 6218  & 3.42 & 3.40 & 0.02 \\
\bottomrule
\end{tabular}
\end{center}
With the cross-entropy, the 95th percentile of the priority was 1.2 times the 5th. The terminal
transitions, 2.4\% of the memory, held 2.6\% of the total priority, so they were sampled as
often as under uniform replay. With the KL divergence the ratio would be 10.3. The terminal
transitions would hold 40\% of the priority, about 16 times their share.

The terminal KL of 3.64 also shows that, one move before the end of a game, the network's
distribution put almost no mass near the actual outcome. Those 2.4\% of transitions are the only
ones whose target is the game result rather than the network's own estimate. Under the
cross-entropy priority the replay did nothing to make them count more.

A second, smaller fault: the 2.5-fold priority boost for battles is decided by the newest step of
a 5-step window. It therefore lands on the transition that \emph{starts} four moves before the
battle, not on the one that starts with the attack. The mirrored copy of a transition gets no
boost at all.

\paragraph{Change.} Made for formal\_6 (Section~\ref{sec:formal6}); at the time of this entry, none. The proposed fix is to compute the priority from
$\mathrm{CE} - H(m)$, and to decide the battle boost by the first step of the window. Whether this
changes play is not yet measured in a run.

\paragraph{Lesson.} A priority must measure surprise, not the size of the loss. With a
distributional loss the two differ by the entropy of the target, which here was almost the whole
loss and made prioritized replay uniform. The distribution of the priorities is cheap to check (the
ratio of high to low percentiles, and the share held by the transitions that carry the outcome),
and it should be checked once whenever the loss changes.

% ------------------------------------------------------------------------------------------
\section{An audit of the rest of the Rainbow learner}

\paragraph{Observation.} Before changing the replay priorities, the rest of the learner was read
and, where a doubt could be measured, probed on formal\_4's network of episode 3000.
\begin{itemize}
  \item \emph{Checked and correct:} the C51 projection, the Double DQN target over legal moves,
        the dueling combination, the factorised NoisyNet layers, the per-lane 5-step buffers
        (mirroring is applied after the $n$-step sum, so the two never mix), the flushing of
        open returns at the end of a game, the move-plus-reply transitions, the asynchronous
        learner's locking and weight copy, and the replay ratio (about 8 samples per stored
        transition, the usual DQN ratio).
  \item \emph{Instance normalisation and constant planes.} The backbone normalises every channel
        over the board, and planes 27--29 (observability, no-battle counter, game progress) are
        constant over the board, so they might have been normalised away. They are not. In 300
        fog-of-war positions, raising the no-battle counter from 0 to 0.99 lowered the network's
        value by 0.15 on average, as it should with a draw costing $-0.5$, and game progress
        by 0.11. Removing one own piece changed it by 0.02. A fresh network learnt a target
        from the counter plane equally well with and without the normalisation layers
        (correlation 0.99 in both cases).
  \item \emph{Dropout in the target.} The attention block has dropout 0.1 and the target network
        is never put in evaluation mode, so its targets carry dropout noise. Dropout alone moved
        the value of the greedy move by a standard deviation of 0.015, more than the NoisyNet
        noise (0.010), and about a third of the spread of the values across a position's legal
        moves (0.043).
  \item \emph{How much NoisyNet explores.} With the acting noise on, the move played differed
        from the greedy move in 4 to 6\% of positions (100 and 300 positions, 8 and 20 noise
        samples).
\end{itemize}

\paragraph{Change.} Made for formal\_6 (Section~\ref{sec:formal6}); at the time of this entry, none. The proposed changes are the evaluation mode for the target
network (or no dropout in the attention block) and the priority fix of the previous section.

\paragraph{Lesson.} A suspicion about an architecture is cheap to test: perturb the input it
concerns and see whether the output moves, or check whether a fresh network can fit a target
that depends only on it. Here the plausible suspicion (normalisation erasing the counters) was
wrong, and two quieter faults (an always-on dropout in the target, and priorities that do not
prioritise) were real.

% ------------------------------------------------------------------------------------------
\section{What AAREN can add to the observation}

\paragraph{Why it is off.} The formal runs set \texttt{AAREN\_ENABLED} to false, so the network sees
the 30 observation planes only. No document records the reason; the code enforces two
constraints. The cold start refuses AAREN, because a demonstration has no history of its own and
zero planes would teach the network what a game without history looks like. The asynchronous
learner also refuses it, because AAREN's supervised training shares the network. The plan was to
make the Rainbow learner work first and change one thing at a time.

\paragraph{Measurement.} The heuristic played 225 games as Player 1, alternately against the random
opponent and the heuristic. A history aggregator tracked Player 2's moves and labelled them at
battles exactly as the trainer does.
\begin{itemize}
  \item \emph{The ``has moved'' bit.} Of Player 1's attacks on a piece that had never moved,
        41\% (1472 attacks, against random) and 49\% (1028, against the heuristic) hit a Bomb
        or the Flag. Of its attacks on a piece that had moved, none did (2503 and 1712 attacks).
        The observation does not contain this bit, and the game state does not record which
        pieces have moved.
  \item \emph{Piece type from the history.} AAREN labels only pieces that moved at least twice.
        Its 6216 labelled reveals were split by game, 80\% for training and 20\% for testing,
        and a fresh AAREN was trained for 4000 batches. On the 1272 test reveals it named the
        type correctly in 35.6\% of cases. Always naming the most common type (Lieutenant) gave
        16.3\%, and ``a long move means a Scout, otherwise the most common type'' gave 23.9\%.
        Its accuracy was the same against the random opponent (34.9\%) as against the
        heuristic (36.0\%). Since the random opponent moves without regard to rank, most of
        what AAREN reads must be where a piece stood in the heuristic setup, rather than how it
        behaves. A setup prior is real information against this opponent pool, but it would
        shrink against varied setups.
\end{itemize}

\paragraph{Reading.} AAREN does carry information the observation lacks: it beats both simple rules
on held-out games. The largest single piece of history information, however, is whether a hidden
piece has moved at all. That bit separates a 45\% chance of attacking a Bomb or the Flag from
none, and it needs one observation plane rather than a 64-plane embedding. The thesis accuracy
of 24.23\% is close to the long-move rule here. Whether AAREN helps the \emph{agent} is a separate
question that only paired runs can answer, with and without it, on the same evaluations.

\paragraph{Change.} Made for formal\_6 (Section~\ref{sec:formal6}); at the time of this entry, none. Proposed order: add a ``has moved'' plane for enemy pieces (and mark
a Scout revealed by a long move), then compare AAREN on and off. The cold start and the
asynchronous learner would have to support AAREN for that comparison.

\paragraph{Lesson.} Before adding a learned module to supply missing information, find out what the
information is and how much of it a hand-made feature already captures. Here one bit carried most
of it, and a module trained on revealed pieces never sees the unmoved ones that bit is about.

% ------------------------------------------------------------------------------------------
\section{Moved-piece planes and three learner fixes (formal\_6)}
\label{sec:formal6}

\paragraph{Change.} Four changes, made together before formal\_6:
\begin{itemize}
  \item \emph{Two observation planes, 32 in all.} The environment now records, for every square,
        whether the piece on it has moved and whether it has ever moved more than one square at
        once, which only a Scout can do. Both facts are public: the opponent sees every move. The
        flags travel with the piece, including when it wins a battle, and are cleared when it is
        taken. Plane 30 marks the enemy pieces of unknown rank that have moved, so they are not a
        Bomb or the Flag. Plane 31 marks those that have made a long move, so they are Scouts.
        The GUI bot and the browser player build the same planes, and tests replay games in both
        and compare them plane by plane with the training environment.
  \item \emph{Replay priorities from the KL divergence.} A transition's priority is now
        $(\mathrm{KL}(m \,\|\, p) + 10^{-6})^{0.6}$, with
        $\mathrm{KL} = \mathrm{CE} - H(m)$, as in Rainbow \cite{RainbowJMLR2021}, instead of the
        cross-entropy. The loss itself is unchanged.
  \item \emph{The battle boost on the battle move.} The 2.5-fold priority of a battle now goes to
        the 5-step transition whose \emph{first} move was the attack. The mirrored copy of a
        transition now gets the same boosts as the original, for battles and for wins.
  \item \emph{No dropout in the target.} The target network is put in evaluation mode whenever
        it is copied, so its targets carry no dropout. Its NoisyNet noise stays on, as in
        Rainbow.
\end{itemize}
Everything else is formal\_5's configuration: material-only potential, cold start, asynchronous
learner, 16 lanes, 20\,000 episodes. The first convolution now reads 32 planes, so checkpoints of
formal\_4 and formal\_5 do not fit the new network; formal\_6 starts from zero like them.

\paragraph{Why.} The three preceding sections. The ``has moved'' bit separated a 41 to 49\% chance
that an attacked hidden piece was a Bomb or the Flag from none, and the observation did not contain
it. Under the cross-entropy priority the terminal transitions, the only ones whose target is the
game result, were sampled no more often than under uniform replay; under the KL divergence they
would hold about 16 times their share of the priority. The battle boost of the thesis code
landed four moves before the attack, and dropout moved the target value of the greedy move by
more than the NoisyNet noise did.

\paragraph{Result, cold start.} The cold start ran as before: the heuristic
won 50, lost 27 and drew 23 of its 100 demonstration games (39\,406 transitions), and 5000
pre-training updates took the loss from 4.13 to 3.95 and the margin loss to 0.076. The checks
set for formal\_6 were: whether attacks on
unmoved hidden pieces, and so losses to Bombs, become rarer than in formal\_4 and formal\_5;
whether the Q-value one move before the end of a game moves towards the outcome (formal\_4's
network put almost no mass there, a KL of 3.64); the spread of the priorities in the replay
memory, which should now be wide rather than flat; and whether the network finally reads the
rank of a revealed piece (the rank-swap test and the share of attacks on revealed pieces it
loses, 59 to 63\% for formal\_4 and formal\_5 against 0\% for the heuristic).

\paragraph{Result, to episode 9000.} The greedy evaluations of stage~1 (episodes 500 to 5000, 40
games each) won 17, 18, 10, 18, 11, 23, 14, 18, 15 and 13 games against the random opponent and
6, 5, 6, 13, 2, 7, 11, 8, 4 and 5 against the heuristic: no trend, and no better than formal\_4
or formal\_5. No gate was passed. At episode 5000 the stage's episode cap moved the opponent
schedule to stage~2, whose training opponents are the frozen and the smart heuristic only; the
random opponent and the plain heuristic drop out. The evaluations against the heuristic then won
7, 6, 7, 1, 5, 4, 1 and 7 of 40 (episodes 5500 to 9000). In the training games of stage~2 the
agent won about 17\% against the frozen heuristic and drew about 80\% against the smart one, and
games grew from about 350 to 550 moves, most of them ending in no-battle draws.

The probes of the previous runs were repeated on the networks of episodes 5000 and 9000, in 30
greedy games against the random opponent (and 300 positions for the rank-swap test):
\begin{center}
\begin{tabular}{lccc}
\toprule
 & ep.\ 5000 & ep.\ 9000 & heuristic \\
\midrule
games won / lost / drawn                      & 11 / 15 / 4 & 0 / 3 / 27 & 23 / 4 / 3 \\
attacks per game                              & 34          & 9          & 39 \\
attacks on revealed pieces lost               & 66\%        & 73\%       & 0\% \\
rank-swap: Scout $-$ Marshal, positive in     & 48\%        & 33\%       & -- \\
share of attacks on unmoved hidden pieces     & 27\%        & 0\%        & 31\% \\
\quad of which Bomb or Flag                   & 44\%        & --         & 40\% \\
\bottomrule
\end{tabular}
\end{center}
formal\_5's network (episode 2750, 30 planes) put 30\% of its attacks on unmoved hidden pieces,
37\% of them Bombs or the Flag. So by episode 5000 the moved-piece planes had not changed whom the
agent attacks, and it still did not read the revealed rank (the Scout and Marshal versions of an
attack differed by $-0.001$ on average). By episode 9000 it no longer attacked unmoved pieces at
all, but it attacked hardly anything: it drew 27 of 30 games against the random opponent, which
formal\_6 had beaten 11 times at episode 5000. Its value of an attack was about $-0.45$ whatever the
target, close to the $-0.5$ of a draw.

The C51 loss reported by the trainer fell from 2.47 (first 500 episodes) to 0.38 (last 500), far
below the 3.0 to 3.7 of formal\_4 and formal\_5. The reported loss is the cross-entropy weighted by
the importance-sampling weights of the prioritized replay, and those weights change with the
priorities, so the number is not comparable across the change of priority and shows nothing about
play on its own.

\paragraph{Reading.} Two separate things. In stage~1, the moved-piece planes and the learner fixes
changed nothing that the probes measure: the network still does not learn the outcome of a battle
from what it can see, the failure of formal\_4 and formal\_5. In stage~2, a change of opponents
alone, with the environment unchanged, turned an agent that beat the random opponent about a third
of the time into one that only draws: the agent trained against a pool in which it mostly draws
(the smart heuristic) or loses (the frozen one), its values moved to the value of a draw, and it
stopped attacking. This is the forgetting that the removal of the curriculum's environment changes
was meant to avoid, caused here by the opponent schedule that was kept.

\paragraph{Lesson.} An episode cap that moves an agent to the next stage of an opponent schedule
when it has not passed the gate turns the schedule into a fixed timetable, and dropping the weak
opponents from the pool removes the games the agent can win; with only draws and losses left, a
value learner learns that nothing it does matters and plays for the draw. Keep the earlier
opponents in the pool, and evaluate against them throughout, so that forgetting shows up when it
starts. (formal\_7 kept them, and the collapse happened anyway, so dropping them was not the cause:
Section~\ref{sec:formal7}.)

\paragraph{Lesson.} When a run is stopped to fix several faults at once, the next run cannot say
which fix mattered; it can only say whether they did together. Here that trade was accepted
because each fix restores the standard algorithm rather than adding something new, and the
paired comparison that matters more (AAREN on and off) still has to be run with one difference.

% ------------------------------------------------------------------------------------------
\section{Why the network does not learn the rank of a revealed piece}
\label{sec:revealed-signal}

\paragraph{Observation.} In every formal run the agent lost most of its attacks on enemy pieces
whose rank it could see (59 to 73\%), and its value of an attack did not depend on that rank
(Section~\ref{sec:formal6}). Three explanations were possible: the information does not reach
the training targets, the network cannot represent the rule, or the training signal is too weak
for the learner to pick up. Three probes on formal\_6's network of episode 5000 separate them.

\paragraph{Measurement 1: the targets.} The network played 32 greedy games (6 won, 22 lost, 4
drawn) through the trainer's pipeline, which stored 11\,426 transitions. Each one was classed by
its first move; an attack on a revealed piece was classed won, tied or lost by the battle rule.
For the non-terminal ones:
\begin{center}
\small
\begin{tabular}{lrrccc}
\toprule
First move & count & share & 5-step reward & $\mathbb{E}[m]$ & network $Q$ \\
\midrule
quiet move                   & 9524 & 83.4\% & $-0.016$ & $-0.110$ & $-0.108$ \\
attack on a hidden piece     & 1304 & 11.4\% & $-0.002$ & $-0.035$ & $-0.017$ \\
attack on a revealed piece, won  & 146 & 1.3\% & $+0.023$ & $-0.040$ & $-0.079$ \\
attack on a revealed piece, lost & 388 & 3.4\% & $-0.027$ & $-0.104$ & $-0.074$ \\
\bottomrule
\end{tabular}
\end{center}
The targets do carry the outcome of the exchange: their means differ by 0.064 between won and
lost attacks. The network's $Q$ differs by $-0.005$. The agent also chose attacks on revealed pieces
that it lost 2.7 times as often as ones it won. Each target distribution $m$ is wide, with a standard
deviation of 0.77 for these transitions, because it carries the uncertainty of an outcome hundreds
of moves away. The target \emph{means} of one class varied much less across transitions (0.13).

\paragraph{Measurement 2: capacity.} A fresh network of the same size was trained on 3267
positions from these games. In each, one attack on an adjacent enemy piece was chosen, and the
target square was edited to a revealed piece of random rank. The C51 target was a function of
the battle outcome $o \in \{1, 0, -1\}$ of that attack alone. Every batch of 128 consisted of such
attacks. The measure is the gap between the network's values of the attack when the target is
shown as a Scout (a win) and as a Marshal (a loss), over 300 held-out positions:
\begin{center}
\small
\begin{tabular}{p{7.2cm}lcc}
\toprule
Target & l.r. & Scout $>$ Marshal & gap (ideal) \\
\midrule
$0.5\,o$, no noise & $3\cdot 10^{-5}$ & 100\% & 0.44 (1.00) \\
$0.5\,o$, no noise & $10^{-3}$ & 100\% & 0.97 (1.00) \\
$0.07\,o$ plus sampled noise of s.d.\ 0.77 & $3\cdot 10^{-5}$ & 68\% & 0.003 (0.14) \\
\quad the same with signal $0.35\,o$ & $3\cdot 10^{-5}$ & 94\% & 0.023 (0.70) \\
\quad the same with 12 ``attack wins / ties / loses'' planes & $3\cdot 10^{-5}$ & 72\% & 0.005 (0.14) \\
distributions of s.d.\ 0.77, mean $0.07\,o$ plus noise of s.d.\ 0.13 & $3\cdot 10^{-5}$ & 99\% & 0.028 (0.14) \\
\quad the same plus a squared error on the mean & $3\cdot 10^{-5}$ & 99\% & 0.052 (0.14) \\
\bottomrule
\end{tabular}
\end{center}
(After 2000 batches. ``Scout $>$ Marshal'' is the share of held-out positions in which the attack
is valued higher with the target shown as a Scout.) With a clean target the network learnt the
rule within 1000 batches, even at the run's learning
rate, so it can read a revealed rank and compare it with the attacker's. Planes that state the
outcome of every possible attack outright hardly helped, so finding the relation is not the
difficulty either. What decides is the size of the signal relative to the width of the target. The
win--loss difference of the material potential is under one atom of the support (0.08 of $\pm 2$
over 51 atoms) and a twelfth of the target's width. Under the cross-entropy, a shift of the mean
that small changes the loss by little, so it is learnt slowly. With targets shaped like the real
ones, 2000 batches of nothing but these attacks learnt a fifth of the gap. Adding a squared error
between the predicted and target means doubled that.

\paragraph{Reading.} The rank information reaches the network, and the network can use it. The
training signal for using it is small and diluted. Attacks on revealed pieces were 5\% of the stored
transitions, so the real learner sees about 6 in a batch of 128. Those 6 compete with everything
else the network fits, against a target network that changes every 1000 updates. The probes show
the signal is learnable when concentrated. That dilution and the moving targets account for the
rest is an inference, not a measurement. This is also why the earlier fixes did not change the
behaviour. The material-only potential made the signal clean, but no larger than a fraction of an
atom. The moved-piece planes and the priority fixes do not touch it. Nor would AAREN: it supplies
information about hidden pieces, while this failure concerns pieces that are already revealed.
(This reading was incomplete. formal\_8 made the signal larger and learnt nothing more
(Section~\ref{sec:formal8}). Training this probe inside the run's own learner showed that the
missing ingredient was contrast in the data rather than the size of the signal
(Section~\ref{sec:ladder}); the outlook planes then supplied the comparison directly
(Section~\ref{sec:outlook}).)

\paragraph{Change.} The first and third were made for formal\_7 (Section~\ref{sec:formal7}). The
options the probes support were, by cost: an expected-value loss beside the C51 loss (it doubled
the learnt gap); a larger weight on the material potential with a wider value support (a five-fold
signal got the ordering right in 94\% of positions); and a priority boost for attacks on revealed
pieces, which raises their share of the batch.

\paragraph{Lesson.} A distributional critic learns differences in the mean slowly when they are
small next to the width of the return distribution, even when they are exact. Before relying on
a reward term to teach a decision, compare that term's difference between the right and the wrong
choice with the spread of the targets. Then check, on a fresh network trained on that decision
alone, how much of the difference it learns in a realistic number of updates.

% ------------------------------------------------------------------------------------------
\section{A mean loss, revealed attacks in the batch, and the opponent pool (formal\_7)}
\label{sec:formal7}

\paragraph{Change.} formal\_6 ended at episode 11\,680, when its process was killed from outside
(its log shows no error; it was started from the desktop app's shell, and formal\_7 is started
detached from it). Its evaluations against the heuristic from episode 9500 on were 5, 6, 8, 5 and
5 wins of 40, with no gate passed. formal\_7
starts from zero with formal\_6's configuration and three changes:
\begin{itemize}
  \item \emph{A loss on the mean.} The squared error between the means of the predicted and the
        target distribution, $\bigl(\sum_i p_i z_i - \sum_i m_i z_i\bigr)^2$ with weight 1, is
        added to each transition's C51 cross-entropy. Both are weighted by the importance-sampling
        weight. The priorities are still the KL divergence. The loss reported by the trainer now
        includes this term, so it is not comparable with that of earlier runs.
  \item \emph{Attacks on revealed pieces in the batch.} A transition whose first move attacks an
        enemy piece of known rank, and its mirror image, keeps four times its priority through
        every priority update. The importance-sampling weights leave this factor uncorrected, so
        these transitions count four times as much in the loss, not just four times as often.
        At about 5\% of the replay, they become about 17\% of each batch.
  \item \emph{The opponent pool after stage 1.} Stages 2 and 3 keep the random opponent and the
        heuristic in their pools beside the frozen and smart heuristics (or the self-play
        snapshot). The usual weighting (more games against the opponents it beats less often)
        still gives the random opponent at least 15\%. Every evaluation now also plays the random
        opponent, so forgetting shows as soon as it starts. Stage 1 and the episode caps are
        unchanged.
\end{itemize}

\paragraph{Why.} Section~\ref{sec:revealed-signal}. The win--loss difference of an exchange is
present in the targets but is a twelfth of their width, so the cross-entropy learns it slowly.
A mean loss doubled the learned gap in the probe. The attacks that carry the difference were 5\%
of each batch. Section~\ref{sec:formal6}: once the stage-1 cap moved formal\_6 to a pool it could
only draw or lose against, it stopped attacking.

\paragraph{Result, to episode 2500.} The cold start ran as in formal\_6 (the heuristic won 53,
lost 27 and drew 20 of the 100 games; 5000 pre-training updates, loss 4.17 to 3.94 with the mean
term, margin loss 0.074). Wins in the 40-game evaluations, with formal\_6 at the same episodes in
brackets:
\begin{center}
\begin{tabular}{lccccc}
\toprule
Episode & 500 & 1000 & 1500 & 2000 & 2500 \\
\midrule
random    & 16 (17) & 13 (18) & 12 (10) & 23 (18) & 17 (11) \\
heuristic & 4 (6)   & 4 (5)   & 10 (6)  & 7 (13)  & 7 (2) \\
\bottomrule
\end{tabular}
\end{center}
Summed over the five evaluations, formal\_7 won 81 games against the random opponent to formal\_6's
74, and 32 against the heuristic to 32. One evaluation of 40 games has a standard error of about
3 wins at these rates, so stage 1 shows no difference yet. That is expected: two of the three
changes act on stage 2 and later, and the mean loss acts on a signal the evaluations measure only
indirectly. The test is still the rank-swap probe and the share of revealed attacks lost, and,
after episode 5000, whether the evaluations against the random opponent hold up.

\paragraph{Result, the revealed-rank probe at episodes 2500 and 3500.} Each checkpoint of both
runs played 32 greedy games (same seed). The attacks on revealed pieces in them were classed by
the battle rule, and the rank swap of Section~\ref{sec:revealed-signal} was run on one fixed set
of 300 positions for all four checkpoints:
\begin{center}
\small
\begin{tabular}{lcccc}
\toprule
 & formal\_6, 2500 & formal\_7, 2500 & formal\_6, 3500 & formal\_7, 3500 \\
\midrule
attacks on revealed pieces & 478 & 484 & 512 & 470 \\
\quad won / tied / lost & 26 / 10 / 64\% & 32 / 6 / 62\% & 30 / 9 / 62\% & 29 / 8 / 63\% \\
Scout $>$ Marshal & 65\% & 64\% & 66\% & 67\% \\
mean gap & 0.021 & 0.020 & 0.032 & 0.022 \\
$\mathbb{E}[m]$ gap, won $-$ lost & 0.088 & 0.111 & 0.101 & 0.115 \\
\bottomrule
\end{tabular}
\end{center}
By episode 3500 (about 650\,000 updates) the two runs do not differ. Both still lose about
two-thirds of the attacks they choose on pieces whose rank they can see, and both rank the Scout
above the Marshal in two-thirds of the positions, with a gap of a fifth of the target's 0.1. The
targets carry the difference as before. So the mean loss and the four-fold priority have not yet
made the network learn it faster in the real run. The probe had predicted twice the learning
speed on a batch made only of these attacks. In the run they are about a sixth of the batch,
and the targets move. The value of every attack was also lower in formal\_7 at episode 3500
($-0.22$ to $-0.26$, against $-0.05$ to $-0.12$ in formal\_6). That checkpoint comes after the
resume, which emptied the replay of its own games, so the cause is not clear.

\paragraph{Result, to episode 5500: stage 2 collapses as before.} Stage 1 ended at its cap of
5000 episodes. Its last evaluations were 15, 18 and 17 wins of 40 against the random opponent
and 7, 6 and 6 against the heuristic (episodes 4000 to 5000), as in formal\_6. 500 episodes into
stage 2, the greedy agent drew 37 of 40 games against the random opponent, won none and lost 3.
It won 7 of 40 against the heuristic. Episodes slowed from about 1.9 to 6.9 seconds, as the
games ran on to the draw rules. The random opponent was in the stage-2 pool: it got 112 of the
500 training games. The agent won the first of them about half the time and then stopped winning
against it. So keeping the opponents it can beat in the pool did not prevent the collapse;
Section~\ref{sec:formal6}'s explanation, that the agent forgot them, is wrong.

Only the opponents change at the stage switch: the game, its limits and the shaping stay the
same. A reading that fits the numbers is that stage 2 made avoiding battle the better policy. Against
the frozen and smart heuristics, three-quarters of the stage-2 pool, the agent won about 15\% and
lost about 70\% of its training games. A game it fought out was then worth about
$0.15 - 0.70 = -0.55$, no more than a draw ($-0.5$). An exchange it lost also cost material
shaping at once, while avoiding battle cost nothing until the game ended. The observation does
not say which opponent the agent is playing, so the same passive policy was played against the
random opponent. The underlying weakness is the one measured above: it loses two-thirds of the
attacks it chooses on pieces whose rank it can see. Until that changes, fighting does not pay
against a competent opponent.

\paragraph{Lesson.} When a curriculum moves the agent to opponents it mostly loses to, check the
value of a draw against the value of playing on. If the draw is worth as much, the agent will learn
to draw, and it will do so against every opponent it cannot tell apart from them.

The process was killed at about episode 3000 by a console window-close event
(Section~\ref{sec:engineering}). It was resumed from the checkpoint of episode 2750 with the
networks, the target network, the optimiser and the curriculum state, but not the replay
buffer. On resuming, the trainer played 100 new heuristic games (48 won, 24 lost, 28 drawn;
36\,214 demonstration transitions) to refill it, and skipped the pre-training updates. Its own
transitions of episodes 1 to 2750 are gone, so for a while after the resume the replay holds a
larger share of demonstrations than it did before. Results after episode 2750 are of a resumed
run.

\paragraph{Lesson.} When a probe shows that a learner cannot pick up a small but exact signal, the
fix should act on that signal, not on the agent's inputs. Two levers do: the loss, so the signal
weighs more against the noise it competes with, and the data, so the cases that carry it are a
larger share of each batch.

% ------------------------------------------------------------------------------------------
\section{What the agent has to learn, and in what order}
\label{sec:skill-ladder}

\paragraph{Observation.} Read together, the formal runs show the agent stuck at one step of a
sequence of skills that Stratego asks for. Each skill pays off only through the one before it,
and the opponent schedule had assumed skills the agent did not have.
\begin{enumerate}
  \item \emph{The rules.} Which moves are legal. The action mask supplies this; nothing has
        to be learned.
  \item \emph{Engagement.} Moving pieces forward and making contact, instead of shuffling until
        a draw rule ends the game. Learned in every run since the draw was given a cost: in stage 1
        the agent attacks, and it wins 30 to 58\% of its evaluation games against the random
        opponent, about 40\% on average.
  \item \emph{Comparing two known ranks.} Attacking a revealed piece it beats and leaving or
        fleeing one that beats it. This is the whole game under full information, and it needs
        no inference. It is \emph{not} learned: from formal\_4 to formal\_7 the agent lost 59 to
        73\% of the attacks it chose on revealed pieces, and its values barely depend on the rank it
        sees (Sections~\ref{sec:revealed-signal} and~\ref{sec:formal7}). From formal\_9 on the
        observation states each attack's outcome (Section~\ref{sec:outlook}). An agent with this skill
        should beat a random opponent nearly every time, because the random opponent offers its
        pieces at random. Winning 40\% of those games is itself a sign that the skill is missing.
  \item \emph{Inference about hidden pieces.} A piece that has moved is neither Bomb nor Flag. A
        piece that moved more than one square is a Scout. A low piece can be sent to make a
        strong one reveal itself. formal\_6 gave the network these facts as planes 30 and 31.
        This skill pays off only through skill~3: knowing more about a piece is worth something
        only to an agent that attacks the right pieces once it knows.
  \item \emph{Managing information.} Keeping its own strong pieces hidden, protecting revealed
        ones, and the Spy against the Marshal. It pays off by denying the opponent skills 3
        and 4.
  \item \emph{Strategy.} Finding and defending the Flag, tempo, the setup and the endgame. Only
        the outcome rewards it.
\end{enumerate}

\paragraph{Reading.} The reward signal for each skill weakens down the list. The material
potential rewards skill~3 at once, as each exchange is won or lost. Skill~4 gets no immediate
reward; it is paid later, through better exchanges. Skill~6 is paid only by the outcome, hundreds
of moves later. This is why skill~3 is the bottleneck. It has the most direct signal of the skills
that have to be learned, and everything above it is built on it. (Why that signal is not learned
is still open. It is not only its size: formal\_8 made it two and a half times larger in the
targets, and the network still did not respond to the rank; Section~\ref{sec:formal8}.)

The opponent schedule assumed the skills in a different order. Stage 2's frozen and smart
heuristics play skills 3 to 5 well, so playing them presumes the agent can already compare ranks.
The stage-1 gates measured win rates against random and the heuristic, which mix every skill
together. And the episode cap moved the agent on whether or not it had passed them. An agent
without skill~3 that meets opponents with it loses most exchanges. It can then do no better than
draw, which is what formal\_6 and formal\_7 learned to do.

\paragraph{Lesson.} Name the skills a game requires, the order in which their rewards depend on
each other, and a direct measure of each. Then advance a curriculum on the measure of the skill the
next stage needs, not on a win rate that mixes them all. The direct measure of skill~3 is the
share of attacks on revealed pieces that the agent loses in its own games. The evaluations report
it from formal\_8 on.

% ------------------------------------------------------------------------------------------
\section{A five-fold material signal and stages without a cap (formal\_8)}
\label{sec:formal8}

\paragraph{Change.} formal\_7 was stopped at episode 5690, in stage 2. formal\_8 starts from zero
with formal\_7's configuration (mean loss, priority boost on attacks on revealed pieces, earlier
opponents kept in the pool) and three changes:
\begin{itemize}
  \item \emph{The material potential is five times larger} (weight 5 instead of 1). The support of
        the value distribution is widened from $\pm 2$ to $\pm 6$ (0.24 per atom) so that it still
        covers every return (outcome plus $-\Phi(s_0)$ with $\Phi \in [-5, 5]$). A won exchange
        of average size now moves the target by about 0.3, more than an atom, against the same
        outcome noise as before.
  \item \emph{No episode cap on any stage.} A stage ends only when its evaluation gates are
        passed (stage 1: 70\% against the random opponent and 50\% against the heuristic).
  \item \emph{A direct measure of skill~3 in the evaluations.} Every evaluation game counts the
        agent's attacks on revealed pieces and how many it won or lost by the battle rule; the log
        shows the share lost beside each opponent's result.
\end{itemize}

\paragraph{Why.} Section~\ref{sec:skill-ladder}. The capacity probe of
Section~\ref{sec:revealed-signal} put the rank-swap ordering right in 94\% of positions when the
signal was five times larger, against 68\% at the run's size. Stage 1 is the only stage in which
no run has collapsed, so the agent stays there until it has the skill the next stage requires.

\paragraph{Caveat.} The shaping is potential-based, so it leaves the optimal policy unchanged
\cite{NgPBRS}. The optimal shaped values are the unshaped ones minus $\Phi(s)$, a term that
depends only on the state, so in the limit the gap between two actions is the same at any weight.
A larger potential makes the exchange visible in the early targets, before the bootstrapped values
have learned $-\Phi$. It speeds up learning and does not change what is learned. This is the
intended effect: the probes showed the agent learning skill~3 too slowly, not learning the wrong
thing.

\paragraph{Result.} The cold start ran as before (the heuristic won 51, lost 23 and drew 26 of
the 100 games; 5000 pre-training updates, loss 4.40 to 4.00, margin loss 0.084). The run was
stopped at episode 2386. Its evaluations, with formal\_7 at the same episodes in brackets:
\begin{center}
\small
\begin{tabular}{lccccc}
\toprule
Episode & 500 & 1000 & 1500 & 2000 & sum \\
\midrule
wins against random    & 10 (16) & 11 (13) & 18 (12) & 18 (23) & 57 (64) \\
wins against heuristic & 4 (4)   & 12 (4)  & 5 (10)  & 5 (7)   & 26 (25) \\
revealed attacks lost, against random    & 63\% & 69\% & 62\% & 60\% & \\
revealed attacks lost, against heuristic & 77\% & 72\% & 72\% & 78\% & \\
\bottomrule
\end{tabular}
\end{center}
The wins do not differ from formal\_7's. The new measure of skill~3 does not improve over the
run. In each 40-game evaluation the agent attacked 450 to 510 revealed pieces against the random
opponent and 160 to 180 against the heuristic, and it lost 60 to 78\% of those attacks.

The probe of Section~\ref{sec:formal7} on the checkpoint of episode 2000 separates the target from
what the network learned. The won--lost gap of the target means was 0.25 (about 0.1 in formal\_7):
the five-fold potential did make the signal in the targets two and a half times larger. The
network's $Q$ gap between won and lost attacks was 0.06. In the rank swap the network valued
the attack higher on a revealed Scout than on a revealed Marshal in 53\% of the 300 positions,
with a mean gap of 0.003. That is no better than chance, and below formal\_7's 66\% at the same
episode. (The probe ran with formal\_8's support and potential. Only the 66\% is comparable for the
formal\_7 checkpoint; its values are on the wrong scale.)

So the size of the signal relative to the target's width was not what kept the network from
learning skill~3. Section~\ref{sec:revealed-signal} explained the failure that way, from a fresh
network trained on batches made only of these attacks; that network learned the rule. In the
run, a signal two and a half times larger left the network's response to the rank at chance.
Something that differs between the probe and the run stops it: the mix of the batch, the moving
bootstrapped targets, the NoisyNet head, or the rest of the training pipeline. Which one is the
next question. (Answered in Section~\ref{sec:ladder}: none of them stops it; the run's data lack the
contrast between ranks that the probe's random targets have.) A direct test is to train the probe's fresh network on batches drawn from the
run's own replay buffer, adding one difference at a time, and watch the rank swap.

\paragraph{Lesson.} A probe that isolates one factor can show that the factor is \emph{sufficient}
to explain a failure, not that it is the explanation. The capacity probe showed a small signal
learnt slowly; the run with a larger signal learnt nothing more. Before scaling a fix up to a full
run, reproduce the failure inside the probe by adding back the run's conditions one at a time,
until the probe fails as the run does.

% ------------------------------------------------------------------------------------------
\section{Checking the inputs before the learner}
\label{sec:inputs-check}

\paragraph{Check.} Before looking further for why skill~3 was not learnt
(Section~\ref{sec:formal8}), the inputs it depends on were checked. In 30 games between two
random players (26\,450 plies), both players' visible boards were compared with the real board
after every ply. For every attack on a square whose rank plane was set, the outcome predicted
from the planes was compared with the real battle. Three baselines were also measured, 40 games
each under the evaluation rules: random against random, the heuristic against random, and the
heuristic against itself.

\paragraph{Result.} No square showed a piece that was not there, a wrong rank, or a hidden marker
on an empty square. In all 457 attacks on revealed pieces the planes predicted the real outcome. The
random player lost 70\% of those attacks (320 lost, 32 tied, 105 won). A revealed piece is a
survivor of an earlier battle, so it is usually strong or a Bomb. The 60 to 78\% that formal\_4 to
formal\_8 lost is therefore what choosing attacks at random gives. The baselines:
\begin{center}
\small
\begin{tabular}{lccc}
\toprule
Pairing (40 games) & wins / losses / draws & revealed attacks lost \\
\midrule
random against random & 6 / 6 / 28 & 65\% \\
heuristic against random & 29 / 8 / 3 & 0\% \\
heuristic against heuristic & 8 / 16 / 16 & 0\% \\
\bottomrule
\end{tabular}
\end{center}
formal\_8 won 10 to 18 and lost 18 to 23 of its 40 evaluation games against the random player
between episodes 500 and 2000. It lost to random more often than random loses to random:
engagement without skill~3 is worse than not fighting.

\paragraph{Lesson.} Before hunting a learning failure, check that the inputs the decision needs
are correct, and measure what a random choice scores on the same measure. Here the inputs were
exact, and the agents' ``skill'' measure sat at the random player's value.

% ------------------------------------------------------------------------------------------
\section{The run's learner on the capacity probe}
\label{sec:ladder}

\paragraph{Question.} Section~\ref{sec:formal8} proposed training the probe's network inside the
run's own learner, adding the run's conditions one at a time until it fails as the run does.

\paragraph{Measurement.} 64 heuristic games gave 22\,834 transitions (with their mirror images)
and 7\,014 positions with an attack available; 300 were held out. Synthetic attacks were made as in
Section~\ref{sec:revealed-signal}: the target square edited to a revealed piece of random rank,
a terminal transition rewarded $0.3\,o$ for the battle outcome $o$, with the mirror image and the
fourfold revealed-attack priority. A fresh agent was trained with the run's update step
unchanged (prioritized replay, mixed precision, NoisyNet, the mean loss, learning rate
$3\cdot10^{-5}$, batches of 128). The measure is the rank swap of Section~\ref{sec:formal7}.
\begin{center}
\small
\begin{tabular}{lcccc}
\toprule
Replay memory & 250 & 1000 & 2000 & 4000 updates \\
\midrule
synthetic attacks only & 82\% & 100\% & 100\% & \\
5\% synthetic among real transitions & 64\% & 80\% & 82\% & 82\% \\
\quad the same, NoisyNet off & 64\% & 81\% & 77\% & 75\% \\
20\% synthetic among real transitions & 67\% & 96\% & 99\% & 97\% \\
\bottomrule
\end{tabular}
\end{center}
With synthetic attacks only, the learnt gap reached 0.57 of the ideal 0.60. At 5\% it was about
0.15, at 20\% 0.45 to 0.49. A second probe asked whether the network can imitate the heuristic: on
2000 held-out positions (17.7 legal moves on average, chance 5.7\%) the cold start's update step
agreed with the heuristic's move in 47\% of positions after 500 updates and 53\% after 5000. A
cross-entropy over the legal moves at learning rate $3\cdot10^{-4}$ reached 53 to 55\%; the
heuristic breaks ties at random, so its ceiling is below 100\%.

\paragraph{Reading.} The update step, the network, NoisyNet and mixed precision all learn the rule
quickly when the batch shows it; dilution slows it, and NoisyNet plays no part. Even at 5\% the
rule was learnt to about 80\%, while the run sat at chance (53\%). What the probe has and the run's
data lack is \emph{contrast}. In the probe the target's rank is drawn at random; in real games a
revealed piece next to one's own is almost always a strong survivor or a Bomb. A revealed Scout
beside one's Miner hardly ever occurs, so the run's data barely show that the rank matters. The
rank swap asks about positions the run never trains on. This also explains formal\_8: a larger
reward for an exchange does not create the missing variation.

\paragraph{Lesson.} A decision is learnt from the variation the data show. When the variation that
matters is rare in the agent's own games, scaling its reward does not help; it has to reach the
learner through the data or through the input.

% ------------------------------------------------------------------------------------------
\section{Battle-outlook planes (formal\_9)}
\label{sec:outlook}

\paragraph{Change.} The observation grew from 32 to 41 planes. For every own piece and each of
the four directions of the action encoding, two planes give the probability that attacking the
enemy piece it would reach wins and that it loses (a Scout reaches the first piece along the
line, every other piece the adjacent square). A ninth plane gives, for every own piece, the
largest probability that an adjacent enemy piece would beat it by attacking (a Bomb or the Flag
never attacks). For a revealed target the probability is the battle rule. For a hidden one it is
the rank distribution of its class: the rank counts of all hidden enemy pieces are public, since
every captured piece was revealed in its battle. A piece that made a long move is a Scout, a piece
that moved is not a Bomb or the Flag, and the unmoved pieces share what is left. The environment
now records these counts with each state; the GUI bot computes them from its board, and the
browser engine computes all nine planes itself. The thesis's observation and the earlier probes
of Section~\ref{sec:revealed-signal} gave the network only the ranks.

\paragraph{Why.} Section~\ref{sec:ladder}. Skill~3 of Section~\ref{sec:skill-ladder} needs no
inference, only the battle rule; like the action mask for skill~1, the observation can supply it.
The per-square advantage head then reads, at the square of each piece, how each of its attacks
would end. The capacity probe had found that planes stating the outcome hardly helped a network
trained on synthetic attacks alone (72\% against 68\%); in the run the question is a different
one, whether the agent's own games teach it to use the outcome.

\paragraph{Check.} In 40 random games all 636 attacks on revealed pieces ended as the planes said.
In 832 attacks on hidden pieces the mean predicted chance of winning was 0.40 and 36\% were won;
by bins of the predicted chance (0--0.2, 0.2--0.4, 0.4--0.6, 0.6--0.8, 0.8--1) 6, 18, 41, 58 and 92\%
were won. The browser engine's planes matched Python's within $10^{-4}$ in six replayed games.
Computed with precomputed ray indices and two gathers, the planes take 3.7\,ms for 16 states on
the GPU (14.6\,ms for the first version, which shifted the board once per square of each ray).

\paragraph{Result.} formal\_9 is formal\_8 with the 41 planes. At episode 500 it beat the random
opponent 17--12 (formal\_8: 10--23) and lost 46\% of its attacks on revealed pieces against it
(63\%); against the heuristic it lost 44\% of them (77\%). A probe of that checkpoint in 12 greedy
games against random: of 1866 positions with an attack the planes mark as a sure loss, the agent
chose one in 64; of 272 with a sure win available it took it in 86. Its values ranked the classes
correctly (sure-loss attacks $-0.128$, quiet moves $-0.078$, sure wins $-0.043$), with gaps about
half the immediate material reward (a revealed Sergeant is worth $5 \times 4/148 = 0.14$). The
measure did not improve further from episode 500 to 1500 (44 to 46\%).

\paragraph{Lesson.} When a decision depends on a rule applied to information already in the
observation, put the rule's result in the observation. The learner is then left with what the
result is worth, which is the part that needs the reward.

% ------------------------------------------------------------------------------------------
\section{A higher learning rate (formal\_10)}
\label{sec:formal10}

\paragraph{Change.} The learning rate rose from $3\cdot10^{-5}$ to $10^{-4}$; everything else is
formal\_9.

\paragraph{Why.} formal\_9 ranked the moves correctly but by small gaps, and the probes of
Section~\ref{sec:revealed-signal} learnt small differences of the mean slowly at
$3\cdot10^{-5}$.

\paragraph{Result.} At episode 500 formal\_10 beat random 24--10 and the heuristic 11--15 (formal\_9:
17--12 and 3--22), losing 40 and 38\% of its revealed attacks. In the same probe, sure wins were taken
in 45\% of positions (formal\_9: 32\%) and the gap between sure wins and quiet moves was 0.125
(0.035). Most of the sure losses it still chose were deliberate: of 33 in six games, 20 were made
by a Scout and 16 by a piece the threat plane marked as about to be taken anyway. A Scout's
sacrifice also resets the no-battle counter. The share of revealed attacks lost therefore counts
cheap sacrifices as mistakes; sure-loss attacks by a non-Scout that is not threatened would be a
cleaner measure of skill~3. Against random it then stayed level (21--16 at episode 1000, 21--11 at
1500), and all of its 10 losses in 30 games at episode 1000 were by running out of movable pieces,
with 6.9 own pieces left against 11.4 (Section~\ref{sec:formal11}).

\paragraph{Lesson.} A measure of a skill should count only the errors of that skill. Losing an
attack is a mistake for a Captain that had a safe move, not for a Scout sent to spend itself.

% ------------------------------------------------------------------------------------------
\section{A potential that paid for attrition (formal\_11)}
\label{sec:formal11}

\paragraph{Observation.} formal\_10 lost to the random opponent only by running out of movable
pieces. The material term of the potential explained part of it. A hidden enemy piece counted as
the mean rank of an army, 3.7, and a revealed Bomb or Flag as 0. A Scout (2) that attacked a
hidden Bomb and died therefore \emph{raised} the material term by $3.7 - 2 = 1.7$; a Sergeant (4)
lowered it by only 0.3. Any lost attack on a hidden piece weaker than 3.7 gave back part of the
loss in the same way. Throwing cheap pieces at unmoved pieces, 41 to 49\% of which are Bombs or the
Flag (Section~\ref{sec:formal6}), was rewarded.

\paragraph{Change.} The enemy's material is now computed from the public rank counts of its hidden
pieces, so it is exact: a reveal leaves the potential unchanged, and only captures move it. A
test checks that a Scout lost to a hidden Bomb lowers the potential by exactly $5 \times 2/148$.

\paragraph{Result.} At episode 500 formal\_11 beat random 23--15 and the heuristic 12--12 with 30
and 31\% of revealed attacks lost; at 1000 23--12 and 10--13: about formal\_10's level. It still lost
11 of 30 games against random by running out of movable pieces. Counted over 20 games, it attacked
about 40 times a game: it won 8.4 attacks on hidden moved pieces, 4.0 on hidden unmoved ones and 6.3
on revealed ones, and lost or tied 7.8, 8.3 (3.35 of them to Bombs) and 4.8 (1.0 to a revealed Bomb).
Another 6.2 pieces a game were taken by the opponent's attacks. It spent about 21 of its 33 movable
pieces in its own attacks.

\paragraph{Lesson.} Potential-based shaping leaves the optimal policy unchanged \cite{NgPBRS}, but a
learner far from convergence follows the immediate shaping. A potential that estimates what it
cannot see must not reward revealing it: compute the hidden part exactly from public information
where possible, or make the estimate unbiased for the pieces the agent actually attacks.

% ------------------------------------------------------------------------------------------
\section{A mobility term in the potential (formal\_12)}
\label{sec:formal12}

\paragraph{Change.} The potential gained the term $g(n_{\mathrm{own}}) - g(n_{\mathrm{enemy}})$, with
$n$ the number of a side's movable pieces and $g(n) = \ln(1+n)/\ln 34$, at weight 1. It is concave,
so the last movable piece costs 0.20, the second-last 0.12, the fifth-last 0.05 and one of 33 only
0.008, on top of its material. Both counts are public. The potential then lies in $[-6, 6]$, so the
support of the value distribution was widened to $\pm 7$ (0.28 per atom).

\paragraph{Why.} Section~\ref{sec:formal11}: a side without a movable piece loses, but the material
term priced the last Sergeant like the first, and an even exchange was worth nothing.

\paragraph{Result.} At episode 500 formal\_12 beat random 30--7 (75\%), the first run past the 70\%
gate of stage~1; against the heuristic it won 10, lost 12 and drew 18. In 30 games against random it lost 5 by
running out of movable pieces (formal\_11 at the same episode: 11). Against the heuristic, 7 of 30
games were lost by the capture of its own Flag (after 212 turns on average, with 13 own pieces left)
and 12 were no-battle draws. After episode 500 it did not improve: random 27--10 at episode 1000,
17--19 at 1500, 26--10 at 2000; the heuristic 8--14, 12--19 and 9--13. Section~\ref{sec:cold-start-level}
shows why.

\paragraph{Lesson.} Value a resource by what losing it does to the game, not by a linear count: a
concave term in the potential makes the last units count most without changing the optimal policy.

% ------------------------------------------------------------------------------------------
\section{The cold start sets the level}
\label{sec:cold-start-level}

\paragraph{Observation.} formal\_12's training games did not improve from the first 250 episodes on.
Against the heuristic it won 36, 36, 38, 36, 39 and 32 games of each 250-episode block (about 23\%
of those games), while its training loss fell from 2.22 to 1.42. In 1500 episodes the learner made
33\,585 updates, 22 per episode at about 55\,ms each; with a target copy every 1000 updates, the
outcome had been bootstrapped through about 34 target networks.

\paragraph{Measurement.} A fresh agent with formal\_12's code was given only the cold start: 100
heuristic games (35\,198 transitions) and pre-training, without a single training episode. In 40
evaluation games against each opponent:
\begin{center}
\small
\begin{tabular}{lcc}
\toprule
Agent & against random & against the heuristic \\
\midrule
cold start only, 5000 updates (as in every run) & 27 / 6 / 7 & 7 / 13 / 20 \\
cold start only, 20\,000 updates & 23 / 14 / 3 & 9 / 12 / 19 \\
formal\_12, episodes 500 / 1000 / 1500 / 2000 & 30, 27, 17, 26 wins & 10, 8, 12, 9 wins \\
\bottomrule
\end{tabular}
\end{center}
(wins / losses / draws.) All the runs of this section's sequence, at equal episodes:
\begin{center}
\small
\begin{tabular}{llcccc}
\toprule
Run & episode & against random & lost & against the heuristic & lost \\
\midrule
formal\_8  & 500  & 10 / 23 / 7  & 63\% & 4 / 23 / 13  & 77\% \\
formal\_9  & 500  & 17 / 12 / 11 & 46\% & 3 / 22 / 15  & 44\% \\
formal\_9  & 1500 & 25 / 8 / 7   & 44\% & 7 / 19 / 14  & 47\% \\
formal\_10 & 500  & 24 / 10 / 6  & 40\% & 11 / 15 / 14 & 38\% \\
formal\_10 & 1500 & 21 / 11 / 8  & 35\% & 10 / 15 / 15 & 31\% \\
formal\_11 & 500  & 23 / 15 / 2  & 30\% & 12 / 12 / 16 & 31\% \\
formal\_11 & 1000 & 23 / 12 / 5  & 34\% & 10 / 13 / 17 & 27\% \\
formal\_12 & 500  & 30 / 7 / 3   & 28\% & 10 / 12 / 18 & 40\% \\
formal\_12 & 1000 & 27 / 10 / 3  & 41\% & 8 / 14 / 18  & 41\% \\
\bottomrule
\end{tabular}
\end{center}
(``Lost'': the share of the agent's attacks on revealed pieces that it lost.) With 40 games the
standard deviation of a win rate near one half is about 8 percentage points, so most of the
differences after formal\_9 lie within noise.

\paragraph{Reading.} The reinforcement learning phase added nothing measurable in 2000 episodes.
The agent plays at the level of the heuristic it imitated (the heuristic beats random 29--8 and
itself 8--16, Section~\ref{sec:inputs-check}). The gains since formal\_8, from 10--23 to about 27--8
against random, are a better starting point: the outlook planes let the network take up the
heuristic's comparisons of ranks. Whether formal\_8's own cold start started that low or its
reinforcement learning pulled it down was not measured. Pre-training four times longer on the same
35\,198 transitions did not help; against random it did worse.

\paragraph{Lesson.} Measure the starting point alone before crediting a run's level to its
reinforcement learning. While the learning signal is this weak, the demonstrator sets the ceiling.

% ------------------------------------------------------------------------------------------
\section{A stronger teacher: the outlook expert}
\label{sec:expert}

\paragraph{Change.} A new scripted player, the outlook expert, can play the demonstrations and
serve as an opponent. It looks one move ahead and reads only public information, through the
network's own observation and the rank distribution of every enemy piece. Each legal move gets a
score in units of piece value (Marshal 11, General 8, Colonel 6, Major 4.6, Captain 3.4, Lieutenant
2.6, Sergeant 2, Scout 1.2, Spy 3, Miner 3, Bomb 1, Flag 25). An attack scores the expected value of
the battle over the target's rank distribution, plus the risk of standing on the square afterwards.
Every move scores the change in risk to the moved piece: the strongest adjacent enemy's chance to
beat it, times its value, times 0.6. Cheap pieces advance. Pieces are drawn toward the unmoved
hidden pieces, where the Flag can be, more strongly in the endgame (15 enemy pieces or fewer) and as
the no-battle counter runs. Attacks near the own Flag count more, and a piece guarding the Flag
stays. Its score terms explain each move it makes.

\paragraph{Why.} Section~\ref{sec:cold-start-level}: the teacher sets the level.

\paragraph{Result.} Under the evaluation rules (100 games against each, about 0.4\,s a game on the
CPU) the expert beat random 86--3 with 11 draws and the heuristic 79--11 with 10 draws, and lost none
of its attacks on revealed pieces. The Flag is valued at 25 rather than at a game-winning amount:
early on its probability is spread over some 30 unmoved pieces, and a large value would make every
one of them look worth attacking. At 25 an early probe is worth little and the last few unmoved
pieces a great deal. (This value was chosen by that argument; other values were not tried.)

A cold start from 400 expert games (the expert won 350 of them; 126\,368 transitions), with the
trainer's pre-training and no training episodes, gave in 40 games against each opponent:
\begin{center}
\small
\begin{tabular}{lcccc}
\toprule
Pre-training updates & random & heuristic & expert & revealed attacks lost \\
\midrule
5000   & 27 / 11 / 2 & 16 / 16 / 8 & 7 / 27 / 6 & 41 to 54\% \\
10\,000 & 28 / 10 / 2 & 21 / 15 / 4 & 7 / 31 / 2 & 40 to 55\% \\
20\,000 & 24 / 16 / 0 & 21 / 16 / 3 & 9 / 30 / 1 & 39 to 57\% \\
\bottomrule
\end{tabular}
\end{center}
Against the heuristic the agent more than doubled its wins (7 to 9 with the heuristic as teacher).
But it lost about half of its attacks on revealed pieces, where its teacher lost none: it took up
the expert's appetite for battle, not its choice of battles. (In the runs taught by the heuristic
the large-margin loss ended pre-training at 0.063 to 0.084, against its margin of 0.1, so the
pre-training did not rank the demonstrator's move first by the margin; it was not recorded here.) A behaviour
cloning loss (the cross-entropy of the expert's move over the legal moves) was measured in the
subsequent probe (Section~\ref{sec:behaviour_cloning}).

\paragraph{Lesson.} A better teacher raises the starting point, but only as far as the imitation
loss transfers the teacher's decisions. Check that the student makes the teacher's key decisions,
not only that it wins more often.

% ------------------------------------------------------------------------------------------
\section{Playing in the browser, with reasons}
\label{sec:browser}

\paragraph{Change.} The browser player (\texttt{web/play.html}) now runs formal\_12's network of
episode 500 (exported to ONNX, 41 planes, 6.1\,MB) and computes the outlook planes itself. Next to
each of the bot's six best moves it shows the move's value and a reason read from the planes the
network saw: ``attack: wins'', ``attack: loses'', ``attack: 5\% win / 75\% lose'', or ``threatened
piece moves away''. One network call per move runs in the page, with no server.

\paragraph{Check.} The page loaded the model with 41 input planes. The bot's first move was a7--a4,
a Scout's long attack shown as ``Q 0.032, attack: 5\% win / 75\% lose'' above the quiet moves (Q
$-0.05$). The Scout met a Lieutenant: a scouting sacrifice that revealed a piece.

\paragraph{Lesson.} An observation built from explicit quantities makes a learnt policy easier to
explain: the same planes that helped it learn are the vocabulary of its reasons.

% ------------------------------------------------------------------------------------------
\section{The teacher's moves: behaviour cloning on legal actions}
\label{sec:behaviour_cloning}

\paragraph{Change.} The demonstration loss is replaced with behaviour cloning over the state's
legal moves:
\[
L_{\text{BC}} = -\log \frac{\exp(Q(s, a_E) / \tau)}{\sum_{a \in \mathcal{A}(s)} \exp(Q(s, a) / \tau)}
\]
with temperature $\tau = 0.25$ and weight $1.0$ (disabling the large-margin hinge loss, which
punished only the single argmax runner-up and left illegal or suicidal moves unpenalised). At the
same time, dropout in the online spatial self-attention block (which stayed active in train mode
during replay updates) was set to $0.0$, eliminating rank-order jitter during replay and Double DQN
target evaluations. A third, unused copy of the network (\texttt{reference\_network}) was also made
conditional on KL regularisation being enabled.

\paragraph{Why.} In formal\_13 (2\,000 episodes), the agent remained stuck in Phase 1 (Physics of
War), winning only 27.5\% of evaluation games against the heuristic and losing 37 to 55\% of its
revealed attacks. The hinge loss \(\max_{a \neq a_E}[Q(a) + 0.1] - Q(a_E)\) stopped improving at
0.085, unable to rank the expert's move above 30+ other legal choices when attention dropout jittered
the Q-values at every gradient step.

\paragraph{Result.} On 66\,576 transitions from 200 held-out expert games, five variants were
pre-trained for 5\,000 updates and evaluated greedily over 40 games against each opponent:

\begin{center}
\small
\begin{tabular}{lcccc}
\toprule
Variant & vs Random & vs Heuristic & vs Expert & Revealed lost \\
\midrule
Margin (dropout 0.1) & 27 / 9 / 4 (67.5\%) & 11 / 20 / 9 (27.5\%) & 9 / 27 / 4 (22.5\%) & 47\% \\
Margin (dropout 0.0) & 27 / 10 / 3 (67.5\%) & 13 / 17 / 10 (32.5\%) & 10 / 30 / 0 (25.0\%) & 37\% \\
BC ($\tau = 0.1$) & 38 / 2 / 0 (95.0\%) & 26 / 9 / 5 (65.0\%) & 0 / 2 / 38 (95.0\% D) & 14\% \\
\textbf{BC ($\tau = 0.25$)} & \textbf{37 / 2 / 1 (92.5\%)} & \textbf{29 / 9 / 2 (72.5\%)} & \textbf{0 / 1 / 39 (97.5\% D)} & \textbf{17\%} \\
BC ($\tau = 0.1$) + Margin & 37 / 2 / 1 (92.5\%) & 25 / 10 / 5 (62.5\%) & 2 / 1 / 37 (92.5\% D) & 18\% \\
\bottomrule
\end{tabular}
\end{center}
Behaviour cloning with $\tau = 0.25$ immediately clears the Phase 1 curriculum threshold (50\% vs
heuristic) straight out of cold start (72.5\% wins), draws 39 of 40 games against its expert
teacher (losing only once), and cuts revealed attack blunders to 17\% (5 blunders in 15 attacks
against the expert).

\paragraph{Lesson.} A hinge margin pushes down only the closest rival; in games with high branching
factors, cross-entropy over legal moves suppresses the entire distribution of bad moves at once.

% ------------------------------------------------------------------------------------------
\section{Ten-thousand-episode run: formal\_14}
\label{sec:formal14}

\paragraph{Change.} Launching a full 10\,000-episode run (\texttt{runs/formal\_14}) incorporating:
\begin{itemize}
  \item Cold start from \texttt{OutlookExpert} with behaviour cloning ($\tau = 0.25$, weight 1.0,
        5\,000 pre-training updates, 10\% batch share decaying to 0 by episode 5\,000).
  \item Spatial self-attention dropout set to 0.0 in the online network.
  \item Replay update interval reduced from 8 to 4 steps (providing $\sim$440\,000 gradient updates
        across 10k episodes without learner lag).
  \item Replay buffer capacity doubled to 150\,000 transitions ($\sim$430 games of history,
        consuming $\sim$1.8\,GB of the 6.44\,GB GPU VRAM).
  \item Learning rate set to $0.0001$, checkpointing every 500 episodes, and adding \texttt{OutlookExpert}
        to the 500-episode periodic evaluation suite.
\end{itemize}

\paragraph{Result.} Read at episode 5\,000 (393\,845 updates). The 40-game evaluations were flat over
the run: wins against the heuristic were 25, 23, 25, 21, 24, 26, 21, 20, 18 and 22 at episodes 500 to
5\,000, the level the cold start had set, and wins against the expert never exceeded 8.

The plotted loss fanned out from about update 280\,000 and fell off a cliff near update 375\,000. Both
come from the demonstrations. The batch held $\mathrm{round}(128 \times 0.10 \times (1 - e/5000))$ of
them, 12 falling to 0, and the cloning term is a mean over those rows at a fixed weight of 1, so its
noise grew as the rows ran out. The 99th percentile of the loss was 2.85 with five rows, then 3.11,
3.70, 5.03 and 8.25 with one row (maximum 19.1, one position against 127 in the C51 term). The formula
reaches zero rows at episode 4\,805; the largest loss in five episodes fell from 12.6 to 1.44 between
episodes 4\,805 and 4\,810 and stayed below 1.5. The C51 loss itself did not fail: its median kept
falling (1.22 before, 0.91 at the end).

Removing the term did collapse the policy. Comparing the checkpoints on either side of it:

\begin{center}
\small
\begin{tabular}{lcc}
\toprule
Measure & Episode 4\,500 & Episode 5\,000 \\
\midrule
Best minus second-best legal $Q$ (3\,000 own positions) & 0.126 & 0.006 \\
Best minus worst legal $Q$ (same positions) & 0.587 & 0.049 \\
Agreement with the expert's move (2\,500 positions) & 45.8\% & 14.6\% \\
Vs heuristic, 100 games (W/L/D) & 64 / 12 / 24 & 41 / 23 / 36 \\
Vs expert, 100 games (W/L/D) & 7 / 31 / 62 & 0 / 52 / 48 \\
Vs random, 100 games (W/L/D) & 86 / 6 / 8 & 82 / 6 / 12 \\
\bottomrule
\end{tabular}
\end{center}
The two checkpoints chose the same greedy move in 13\% of those positions. The NoisyNet scale was not
the cause: the mean $|\sigma|$ fell from 0.0187 to 0.0153 after episode 4\,800, in the value layers
(0.0206 to 0.0168 and 0.0157 to 0.0130) while the last advantage layer did not move (0.0284 to 0.0283).
Noise nevertheless changed the greedy move in 31\% of the expert positions at episode 5\,000, against
14\% at 4\,500, because the gaps it competes with had shrunk twentyfold. The share of training games won by flag
capture also fell, from 0.34 in episodes 4\,400 to 4\,600 to 0.14 in 4\,600 to 4\,800, before the last
demonstration left; that fall was not explained.

\paragraph{Later reading: the evaluations against the expert.} The wins hide a change the losses
show. In the 40-game evaluations against the expert the agent went 2 / 3 / 35 (wins / losses /
draws) at episode 500 and 0 / 4 / 36 at episode 1\,000: it drew with its teacher, as the cold start
alone had (Section~\ref{sec:behaviour_cloning}). From episode 1\,500 to 4\,500 it went 6 / 16 / 18,
6 / 19 / 15, 8 / 22 / 10, 2 / 27 / 11, 3 / 15 / 22, 7 / 12 / 21 and 3 / 19 / 18: 12 to 27 losses in
every evaluation, against 3 and 4 before. The demonstrations were still 9 of the 128 rows of a
batch at episode 1\,500, so this began some 3\,300 episodes before the last demonstration left.
Against the heuristic the same evaluations stayed level (18 to 26 wins). Two things changed between
episodes 1\,000 and 1\,500, and the run does not separate them: the demonstration share kept
falling, and the third stage added a 500-episode-old copy of the agent to the opponents. The
reinforcement learning therefore did not merely add nothing to the cloned policy. Against the
strongest fixed opponent it made the policy worse while the cloning term was still in the loss,
and the evaluation against the weaker heuristic did not show it.

\paragraph{Lesson.} A falling loss can mean the network has stopped distinguishing actions. Here the
differences between moves were held up by the cloning term and the C51 targets did not sustain them,
so a demonstration schedule that decays to zero removes the policy with it; a mean over the rows of a
shrinking subset also needs its weight scaled with the subset. Verify all architectural assumptions and cold-start transfer in fast, isolated
probes before committing GPU hours to an overnight multi-thousand episode run.

% ------------------------------------------------------------------------------------------
\section{A batched game engine on the GPU}
\label{sec:vec-engine}

\paragraph{Change.} A second implementation of the game (\texttt{vec/engine.py}) holds any number
of games as tensors on the GPU and steps them together. The board is stored in the frame of the
side to move and rotated by 180 degrees after every ply, so no rule depends on which player is
moving. The rules, the 41 observation planes and the 3\,600-action encoding are those of the
original environment. A test replays games of random and attacking moves on both implementations
and compares the legal moves, every observation plane, the ending and the winner on every ply; on
the GPU the same check ran over 32 games and 29\,955 plies. The outlook expert was ported the same
way: one call scores all 3\,600 actions of every game, and the scores equal the original's to
within $10^{-3}$ on more than 1\,000 moves.

\paragraph{Why.} formal\_14 made 393\,845 updates and 1\,575\,961 environment steps in about ten
hours: about 95\,ms per update and 42 steps per second. The learner was the cost, with three passes
of a 3\,600 $\times$ 51 output for a batch of 128, so faster environment code could not help (the
author reports that an earlier port of the environment to Rust gained about 5\%; this was not
re-measured). A different learner needs far more games per second, and then the game itself
becomes the limit.

\paragraph{Result.} The first batched version played about 27\,000 plies per second with the
observation, or 38\,ms per ply of 1\,024 games, and did not speed up from 1\,024 to 4\,096 games.
A profile by operation found three causes. A cumulative sum over a boolean tensor is promoted to
64-bit integers and was most of the cost of the legal-move mask; eight chained logical ANDs along
the nine squares of a line replaced it. Indexing with a three-dimensional index tensor became an
index-select on a flat index. Three general tensor contractions over one-hot planes became a table
lookup by rank and one sum. After these changes 1\,024 games cost 0.9\,ms for the legal-move mask,
5.0\,ms for the observation and 4.1\,ms for a step, and the outputs were unchanged (the tests
above). The network's forward pass, 10\,ms for 1\,024 positions, is the larger cost again.

One difference from the original was found and kept as a correction. When a Scout on one edge of
the board looks down a line that is empty to the far edge, the original outlook plane reads the
contents of square 0 instead of nothing (3 of the 29\,955 positions).

\paragraph{Lesson.} Vectorised code has to be profiled operation by operation: one promoted data
type cost more than the rest of the step together. A second implementation is only worth its speed
if a test holds it to the first one on every ply.

% ------------------------------------------------------------------------------------------
\section{The outlook expert is a passive defender}
\label{sec:passive-expert}

\paragraph{Observation.} With the batched engine a match of 1\,000 games takes about half a
minute. The outlook expert beat the random player 950 to 11 with 39 draws. Against itself it won 4,
lost 5 and drew 991, all 991 by the 200-ply no-battle rule, after 289 plies on average. Two random
players went 209 / 201 / 590.

\paragraph{Reading.} In the expert's score, attacking a hidden piece has a negative expected
value and a quiet move has no term that leads towards a battle, so two experts wait for each other.
This revises the reading of Section~\ref{sec:behaviour_cloning}, where the cloned network ``drew
39 of 40 games with its teacher'': that was not parity of skill but two players refusing to
attack. The expert beats opponents that attack it and has no way to win against one that does not.

\paragraph{Change.} Three optional terms were added to the batched expert, all zero by default so
that it still equals the original. \emph{Pursuit} draws a quiet move towards enemy pieces the moved
piece could attack with positive expected value, weighted by the expected gain and falling with
the square of the distance. \emph{Repel} pushes it away from pieces likely to beat it, beyond the
adjacent square. \emph{Probe} rewards attacking a hidden piece with a cheap piece.

\paragraph{Result.} Against the original expert, 600 games each (wins / losses / draws):
\begin{center}
\small
\begin{tabular}{lc}
\toprule
Terms & against the original expert \\
\midrule
pursuit 1.0 & 142 / 53 / 405 \\
pursuit 3.0 & 190 / 82 / 328 \\
probe 1.0 & 12 / 10 / 578 \\
probe 2.5 & 13 / 17 / 570 \\
pursuit 1.0, probe 1.0 & 168 / 42 / 390 \\
pursuit 2.0, probe 1.5, repel 0.5 & 221 / 62 / 317 \\
pursuit 3.0, probe 2.5, repel 1.0 & 193 / 65 / 342 \\
\bottomrule
\end{tabular}
\end{center}
Probing alone changed nothing, because the pieces never met. A hill climb over these weights and
the expert's own constants, in which a candidate replaces the champion when it wins a match of 800
games by a net four per cent, accepted few candidates in 93 trials (most single changes lie within
the noise of a match). Its last champion beat the original expert 348 to 96 with 356 draws and the
random player 385 to 8.

\paragraph{Lesson.} A one-ply evaluation needs a term whose value changes with distance to what
it wants, or it values captures it will never stand next to. And a draw between a student and its
teacher shows that they play alike, not that they play well.

% ------------------------------------------------------------------------------------------
\section{A policy-value network, cloned in minutes}
\label{sec:pv-clone}

\paragraph{Change.} The Q-network was replaced by a network with a policy of its own
(\texttt{vec/net.py}): a residual tower of six blocks and 64 channels without normalisation
layers, a policy head with one logit for each of the 3\,600 actions, and a value head with one
output in $[-1, 1]$ for the side to move. It is cloned from the expert on a stream of batched
games (\texttt{vec/clone.py}). On every ply of 2\,048 games the expert scores the legal moves and
the network is trained towards the softmax of the scores at temperature 0.03. Each side of a game
is driven by the expert (30\%), the expert with a quarter of its moves random (20\%), the random
player (10\%) or, after 300 updates, the network itself (40\%), so the network is also labelled on
the positions its own mistakes lead to \cite{DAgger}.

Three additions came with it. The observation grew from 41 to 59 planes, all public information:
the share of each enemy rank that is still hidden (12 planes), the own pieces that have moved and
that have made a long move, and the from- and to-squares of the last two moves. A captured piece
only vanishes from the 41 planes, so whether the enemy Marshal or Spy was still in the game could
not be read from them. A third head predicts the rank of the enemy piece on each square and is
trained on the hidden ones against the true board, which the trainer knows. This is the task the
thesis gave to AAREN, with a label for every hidden piece on every ply instead of only for the
pieces battles reveal. Finally the setups are drawn from a bank of 40\% thesis setups (Flag in a
back corner with Bombs beside it), 45\% structured ones (Flag anywhere in the back two rows with
two or three Bombs around it, the other pieces by row preferences) and 15\% random ones.

\paragraph{Why.} Section~\ref{sec:formal14}: in the Q-network the policy lived in differences
between action values of a few hundredths, held up by the cloning term and lost with it.

\paragraph{Result.} About 22\,500 labelled positions per second. After 70 seconds (447 updates)
the greedy network beat the random player in 188 of 200 batched games, with 4 losses; the expert
itself wins 95\%. After ten minutes (5\,588 updates, 59 planes, the mixed setups) it won 379 of
400 against random and drew 396 of 400 with the expert, chose the expert's best move in 80\% of
positions (equal-scored quiet moves keep this below 100\%), and named the rank of a hidden piece
correctly 39.5\% of the time, where always naming the Scout would score 20 to 24\%
(Section~\ref{sec:cold-start-level} gives the cold start of the earlier runs: 5\,000 updates on
about 35\,000 positions in some nine minutes).

\paragraph{Lesson.} When the teacher is a function that can be evaluated on any position, cloning
is a streaming supervised problem and costs minutes; it should be treated as a starting point to
be measured, not as a phase of the run.

% ------------------------------------------------------------------------------------------
\section{Self-play with a policy gradient improves on the clone}
\label{sec:ppo}

\paragraph{Change.} The cloned network was trained further by proximal policy optimisation
\cite{PPO} on batched self-play (\texttt{vec/ppo.py}). Each iteration plays 128 plies of 1\,024
games and makes two passes over the learner's decisions in minibatches of 4\,096 (learning rate
$3 \times 10^{-4}$, clip 0.2, entropy weight 0.005, value weight 0.5, belief weight 0.25). The only
reward is the result of the game: $+1$, $-1$, and $-0.5$ to both sides for a draw. Advantages are
generalised advantage estimates \cite{GAE} ($\gamma = 0.998$, $\lambda = 0.95$) over each side's
own sequence of decisions, so a result reaches a move through the returns that followed it and not
through a chain of target networks. Half of the games are self-play, a quarter are against the
original expert, 5\% against the random player and 20\% against frozen earlier copies of the
learner, one kept every 25 iterations. For the first four iterations only the value head is
trained.

\paragraph{Why.} The first direction named in Chapter~\ref{chap:recommendations}: a learner whose
policy does not depend on small differences between action values. A first version of the warm-up
trained the value loss through the shared trunk. In one iteration the policy moved by a KL
divergence of 0.26 and the greedy win rate against the random player fell from 94\% to 60\%, with
the policy loss switched off: the cloned policy head depends on features the value loss was free
to rewrite. The warm-up now steps the value head's own parameters only.

\paragraph{Result.} From the ten-minute clone, on a GPU shared with the expert tuning, in 300
batched games per line (wins / losses / draws):
\begin{center}
\small
\begin{tabular}{rrcccc}
\toprule
Iteration & decisions & greedy vs random & greedy vs expert & sampled vs expert & sampled vs its start \\
\midrule
0 (clone) & 0 & 379 / 18 / 3 of 400 & 2 / 2 / 396 of 400 & -- & -- \\
40  & 3.8 M  & 292 / 7 / 1 & 200 / 37 / 63 & 240 / 59 / 1 & 238 / 60 / 2 \\
80  & 7.7 M  & 298 / 2 / 0 & 261 / 19 / 20 & 256 / 40 / 4 & 269 / 29 / 2 \\
120 & 11.7 M & 295 / 5 / 0 & 241 / 16 / 43 & 278 / 18 / 4 & 281 / 16 / 3 \\
\bottomrule
\end{tabular}
\end{center}
Iteration 40 was reached after about 12 minutes. In the training games the self-play went from 30
decisive and 680 drawn games in iterations 1 to 5 to 1\,078 decisive and one drawn in iterations
36 to 40. The run stopped after 45 minutes at iteration 132. That network was then evaluated in
the original environment by the thesis's own routine, greedy, in 200 games against each opponent:
\begin{center}
\small
\begin{tabular}{lcc}
\toprule
Opponent & wins / losses / draws & revealed attacks lost \\
\midrule
random          & 199 / 0 / 1  & 14\% \\
heuristic       & 193 / 4 / 3  & 9\% \\
smart heuristic & 152 / 1 / 47 & 9\% \\
outlook expert  & 177 / 1 / 22 & 18\% \\
\bottomrule
\end{tabular}
\end{center}
The best earlier result against the heuristic was 29 / 9 / 2 in 40 games, straight after cloning
(Section~\ref{sec:behaviour_cloning}), and formal\_14 never won more than 8 of 40 against the
expert.

What the learner changed in its play, from 400 sampled games against the expert before and after
40 iterations:
\begin{center}
\small
\begin{tabular}{lcc}
\toprule
Measure & clone & iteration 40 \\
\midrule
Result (wins / losses / draws) & 7 / 5 / 388 & 309 / 84 / 7 \\
Attacks per game & 2.7 & 20.2 \\
Attacks on hidden pieces that have moved, won & 41\% of 0.9 per game & 76\% of 9.4 per game \\
Attacks on hidden unmoved pieces, won & 8\% of 1.5 per game & 52\% of 7.8 per game \\
Marshal lost, per game & 0.03 & 0.62 \\
\bottomrule
\end{tabular}
\end{center}
The first thing self-play found is that a hidden piece which has moved cannot be a Bomb and is
usually weak, so a strong piece wins three such attacks in four. At that point it was bold and not
yet careful: it lost its Marshal in most games, a third of those into a Bomb and a quarter to the
Spy.

\paragraph{Lesson.} Reinforcement learning improved on its starting point here for the first
time in this project, within minutes. Four things differ from the Q-learner, and this run does not
separate them: the policy is its own output, trained by a policy gradient; the game's result
reaches a move through measured returns; twelve minutes hold more experience than formal\_14's ten
hours (3.8 against 1.6 million decisions); and a draw costs both sides while half the games are
self-play, so mutual stalling is the worst joint outcome. A plateau of a value-based learner at a
few dozen updates per second said little about what reinforcement learning could do in this game.

% ------------------------------------------------------------------------------------------
\section{Three-hour PPO self-play and post-run audits (ppo\_2)}
\label{sec:ppo2}

\paragraph{Change.} The \texttt{ppo\_1} run of Section~\ref{sec:ppo} was resumed at iteration 132
(12\,940\,127 decisions, six league snapshots) for three hours (\texttt{runs/ppo\_2}) to iteration
771 (75.48\,M learner decisions recorded at iteration 770). Rollout size (1\,024 parallel games,
128 plies), PPO hyperparameters (constant learning rate $3 \times 10^{-4}$, two epochs, minibatch
4\,096, clip 0.2, entropy weight 0.005, value weight 0.5, belief weight 0.25, draw reward $-0.5$),
and the diverse setup bank stayed as in \texttt{ppo\_1}. The opponent mix was broadened to 45\%
self-play, 20\% scripted experts drawn uniformly from three styles---the original passive
\texttt{expert}, the hand-set proactive \texttt{expert2} (pursuit 2.0, repel 0.5, probe 1.5), and
the hill-climbed \texttt{expert3} of Section~\ref{sec:passive-expert} (pursuit 1.721, repel 0.5369,
probe 0.6253, attack chance 0.3724, clock 0.4912, advance 0.0331, pull 0.7336, power 1.7568)---3\%
random play, and 32\% frozen snapshots in a 12-slot league (one added every 40 iterations, sampled
uniformly and evicted at random when full).

\paragraph{Why.} Section~\ref{sec:ppo} showed 45 minutes of PPO self-play lifting the cloned
network above the passive expert, while Section~\ref{sec:passive-expert} produced two proactive
scripted opponents that attack instead of waiting. Continuing the run against all three expert
styles and a deeper snapshot league tested how far pure PPO self-play would climb before hitting
new limits.

\paragraph{Result.} In the periodic 300-game batched evaluations on the diverse setup bank
(wins / losses / draws):
\begin{center}
\small
\begin{tabular}{rrccccc}
\toprule
Iter & decisions & greedy vs random & greedy vs expert & sample vs expert & sample vs expert2 & sample vs expert3 \\
\midrule
200 & 19.56\,M & 298 / 2 / 0 & 236 / 10 / 54 & 292 / 6 / 2  & 217 / 66 / 17 & 220 / 64 / 16 \\
300 & 29.40\,M & 298 / 2 / 0 & 208 / 10 / 82 & 277 / 18 / 5 & 228 / 56 / 16 & 233 / 51 / 16 \\
400 & 39.20\,M & 294 / 6 / 0 & 234 / 10 / 56 & 278 / 16 / 6 & 231 / 46 / 23 & 224 / 57 / 19 \\
500 & 49.15\,M & 300 / 0 / 0 & 246 / 7 / 47  & 289 / 10 / 1 & 234 / 53 / 13 & 246 / 47 / 7  \\
600 & 58.90\,M & 292 / 7 / 1 & 239 / 6 / 55  & 280 / 17 / 3 & 233 / 41 / 26 & 247 / 42 / 11 \\
700 & 68.58\,M & 295 / 5 / 0 & 241 / 14 / 45 & 278 / 20 / 2 & 245 / 42 / 13 & 238 / 49 / 13 \\
\bottomrule
\end{tabular}
\end{center}
Win rates rose against both proactive experts (from 72.3\% to 81.7\% against \texttt{expert2} and
from 73.3\% to 82.3\% at iteration 600, 79.3\% at 700, against \texttt{expert3}). In a 600-game
sampled head-to-head match on diverse setups, the final \texttt{ppo\_2} checkpoint beat
\texttt{ppo\_1} 481 / 102 / 17 (80.2\% wins). Evaluated in the original thesis environment across
400 games (100 against each of random, heuristic, smart heuristic, and outlook expert),
\texttt{ppo\_2} scored 393 / 0 / 7 (98.25\% wins, zero losses) greedily and 397 / 2 / 1 (99.25\%
wins) when sampled, including 100 / 0 / 0 against the outlook expert.

\paragraph{Diagnosis: three limits exposed by auditing ppo\_2.} Although \texttt{ppo\_2} nearly
saturated the fixed thesis opponents, inspecting its logs, draws, and belief head uncovered three
structural weaknesses:
\begin{enumerate}
  \item \emph{Greedy oscillation draws.} Greedy play against the passive \texttt{expert} still
        ended in 45 to 82 no-battle draws per 300 games (15 to 27\%), whereas sampled play drew
        only 1 to 6. Inspecting those greedy draws showed that 59.5\% were pure two-square
        $A \leftrightarrow B$ oscillations and 83.8\% were cycles of at most four moves. None of
        the game engines enforced Stratego's official two-square repetition rule, so a deterministic
        policy could shuffle indefinitely instead of choosing a winning alternative.
  \item \emph{Optimisation churn and non-transitive cycling.} Keeping the learning rate fixed at
        $3 \times 10^{-4}$ left the PPO clip fraction at 0.14 to 0.17 and per-update KL divergence
        at 0.0131 to 0.0190 all the way to iteration 770. Combined with uniform snapshot sampling
        and random eviction and no regularisation towards a slower anchor policy, the learner
        cycled rather than settling: sampled wins against \texttt{expert} moved
        $97.3\% \to 92.3\% \to 96.3\% \to 92.7\%$, and policy entropy swung from 1.216 (iteration
        140) down to 0.838 (iteration 475), back up to 1.055 (iteration 650), and to 0.900--0.964
        at iterations 700--770. Value explained variance stayed steady between 0.75 and 0.84.
  \item \emph{An aggregate belief metric confounded by game length and setup priors.} Logged
        belief-head accuracy floated between 0.360 and 0.410 and appeared to fall during aggressive
        stretches (0.410 at iteration 445 when games averaged 322 plies, versus 0.360 at iteration
        650 when games averaged 231 plies). Stratifying by ply and setup bank showed that per-ply
        accuracy did improve over training---shorter games simply concentrated more decisions in
        the opening, where fewer enemy pieces have moved or been captured---yet accuracy on
        uniformly random setups at plies 0 to 40 remained below 20\%. The aggregate 38 to 40\%
        figure therefore reflects setup priors (heuristic and structured placements) and late-game
        elimination far more than early behavioural inference.
\end{enumerate}

\paragraph{Lesson.} Once fixed opponents approach saturation, aggregate win rates stop diagnosing
the learner. Direct matches against frozen checkpoints, cycle audits of greedy draws, and
ply-stratified auxiliary metrics expose cycling, missing rule constraints, and composition shifts
that headline averages hide.

% ------------------------------------------------------------------------------------------
\section{Codebase audit fixes across the Rainbow learner, setup generator, and GUI}
\label{sec:audit-fixes}

\paragraph{Change.} A full audit of the training pipeline, setup generator, and desktop GUI
uncovered six silent bugs; all six were fixed and locked down with unit tests:
\begin{itemize}
  \item \emph{Horizontal-flip corruption of directional battle-outlook planes}
        (\texttt{network/drqn\_agent.py}). Mirror augmentation previously reflected columns
        $c \mapsto 9 - c$ with \texttt{torch.flip} without swapping the directional battle-outlook
        planes $34 \leftrightarrow 35$ ($P(\text{win})$ left/right) and $38 \leftrightarrow 39$
        ($P(\text{lose})$ left/right), corrupting half of all demonstration and replay transitions
        from \texttt{formal\_9} onwards. A helper (\texttt{\_flip\_state\_horizontal}) now reflects
        columns and swaps both plane pairs whenever at least 41 planes are present, applied to
        both \texttt{state} and \texttt{next\_state} in demonstration recording and replay storage.
  \item \emph{PBRS breaking the ten-fold PER win boost} (\texttt{network/drqn\_agent.py}).
        Winning transitions were previously detected by \texttt{reward\_clipped >= 0.8 and done}.
        Under potential-based shaping \cite{NgPBRS}, a terminal transition subtracts
        $\Phi(s_{\text{prev}})$, so dominant wins with $\Phi(s_{\text{prev}}) > 0.2$ fell below
        0.8 and missed the ten-fold priority boost. Win detection now checks the unshaped game
        outcome (\texttt{winner == self.player\_id}) whenever outcome metadata is passed, falling
        back to the clipped-reward threshold only when no unshaped winner is available.
  \item \emph{C51 atom entropy instead of action-policy regularisation}
        (\texttt{network/drqn\_agent.py}). The optional entropy-regularisation, KL-regularisation,
        and dynamic-damping terms previously computed entropy and KL over the 51 C51 return atoms
        \cite{C51} (\texttt{dim=2}), which would push return distributions toward uniform rather
        than regularising the policy. They now compute entropy and KL on the softmax policy over
        expected Q-values (masked to legal moves when available); defaults remain disabled.
  \item \emph{Empty full-length Scout ray in the original outlook code}
        (\texttt{environment/battle\_outlook.py}). When a Scout on an edge looked down an empty
        nine-square ray, the original outlook code read square $(0, 0)$ instead of sentinel index
        100 (the three-position discrepancy noted in Section~\ref{sec:vec-engine}). Guarding empty
        rays with \texttt{torch.where} brings it into exact agreement with \texttt{vec/engine.py},
        and the cross-engine parity test now requires zero tolerated plane differences across all
        59 planes.
  \item \emph{Off-by-one piece groups in the heuristic setup}
        (\texttt{environment/heuristic\_setup.py}). Hardcoded integer rank lists misclassified
        several ranks (\texttt{[7, 8]} omitted Colonel $= 9$; \texttt{[5, 6]} shifted Captain and
        Lieutenant; \texttt{[4]} duplicated Miner $= 4$ as Sergeant). Grouping now uses
        \texttt{PieceType} enum members directly and places the exact 40-piece army without
        duplication or omission.
  \item \emph{Pygame GUI Expectamax utility, chance-node pruning, and policy-value support}
        (\texttt{Python Stratego Game/dqn\_bot\_logic.py}). The Expectamax loss utility
        $-0.5 - ((12 - r)/12) \cdot 0.5$ was inverted---penalising the loss of a Scout more than a
        Marshal---and was corrected to $-0.5 - (r/12) \cdot 0.5$. The \texttt{top\_n = 3}
        chance-node truncation (which discarded 52.5\% of probability mass under the default prior)
        was removed so all ranks with positive probability (at most 12) are evaluated. In
        addition, \texttt{DQNBotLogic} now auto-detects \texttt{PolicyValueNet} checkpoints
        (59-plane encoding, legal-masked acting, value evaluation, multi-move recommendations, and
        belief-head rank probabilities) alongside existing Rainbow checkpoints.
\end{itemize}

\paragraph{Why.} Auditing \texttt{formal\_14} verified that the horizontal-flip defect had a direct
behavioural cost: because vertical outlook planes ($32, 33, 36, 37$) are invariant to horizontal
mirroring while horizontal planes ($34, 35, 38, 39$) were inverted in 50\% of stored transitions,
\texttt{formal\_14} committed sure-loss blunders on horizontal attacks at 5.7 times the vertical
rate (2.32\% versus 0.41\%) and lost 33.1\% of horizontal attacks on revealed pieces versus 18.3\%
vertically. The remaining five fixes eliminated silent failures in priority boosting, optional
policy regularisers, setup generation, engine parity, and GUI search.

\paragraph{Result.} Unit tests verify exact $10^{-6}$ agreement between
\texttt{\_flip\_state\_horizontal} and freshly mirrored encoding across all 41 planes for both
players, unshaped-winner priority boosting under large positive $\Phi$, action-policy entropy and
dynamic damping, zero-tolerance 59-plane parity including all four edge Scout rays, exact 40-piece
heuristic setups across 25 seeds for both sides, Expectamax loss ordering (Marshal $>$ Scout) with
full 12-rank chance nodes, and bit-for-bit 59-plane parity for \texttt{PolicyValueNet} in the GUI
adapter (231 passed, 0 failed across the full test suite, with one pre-existing documented xfail).
Re-running \texttt{formal\_14} with the fixed Rainbow Q-learner was not measured, as active
training has moved to the batched policy-gradient trainer; these fixes ensure the thesis codebase,
setup generator, and desktop GUI are sound.

\paragraph{Lesson.} Spatial data augmentation must transform every directional feature plane
alongside the board axes and action indices, not merely flip coordinates. Comparing directional
error rates (horizontal versus vertical blunders) provides a fast behavioural test for mirror
augmentation bugs, and named enum members should always replace raw rank integers in game logic.

% ------------------------------------------------------------------------------------------
\section{Enforcing the official two-square repetition rule across all four engines}
\label{sec:two-square}

\paragraph{Change.} Stratego's two-square repetition rule (in the box-rule form; the ISF tournament
rule allows one more move) was added identically across all
four game engines (\texttt{vec/engine.py}, \texttt{environment/environment.py},
\texttt{web/play.html}, and \texttt{Python Stratego Game/stratego.py}): if a player's previous two
consecutive quiet moves reversed each other ($A \to B$, then $B \to A$), repeating $A \to B$ on
that player's next turn is illegal and masked out of the legal-move set. Each side tracks its own
last two quiet moves in its own coordinate frame; any battle or game reset clears both trackers,
and any non-reversing move breaks the pattern. If the forbidden move is a player's only remaining
legal move, that player loses by immobility (\texttt{no\_moves}), matching the official rules. The
observation stays at 59 planes, keeping full checkpoint compatibility with \texttt{ppo\_2}. In the
same pass, \texttt{\_is\_position\_repetitive()} in \texttt{environment/environment.py} was fixed
to compare \texttt{recent\_moves[i]} with \texttt{recent\_moves[i - 2]} (the same player's
previous move) rather than \texttt{i - 1} (the opponent's move).

\paragraph{Why.} The \texttt{ppo\_2} draw audit (Section~\ref{sec:ppo2}): 59.5\% of greedy
no-battle draws against the passive expert were pure two-square $A \leftrightarrow B$ cycles and
83.8\% were cycles of at most four moves. Without the rule in the action mask, a deterministic
greedy policy could cycle indefinitely into a 200-ply no-battle draw whenever a local back-and-forth
move outscored attacking moves by a small margin. Implementing the rule in only a subset of engines
would break cross-engine parity: adding it to the training environment immediately caused the GUI
parity test to fail at ply 35 until \texttt{stratego.py} received the matching tracker.

\paragraph{Result.} All four engines pass exact cross-engine replay parity tests
(\texttt{test\_vec\_engine.py}, \texttt{test\_web\_engine.py}, and \texttt{test\_gui\_engine.py}),
including unit tests for non-reversing resets, battle resets, trapped \texttt{no\_moves} losses,
and Scout long-move masking without breaking slide continuation. To isolate the effect of the rule
alone---without changing a single weight---the trained \texttt{ppo\_2} checkpoint
(\texttt{runs/ppo\_2/net\_final.pt}) was re-evaluated on the upgraded batched engine across 300
games per opponent on the diverse setup bank (seed 123):
\begin{center}
\small
\begin{tabular}{lccc}
\toprule
Matchup (300 games, \texttt{ppo\_2} with two-square rule) & W / L / D & win rate & mean plies \\
\midrule
greedy vs \texttt{expert}  & 284 / 11 / 5  & 94.7\% & 183 \\
greedy vs \texttt{expert2} & 237 / 40 / 23 & 79.0\% & 218 \\
greedy vs \texttt{expert3} & 242 / 40 / 18 & 80.7\% & 208 \\
sampled vs \texttt{expert}  & 285 / 13 / 2  & 95.0\% & 197 \\
sampled vs \texttt{expert2} & 234 / 62 / 4  & 78.0\% & 205 \\
sampled vs \texttt{expert3} & 235 / 62 / 3  & 78.3\% & 208 \\
\bottomrule
\end{tabular}
\end{center}
Enforcing the two-square rule alone cut greedy no-battle draws against the passive \texttt{expert}
from 45--82 per 300 games (Section~\ref{sec:ppo2}, 69.3--82.0\% wins at 215--257 plies) down to
5 of 300 (94.7\% wins at 183 plies), closing the entire gap with sampled play (95.0\% wins, 2
draws). Against both proactive experts, greedy play under the rule also matched or slightly
outperformed sampled play in win rate while cutting losses from 62 to 40.

\paragraph{Lesson.} Missing game-rule constraints show up most sharply under greedy evaluation,
where stochastic tie-breaking no longer breaks illegal cycles. Enforcing a rule inside the shared
action mask across every engine---and locking all implementations together with replay parity
tests---both fixes evaluation artefacts immediately and prevents silent drift between training,
browser play, and desktop GUI.

% ------------------------------------------------------------------------------------------
\section{Stabilising PPO self-play: unbiased bootstrap, cosine annealing, PFSP, and magnet KL (ppo\_3)}
\label{sec:ppo3}

\paragraph{Change.} Four algorithmic upgrades were made to the batched PPO trainer
(\texttt{vec/ppo.py}), keeping the 59-plane, 64-channel, six-block \texttt{PolicyValueNet}
architecture and checkpoint format backward-compatible with \texttt{ppo\_2}, together with a
fixed-reference evaluation flag (\texttt{--eval-ref}):
\begin{itemize}
  \item \emph{Unbiased one-ply lookahead \texttt{side\_gae()} bootstrap.} At rollout truncation
        ($t = T$), \texttt{side\_gae()} previously set the waiting side's bootstrap value to
        $-v_{\text{last}}$, negating the moving side's value estimate. Because a draw pays $-0.5$
        to both sides, $V_{\text{mover}} + V_{\text{other}} = -P(\text{draw}) \neq 0$, and against
        asymmetric opponents (\texttt{random}, scripted experts, or weaker snapshots) negating the
        network's valuation of the opponent's observation misestimates the learner's value. The
        rollout now steps one extra lookahead ply at $s_T$: the side to move at $s_T$ bootstraps
        from $V(s_T)$ when it is the learner, and after stepping ply $T$ the waiting side receives
        either its exact terminal reward (\texttt{draw} or $-z_T$) if ply $T$ ended the game, or
        $V(s_{T+1})$ evaluated from its own canonical view. When lookahead tensors are omitted,
        \texttt{has\_boot} zeroes only the unbootstrapped final truncated decision without negating
        $v_{\text{last}}$.
  \item \emph{Cosine learning-rate annealing.} \texttt{cosine\_lr()} anneals the step size smoothly
        from \texttt{--lr} to \texttt{--lr-end} over the run duration (\texttt{--hours}) or
        iteration cap (\texttt{--max-iters}).
  \item \emph{Prioritized Fictitious Self-Play (PFSP) league.} \texttt{League} now tracks
        exponentially windowed per-slot learner W/L/D counts (window 200, draws counted as $0.5$)
        and samples snapshot opponents proportionally to $(1 - \hat{w}_j)^p$
        (\texttt{--pfsp-power}, default $2.0$) with Laplace smoothing ($\text{prior} = 1.0$) so
        new snapshots are explored. The starting anchor (\texttt{nets[0]}) is protected; when full,
        the league evicts the easiest (highest learner win-rate) non-anchor snapshot instead of
        evicting at random.
  \item \emph{Legal-masked magnet KL regularisation.} Inspired by Regularised Nash Dynamics
        \cite{perolat2022stratego} and Magnetic Mirror Descent \cite{Ataraxos2025}, the policy loss
        adds $\beta_{\text{KL}}\,\mathrm{KL}(\pi_\theta \,\|\, \pi_{\text{magnet}})$ computed
        strictly over legal actions (\texttt{masked\_kl\_loss}, masking illegal logits to $-10^9$
        in fp32 and zeroing illegal terms so both value and gradients are invariant to illegal
        logits) against a frozen magnet copy refreshed every \texttt{--magnet-every} iterations
        (default 40; skipped during value warm-up or when \texttt{--kl-coef 0}).
  \item \emph{Fixed reference evaluation (\texttt{--eval-ref}).} Periodic evaluations play sampled
        games ($T = 1.0$) against a fixed reference checkpoint to track head-to-head progress once
        scripted experts saturate.
\end{itemize}

\paragraph{Why.} Section~\ref{sec:ppo2}: in \texttt{ppo\_2} a constant learning rate of
$3 \times 10^{-4}$ left the PPO clip fraction at 0.14--0.17 and per-update KL at 0.0131--0.0190
through iteration 770; uniform snapshot sampling and random eviction without a magnet anchor
allowed non-transitive cycling in both policy entropy and win rate against the expert; and
negating $v_{\text{last}}$ biased truncated rollouts whenever draw probabilities were non-zero or
the opponent was not a symmetric self-play copy.

\paragraph{Result.} Initialised from \texttt{runs/ppo\_2/net\_final.pt} on the upgraded engine with
the two-square rule, \texttt{runs/ppo\_3} trained for 0.38 hours (24 minutes, stopping cleanly at
iteration 64 with 6\,253\,818 learner decisions and four league snapshots) using 1\,024 parallel
games, 128-ply rollouts, two epochs, minibatch 4\,096, cosine learning rate
$1.5 \times 10^{-4} \to 2 \times 10^{-5}$, entropy weight 0.01, magnet KL weight
$\beta_{\text{KL}} = 0.02$ refreshed every 40 iterations, PFSP power $p = 2.0$ (12-slot league,
snapshot every 20 iterations), a two-iteration value-head warm-up, and opponent mix 45\% self-play,
20\% three-style experts, 2\% random, and 33\% PFSP snapshots:
\begin{center}
\small
\begin{tabular}{rccccccccr}
\toprule
Iter & samples & lr & kl & mkl & clip & entropy & ev & belief & nobattle / 5 it \\
\midrule
1  & 0.10\,M & $1.50 \times 10^{-4}$ & 0.0000 & 0.0000 & 0.00 & 1.062 & 0.73 & 0.283 & 0  \\
10 & 0.98\,M & $1.46 \times 10^{-4}$ & 0.0078 & 0.0549 & 0.09 & 0.975 & 0.80 & 0.375 & 14 \\
30 & 2.93\,M & $1.13 \times 10^{-4}$ & 0.0063 & 0.1270 & 0.07 & 0.951 & 0.84 & 0.381 & 17 \\
40 & 3.91\,M & $7.02 \times 10^{-5}$ & 0.0040 & 0.1351 & 0.04 & 0.928 & 0.83 & 0.377 & 19 \\
45 & 4.40\,M & $5.78 \times 10^{-5}$ & 0.0028 & 0.0165 & 0.03 & 0.922 & 0.82 & 0.375 & 20 \\
60 & 5.86\,M & $2.72 \times 10^{-5}$ & 0.0016 & 0.0416 & 0.01 & 0.881 & 0.82 & 0.375 & 24 \\
\bottomrule
\end{tabular}
\end{center}
As the learning rate annealed, per-update KL fell monotonically after warm-up from 0.0078 to
0.0016 and the clip fraction dropped from 0.09 to 0.01 (against 0.013--0.019 KL and 0.14--0.17
clip throughout \texttt{ppo\_2}). Magnet KL (\texttt{mkl}) rose smoothly to 0.1351 at iteration
40, reset to 0.0165 when the magnet refreshed at iteration 40, and reached 0.0416 by iteration 60.
Policy entropy declined smoothly from 1.062 to 0.881 without the oscillations of \texttt{ppo\_2}.
Under the two-square rule, training no-battle draws per five-iteration block fell to 2--24
(\texttt{ppo\_2}: 17--94), while immobilisation (\texttt{nomove}) finishes rose to 276--316 after
warm-up (\texttt{ppo\_2}: 118--247). In the in-run 300-game evaluations, sampled score against the
fixed \texttt{--eval-ref} (\texttt{ppo\_2}) improved from 148 / 147 / 5 (49.3\% wins) at iteration
30 to 176 / 120 / 4 (58.7\% wins) at iteration 60 (and 178 / 122 / 0 = 59.3\% against
\texttt{start}), while greedy play against \texttt{expert} went 289 / 8 / 3 (96.3\%) at iteration
30 and 288 / 9 / 3 (96.0\%) at iteration 60, and sampled play at iteration 60 reached
248 / 49 / 3 (82.7\%) against \texttt{expert2} and 242 / 50 / 8 (80.7\%) against \texttt{expert3}.

\paragraph{Head-to-head and original-environment evaluations.} Across 600-game head-to-head matches
on the diverse setup bank (10\,000 setups, seed 123):
\begin{center}
\small
\begin{tabular}{lcccc}
\toprule
Matchup (600 games, diverse setups) & temperature & W / L / D & win rate & score \\
\midrule
\texttt{ppo\_3} final (iter 64) vs \texttt{ppo\_2} final & $T = 1.0$ & 335 / 262 / 3  & 55.8\% & 56.1\% \\
\texttt{ppo\_3} final (iter 64) vs \texttt{ppo\_2} final & $T = 0.0$ & 311 / 265 / 24 & 51.8\% & 53.8\% \\
\texttt{ppo\_3} \texttt{net\_60} vs \texttt{ppo\_2} final & $T = 1.0$ & 307 / 288 / 5  & 51.2\% & 51.6\% \\
\texttt{ppo\_3} \texttt{net\_60} vs \texttt{ppo\_2} final & $T = 0.0$ & 315 / 256 / 29 & 52.5\% & 54.9\% \\
\texttt{ppo\_3} \texttt{net\_30} vs \texttt{ppo\_2} final & $T = 1.0$ & 316 / 279 / 5  & 52.7\% & 53.1\% \\
\texttt{ppo\_3} final (iter 64) vs \texttt{ppo\_1} final & $T = 1.0$ & 493 / 102 / 5  & 82.2\% & 82.6\% \\
\bottomrule
\end{tabular}
\end{center}
\texttt{ppo\_3/net\_final.pt} beat \texttt{ppo\_2/net\_final.pt} in both sampled play
(335 / 262 / 3, 55.8\% wins, $z \approx 3.0$ on decisive games) and greedy play (311 / 265 / 24,
51.8\% wins, 53.8\% score), and beat \texttt{ppo\_1} 493 / 102 / 5 (82.2\%, up from
\texttt{ppo\_2}'s 481 / 102 / 17 = 80.2\%). No intermediate checkpoint outperformed
\texttt{net\_final.pt} overall.

In the original thesis environment (\texttt{vec/eval\_original.py}, 100 games per opponent across
random, heuristic, smart heuristic, and outlook expert):
\begin{center}
\small
\begin{tabular}{lcccc}
\toprule
Opponent (100 games each) & greedy W / L / D & greedy rev.\ lost & sampled W / L / D & sampled rev.\ lost \\
\midrule
random          & 100 / 0 / 0 & 23\% (66 / 293) & 100 / 0 / 0 & 26\% (75 / 287) \\
heuristic       & 100 / 0 / 0 & 20\% (52 / 256) & 100 / 0 / 0 & 18\% (49 / 272) \\
smart heuristic & 98 / 0 / 2  & 15\% (52 / 340) & 100 / 0 / 0 & 17\% (57 / 337) \\
outlook expert  & 99 / 0 / 1  & 31\% (75 / 239) & 100 / 0 / 0 & 29\% (73 / 256) \\
\midrule
Total (400 games) & 397 / 0 / 3 (99.25\%) & 22\% (245 / 1128) & 400 / 0 / 0 (100.0\%) & 22\% (254 / 1152) \\
\bottomrule
\end{tabular}
\end{center}
\texttt{ppo\_3} won 397 / 0 / 3 (99.25\%, zero losses) greedily and 400 / 0 / 0 (100.0\%) when
sampled, improving on \texttt{ppo\_2}'s 393 / 0 / 7 greedy and 397 / 2 / 1 sampled records.

\paragraph{Behavioural diagnosis and stratified belief audit.} Comparing \texttt{ppo\_3} and
\texttt{ppo\_2} across 600 sampled games each in \texttt{vec/diagnose.py} (seed 0) shows how
\texttt{ppo\_3} gained its edge and what it did not change. Against \texttt{expert}, \texttt{ppo\_3}
went 577 / 20 / 3 versus \texttt{ppo\_2}'s 570 / 30 / 0 (both 16.1 attacks per game; hidden-moved
attacks won 84\% of 7.3 per game versus 82\% of 7.2; hidden-unmoved won 58\% and lost 34\% of 6.5
versus 59\% and 33\% of 6.5; revealed attacks won 71\% of 2.3 versus 73\% of 2.4; sure losses
0.57 versus 0.53 per game, 0.10 by non-Scouts for both). Against the proactive \texttt{expert3},
\texttt{ppo\_3} went 505 / 85 / 10 versus 484 / 114 / 2, cutting losses by 25\% (16.8 versus 16.3
attacks per game; hidden-moved won 79\% for both; hidden-unmoved won 57\% and lost 33\% for both;
non-Scout sure losses 0.09 versus 0.12 per game; Miners lost 2.37 versus 2.55 per game; Marshal
lost 0.51 versus 0.56 per game). However, Bomb deaths of the highest pieces did not fall: against
\texttt{expert}, the Marshal was lost 0.56 times per game (35\% into a Bomb, 33\% to the Spy)
versus 0.57 in \texttt{ppo\_2} (32\% Bomb, 39\% Spy) and the General 0.47 times per game in both
(37\% versus 33\% Bomb); against \texttt{expert3}, the Marshal Bomb share was 31\% for both and
the General was lost 0.57 times per game in both (35\% versus 31\% Bomb).

In the stratified belief-head check (200 sampled games per setup bank between \texttt{ppo\_3} and
\texttt{ppo\_2}, scoring argmax accuracy on all hidden enemy squares of identical observations
against the public rank-count argmax baseline):
\begin{center}
\small
\begin{tabular}{lccc}
\toprule
Setup bank & ply 0 (\texttt{ppo\_3} / \texttt{ppo\_2} / public) & ply 40 (\texttt{ppo\_3} / \texttt{ppo\_2} / public) & ply 100 (\texttt{ppo\_3} / \texttt{ppo\_2} / public) \\
\midrule
heuristic  & 33.8\% / 33.8\% / 20.0\% & 34.8\% / 31.0\% / 20.5\% & 40.4\% / 36.1\% / 27.1\% \\
structured & 26.6\% / 26.6\% / 20.0\% & 26.6\% / 25.7\% / 20.1\% & 30.1\% / 30.0\% / 24.9\% \\
random     & 17.9\% / 17.9\% / 20.0\% & 20.9\% / 19.7\% / 19.8\% & 26.0\% / 25.7\% / 23.3\% \\
\bottomrule
\end{tabular}
\end{center}
\texttt{ppo\_3} improved mid- and late-game belief accuracy over \texttt{ppo\_2} on heuristic
setups ($+3.8$ percentage points at ply 40 and $+4.3$ at ply 100) and edged ahead at ply 40 on
structured ($+0.9$) and random ($+1.2$) setups, while confirming Section~\ref{sec:ppo2}'s
finding: on uniformly random setups at plies 0--40, both networks remain near the 20\% public-count
baseline because early-game belief accuracy comes mainly from setup priors rather than move-pattern
inference.

As the verified winner across head-to-head and original-environment benchmarks,
\texttt{runs/ppo\_3/net\_final.pt} was exported to \texttt{Modular Stratego/web/model.onnx} and
\texttt{web/model.json} (maximum policy logit difference $9.2 \times 10^{-5}$ and value difference
$1.4 \times 10^{-6}$ between ONNX and PyTorch; previous \texttt{ppo\_1} web model backed up in
\texttt{runs/ppo\_3/web\_model\_prev/}), passing \texttt{test\_web\_engine.py} (2 passed in
26.33\,s).

\paragraph{Re-check on fresh seeds.} The head-to-head was replayed on two new setup banks
(\texttt{diverse\_setup\_bank(10000, seed=7)} and \texttt{seed=8}, 1\,000 games each). Sampled,
\texttt{ppo\_3} beat \texttt{ppo\_2} 553 / 441 / 6 and 581 / 412 / 7 (55.3\% and 58.1\% wins);
greedy, 527 / 428 / 45 and 535 / 424 / 41 (54.9\% and 55.5\% score). The edge is real.
\texttt{net\_60} scored 585 / 409 / 6 on seed 7, so the final checkpoint is not better than
\texttt{net\_60}: the two are within noise, and the 51.2\% of \texttt{net\_60} above was one
unlucky bank. Two cautions on what the gain means. First, \texttt{ppo\_3} trained under the
two-square rule and \texttt{ppo\_2} did not, and no run changed one thing at a time, so the 55--58\%
cannot be split between adapting to the rule and the four trainer changes. Second, the rule as
implemented forbids the third move between two squares, which is the classic box-rule wording;
the ISF tournament rule allows three moves and forbids the fourth.

\paragraph{Lesson.} Once self-play has found the basic tactical patterns of a game, further gains
come from damping optimisation churn rather than taking larger steps: cosine learning-rate
annealing, an unbiased lookahead bootstrap, PFSP opponent prioritisation, and legal-masked magnet
KL, together with training under the two-square rule, turned \texttt{ppo\_2}'s cycling into a
55--58\% head-to-head advance in a tenth of the total training budget (6.25\,M additional
decisions); which of these did the work is not yet separated. At the same time, component diagnostics show the
trade-off honestly: stabilising model-free PPO tightened Miner conservation and reduced losses
against proactive opponents, but did not reduce high-rank Bomb deaths or produce above-baseline
rank inference on uniformly random setups before ply 40.

% ------------------------------------------------------------------------------------------
\section{Two more hours at a falling learning rate (ppo\_4)}
\label{sec:ppo4}

\paragraph{Change.} \texttt{runs/ppo\_4} continued \texttt{ppo\_3/net\_final.pt} for two hours
with the \texttt{ppo\_3} settings unchanged (cosine learning rate $1.5\times10^{-4} \to
2\times10^{-5}$, magnet KL 0.02 refreshed every 40 iterations, PFSP power 2, a 12-slot league with a
snapshot every 25 iterations, \texttt{--eval-ref} \texttt{ppo\_3}, seed 4). It stopped at iteration
350 after 34 million learner decisions. Training charts for all runs are drawn by
\texttt{python -m vec.plot\_runs} into \texttt{runs/plots/ppo\_training.png}.

\paragraph{Why.} \texttt{ppo\_3} was still rising against its reference when its 24 minutes ran
out (49\% at iteration 30, 59\% at 60). The question was whether a longer schedule keeps the gain
going or whether the network has reached what it can learn from its input.

\paragraph{Result.} For the first 100 iterations \texttt{ppo\_4} sat at 49--52\% against
\texttt{ppo\_3} (300 sampled games per evaluation), which looked like a plateau. It then rose
steadily as the learning rate fell: 55\% at iteration 150, 58\% at 200, 60\% from 250 on. On two
fresh setup banks (1\,000 games each) the final net beat \texttt{ppo\_3}:
\begin{center}
\small
\begin{tabular}{lcc}
\toprule
\texttt{ppo\_4} vs \texttt{ppo\_3} & seed 7 & seed 8 \\
\midrule
sampled ($T=1$) & 601 / 385 / 14 (60.8\% score) & 603 / 390 / 7 (60.7\%) \\
greedy ($T=0$)  & 556 / 400 / 44 (57.8\%)       & 554 / 401 / 45 (57.6\%) \\
\bottomrule
\end{tabular}
\end{center}
Against the scripted opponents it held or gained ground (\texttt{vec/diagnose.py}, 600 sampled
games, seed 0): 574 / 23 / 3 against the original expert (\texttt{ppo\_3}: 577 / 20 / 3) and
520 / 70 / 10 against the tuned \texttt{expert3} (\texttt{ppo\_3}: 505 / 85 / 10, 18\% fewer
losses). In the original environment it won 398 of 400 greedy games with no losses. The Marshal was
lost 0.47 times per game against the expert (\texttt{ppo\_3}: 0.56), but the share of Marshal and
General losses that came from attacking a Bomb did not fall (32\% and 44\%). One habit got worse:
attacks with a non-Scout on an enemy piece already revealed to beat it rose from 0.10 to 0.22 per
game against the expert and from 0.09 to 0.16 against \texttt{expert3}, and attacks on revealed
pieces lost 30\% of the time instead of 24\%. \texttt{ppo\_4} replaced \texttt{ppo\_3} in
\texttt{web/model.onnx} (ONNX equal to PyTorch within $2\times10^{-5}$; \texttt{ppo\_3}'s export kept
in \texttt{runs/ppo\_4/web\_model\_prev/}).

\paragraph{Lesson.} A flat head-to-head early in a continuation run is not yet a plateau: here the
gain arrived only once the learning rate had annealed, so a run should be judged at its end, on fresh
seeds. Head-to-head strength and the expert scores also hide changes in specific habits; the
per-piece diagnosis caught a doubled rate of attacks that a human would see as plain blunders, which
is the kind of weakness an opponent who adapts would use.

% ------------------------------------------------------------------------------------------
\section{Measuring past the scripted opponents: a ladder, an exploiter, and a memory probe}
\label{sec:phase-a}

\paragraph{Change.} Once \texttt{ppo\_4} beat every scripted opponent, three measurements replaced
them. (1) An Elo ladder (\texttt{vec/ladder.py}): a round robin of twelve players, 400 sampled games
per pair over two setup banks, fitted with a Bradley--Terry model (draws as half a win, the original
expert fixed at 1500) and 1\,000 bootstrap resamples for the intervals. (2) An exploiter test
(\texttt{vec/exploit.py}): a fresh PPO learner, started from the clone, plays only against one
frozen target for 45 minutes; how fast it learns to win measures how easy the target is to beat
by adapting to it, which is what a human does. (3) A supervised memory probe
(\texttt{vec/memory/}): 20\,000 games of \texttt{ppo\_4}, split by game, on which four belief
models of equal size predicted the rank of every hidden enemy piece from the 59 planes alone, the
59 planes plus hand-made history planes (times moved, approaches and retreats, battles), or the 59
planes plus a learned per-piece memory over its last 24 events (the thesis's AAREN, or a GRU).

\paragraph{Why.} Section~\ref{sec:ppo4}: the scripted experts no longer separate the networks, and
the belief head had stayed near the same accuracy since \texttt{ppo\_1}. The working explanation
was that the network lacks memory, which is what AAREN was meant to supply.

\paragraph{Result.} \emph{Ladder.} The clone rates 1460, \texttt{ppo\_1} 1804, \texttt{ppo\_2}
1980, \texttt{ppo\_3} 2026 and \texttt{ppo\_4} 2089, with 95\% intervals of about $\pm 20$ that
do not overlap; the tuned experts rate about 1690--1700. The pairwise scores agree with the
fresh-seed head-to-heads (\texttt{ppo\_4} against \texttt{ppo\_3}: 59.6\% here, 60.7\% in
Section~\ref{sec:ppo4}).

\emph{Exploiter.} Against the original expert the exploiter passed 50\% within four minutes and
ended at 93--96\%, every win by capturing the Flag before the expert committed its strong pieces.
Against \texttt{ppo\_4} it went from 2\% to 8\% in the same time and was still rising. Because the
clone starts about 630 Elo below \texttt{ppo\_4} and almost never wins at first, this shows only that
\texttt{ppo\_4} has no weakness simple enough to find from that level in 45 minutes, not that it has
none.

\emph{Memory probe.} On the test games:
\begin{center}
\small
\begin{tabular}{lccc}
\toprule
Input & accuracy, all hidden pieces & moved pieces & AUC Marshal/General (moved) \\
\midrule
public rank counts only        & 23.3\% & 23.2\% & 0.687 \\
59 planes                      & 31.6\% & 31.6\% & 0.888 \\
59 planes + hand-made history  & 32.0\% & 34.8\% & 0.918 \\
59 planes + AAREN memory       & 32.0\% & 35.6\% & 0.928 \\
59 planes + GRU memory         & 32.0\% & 35.4\% & 0.927 \\
\bottomrule
\end{tabular}
\end{center}
History helps on pieces that have moved (three to four points), but those are about a sixth of the
hidden pieces, so the overall gain is 0.4 points. For unmoved pieces, which is where the Bombs that
kill the Marshal and General are, it adds nothing (AUC 0.743 with or without memory): an unmoved
piece has no history. The three memory designs reach the same ceiling, and \texttt{ppo\_4}'s own
belief head (31.4\%) matches a model trained for this alone, so the flat belief curve is a limit
of the information in the observation, not a failure to learn. The probe reads only public
information, and a leak could only have inflated a gain.

\emph{The blunders of Section~\ref{sec:ppo4}.} Over 600 games against each expert, recording
\texttt{ppo\_4}'s own value estimate before each move, attacks on a revealed piece that beats the
attacker (non-Scouts) occurred 0.3--0.7 times per 1\,000 moves when its value was 0.5--0.8 and
3.7--4.5 times when it was at least 0.8, and \texttt{ppo\_3} shows the same pattern at a lower
level. They did not cluster near the 200-ply no-battle limit (none after 150 quiet plies), so they
are not sacrifices to reset the draw clock. Games with such a blunder were won as often as games
without (95\% against 95\% against the expert).

\paragraph{Lesson.} Before building a module to supply missing information, measure how much of it
there is: here the probe showed in an hour that move history adds almost nothing in these games,
and that the real gap, unmoved Bombs, is one no memory can close. And a reward paid only for the
result leaves the agent indifferent to material once it believes the game is won; against
scripted opponents that costs nothing, but it is the kind of play a human notices and can turn
around when the network's confidence is wrong.

% ------------------------------------------------------------------------------------------
\section{Engineering notes}
\label{sec:engineering}

\begin{itemize}
  \item formal\_6 and formal\_7 were both killed partway, with no Python error. formal\_7's error
        log showed \texttt{forrtl: error (200): program aborting due to window-CLOSE event} once
        for each process. The Intel Fortran runtime that numpy's maths library loads installs a
        console handler that aborts the process when its console window closes. Starting the run
        detached from the app's shell was not enough, because the launcher still gave it a
        console. Runs are now started with \texttt{pythonw.exe}, which has no console, with the
        output sent to log files and \texttt{FOR\_DISABLE\_CONSOLE\_CTRL\_HANDLER=1} set. A
        long run should not own a console window that something else can close.
  \item An \texttt{import time} statement inside the training loop made \texttt{time} a local
        name for the whole training function, so any use of \texttt{time} before that line
        failed. Modules should be imported once, at the top of the file.
  \item A checkpoint loader that silently falls back to a fresh network turns a missing file into
        a misleading result: a probe of a checkpoint that did not exist yet reported the values of
        an untrained network. Diagnostic scripts now assert that the file exists.
  \item Most source files have Windows line endings. An edit script that matches multi-line text
        must normalise them first, or it silently matches nothing; a patch helper that asserts the
        number of matches made every such miss fail loudly instead.
  \item The first run of the capacity probe above passed an all-zero next-state distribution to
        the C51 projection for its terminal targets. The projection moves probability mass, so
        the targets were empty, the loss was zero, and no network learnt anything. Every
        configuration then gave the same numbers to three digits, which exposed it. A probe
        should assert that its targets are distributions (they sum to one) before trusting a
        negative result.
  \item Two training runs and the probes shared one 6\,GB laptop GPU for part of the formal\_9 to
        formal\_10 comparison; an episode then took 4 to 33\,s instead of about 2 to 4\,s. Runs on
        one GPU are compared at equal episodes, and a run that only serves as a control is stopped
        once its comparison point is reached.
  \item The trainer's progress bar rewrites one line with carriage returns, so a log watcher that
        filters lines sees one ever-growing line that matches everything. Watching the plain log for
        the evaluation lines (or turning carriage returns into newlines first) avoids it.
  \item The planned Rust version of the environment could not be built: this machine has no Rust
        toolchain, and installing one is the owner's decision. The parts it would take over (the
        environment step, the shaping, the observation and outlook planes, the scripted players)
        are pure functions of the board and are covered by the tests that compare the Python and
        browser engines, which a Rust port can reuse.
  \item A patch helper made tolerant of trailing whitespace also treated the end of a replaced
        block as optional whitespace, and so ate the indentation of the line that followed. Four
        source files no longer compiled. Compiling every edited file at once, before running any
        test, found all of them in one pass.
  \item A nested rollout-completion helper (\texttt{record\_done} in \texttt{vec/ppo.py}) used
        \texttt{ends += ...} on an outer tensor without \texttt{nonlocal ends}, making
        \texttt{ends} local and raising \texttt{UnboundLocalError} the first time a game finished.
        An 8-ply unit-test rollout never reached the end of a game (\texttt{if not done.any():
        return}) and so passed. Replacing the rebinding with in-place \texttt{ends.add\_(...)} and
        setting \texttt{max\_plies=6} in the smoke test so games terminate during rollout caught
        the path on both initial training and checkpoint resume.
  \item A forward helper in the GUI adapter (\texttt{\_pv\_forward} in \texttt{dqn\_bot\_logic.py})
        returned a 2-tuple when \texttt{with\_belief=False} and a 3-tuple when
        \texttt{with\_belief=True}, while evaluation and recommendation callers unpacked three
        values unconditionally, failing whenever queried for the opponent side
        (\texttt{player\_id != self.player\_id}). Helpers with optional outputs should return a
        fixed tuple arity with \texttt{None} in unused slots.
\end{itemize}

% \include{chapters/conclusion} may also be used

%\appendix
%   \include{appendices/evaltool}
%   \include{appendices/reports}
%   \include{appendices/io}
%   \include{appendices/code}
%   \include{appendices/usersguide}

\nocite{*}
\bibliographystyle{acm}
{
\singlespace
\bibliography{references}
}

\include{misc/vita}

\end{document}
